\documentclass[12pt]{article}
        \usepackage[T1]{fontenc}
        \usepackage[utf8]{inputenc}
        \usepackage[margin=1in]{geometry}
        \usepackage[round,sort]{natbib}
        \usepackage[colorlinks=true,citecolor=blue,urlcolor=blue,linkcolor=black]{hyperref}
        \usepackage{enumitem}
        \usepackage{longtable}
        \usepackage{booktabs}
        \setcitestyle{round}
        \title{Writing-Oriented Literature Review for WeChat Media, Information, and Opinion Formation}
        \author{Codex draft for Ziwen Zu}
        \date{\today}

        \begin{document}
        \maketitle

        \section{Introduction}

        This review reconstructs the literature most relevant to a project on WeChat media, official communication, information use, and opinion formation in China. The corpus comes from the \texttt{Information} collection in the Literature vault and is filtered to unique papers after duplicate detection. The central analytic payoff of a WeChat-focused review is that WeChat cannot be treated as a generic social media platform. It bundles semi-private messaging, official accounts, recommendation channels, service delivery, payment infrastructure, and routine administrative interaction. That architecture changes what digital traces mean. A read, like, share, comment, or subscription on WeChat often combines informational use with social visibility and potential political interpretation.

        The existing literature gives strong pieces of this story but rarely assembles them into a single framework. WeChat-specific studies show that official accounts are used for service delivery, audience acquisition, and covert influence. China-focused work on propaganda and censorship shows that the state manages attention and discourse through a mix of selective suppression, soft persuasion, and strategic amplification. A broader media-effects literature identifies mechanisms of exposure, misinformation, credibility, and networked discussion, but often on platforms with more public and flatter interaction structures than WeChat. The practical implication for this project is that a WeChat-centered theory should distinguish among exposure, instrumental use, visible endorsement, and socially legible recommendation rather than collapsing them into one measure of ``engagement.''

        \section{Major Clusters in the Literature}

        \subsection{Official Accounts as Mixed Governance and Communication Infrastructure}

        A first cluster treats official digital accounts not as pure propaganda outlets, but as hybrid governance infrastructure. Work on Chinese government social media repeatedly shows that these channels are used to provide practical information, route users toward services, and sustain routine contact with citizens rather than only to transmit ideology \citep{Zheng2013,ZhengZheng2014,jiang_2021_government_information_quarterly,guo_2021_information_management,li_2022_policy_internet,jiang_2023_jpart,yuan_2023_government_information_quarterly}. The central claim is that official accounts solve administrative and communicative problems at the same time: they lower search costs for users, increase state visibility, and create a durable audience for subsequent political messaging.

        The strongest evidence in this cluster comes from studies that observe actual use rather than declared intentions. \citet{jiang_2021_government_information_quarterly} shows with clickstream data that a government WeChat official account was used mainly as a lookup tool, with users seeking practical information while many built-in service functions remained underused. That finding is theoretically important because it implies that following or reading official content on WeChat need not mean ideological endorsement. \citet{li_2022_policy_internet} and \citet{jiang_2023_jpart} move from interface use to crisis communication and compliance, arguing that government social media can coordinate information and shape citizen behavior in emergencies. More recent work such as \citet{yu_2026_jcc} pushes the argument further by treating public service information itself as a communication strategy rather than politically residual content.

        What remains unresolved is whether service content merely attracts attention or whether it also changes the political meaning of later contact with official accounts. This is where the literature is especially useful for your project. It suggests that WeChat official accounts are best understood as mixed channels: service delivery may build audience capacity for future persuasion, but it can also crowd out immediate attention to overt propaganda. A WeChat-centered project can therefore contribute by theorizing the trade-off between audience accumulation and political persuasion rather than assuming that all official communication has the same objective.

        \subsection{Semi-Officiality, Decentralized Propaganda, and Covert Influence}

        A second cluster shifts the focus from official messaging to the organizational form of state influence online. Here the literature argues that digital propaganda in China increasingly works through delegated, semi-official, or decentralized channels rather than through a single visible propaganda apparatus \citep{repnikova_2018_jcc,neagli_2021_journal_of_current_chinese_affairs,sun_2022_politics_society,han_2025_cq,lu_2025_ajps,KingPanRoberts2017,LuPan2021}. The key move is away from a simple state-versus-society model and toward a layered ecosystem in which official agencies, affiliated media, semi-official influencers, and platform intermediaries jointly shape discourse.

        The WeChat-specific contribution here comes from \citet{neagli_2021_journal_of_current_chinese_affairs}, which shows that semi-official WeChat public accounts can posture as independent voices while remaining linked to party-state influence. In that account, credibility comes precisely from not looking fully official. This logic resonates with \citet{repnikova_2018_jcc}, who emphasizes participatory persuasion rather than one-way indoctrination, and with \citet{sun_2022_politics_society}, who conceptualizes censorship as delegated, multistage, and layered. Work outside WeChat but within the same broader argument, such as \citet{lu_2025_ajps} on Douyin and \citet{han_2025_cq} on ``Propaganda State 2.0,'' suggests that the Chinese state increasingly adapts propaganda to platform-native content styles and organizational decentralization.

        The unresolved issue is boundary-drawing. At what point does service communication become soft propaganda, and when does a semi-official account become analytically distinct from an official one? This matters directly for a WeChat project because WeChat's account ecology blurs institutional identities more than many public platforms do. The literature makes clear that state influence is often covert, distributed, and embedded in apparently ordinary content. What it still lacks is a clear account of how citizens interpret and behaviorally respond to that ambiguity on WeChat itself.

        \subsection{Affordances, Visibility, and the Meaning of Engagement}

        A third cluster is the most directly relevant to your theoretical contribution. These studies imply that digital traces are not interchangeable because platforms vary in publicness, network structure, and the visibility of endorsement \citep{jiang_2021_government_information_quarterly,medaglia_2017_government_information_quarterly,wukich_2021_public_performance_management_review,wukich_2022_government_information_quarterly,TsaiMen2018,Xu2022WeChat,tai_2020_journal_of_communication}. In the WeChat setting, this point is especially sharp because the platform combines private chats, public accounts, bookmarking, forwarding, and semi-public recommendation mechanisms in one app.

        Existing WeChat-related work, however, only partly solves the problem. \citet{jiang_2021_government_information_quarterly} shows how users navigate official accounts, but it is mainly about service use and information architecture. \citet{Xu2022WeChat} and \citet{TsaiMen2018} help by showing that sharing, engagement, and organizational-public interaction on WeChat depend on the structure of the platform rather than only on message content. Studies of censorship and platform design such as \citet{tai_2020_journal_of_communication} imply that visibility and specificity change what kinds of messages survive and circulate. Yet the literature still rarely distinguishes among low-visibility consumption, quasi-private approval, and socially legible endorsement in a single design.

        This gap is precisely where a WeChat-centered project can intervene. The strongest claim you can make is not just that WeChat matters, but that WeChat changes what counts as political expression. On more public platforms, a like or share is often interpreted as an overt signal. On WeChat, similar actions may range from purely instrumental use to socially visible recommendation. The literature points toward this distinction, but it has not yet theorized or measured it cleanly enough.

        \subsection{Exposure, Selection, and Persuasion in Restricted Information Environments}

        A fourth cluster asks when information exposure changes beliefs, attitudes, or behavior. The strongest papers here move away from a simple ``propaganda works'' versus ``propaganda fails'' dichotomy and instead emphasize selective exposure, framing, and audience predispositions \citep{pan_2022_psrm,shen_2025_psrm,WangWestwood2024,allcott_2020_aer,levy_2021_aer,PanXu2018}. For a WeChat project, the central lesson is that the political effect of content depends heavily on whether exposure is forced, chosen, or embedded in another activity such as service use.

        China-focused work is especially valuable because it separates the effects of state-controlled content from the conditions under which users encounter it. \citet{shen_2025_psrm} explicitly distinguishes forced from selective exposure to propaganda in China. \citet{pan_2022_psrm} shows that government-controlled media can shift attitudes through framing, while \citet{WangWestwood2024} highlights selective reading even in a restricted information environment. Comparative media-effects studies such as \citet{allcott_2020_aer} and \citet{levy_2021_aer} are not about China or WeChat, but they help identify the kinds of designs needed to separate content effects from self-selection.

        The unresolved question is measurement. If a user reads an official WeChat article because it contains utility-relevant information, is that exposure politically meaningful in the same way as a visible recommendation or share? The literature suggests that the answer is no, but it lacks a WeChat-specific design that tracks how different kinds of platform actions map onto different kinds of political inference. Your project can usefully position itself as solving that measurement problem.

        \subsection{Censorship, Self-Censorship, and Opinion Expression under Surveillance}

        A fifth cluster links information control to expression rather than only to content availability. The common claim is that censorship, surveillance, and anticipated audiences shape what citizens say, search for, and circulate \citep{tai_2014_journal_of_east_asian_studies,sun_2022_politics_society,ye_2023_discourse_context_media,ShenTruex2020,Kuran1995,NoelleNeumann1993,xu_2022_jop}. In this view, the most important political outcome is often not persuasion but strategic adaptation: users learn what is safe to post, how to recode sensitive content, and when silence is preferable to visibility.

        This literature is conceptually crucial for a WeChat project even when it is not WeChat-specific. \citet{ye_2023_discourse_context_media} shows how censorship generates recoding practices and a broader ``sensitive word'' culture. \citet{sun_2022_politics_society} emphasizes a layered regime of information control rather than one-shot deletion. Foundational work on self-censorship and preference falsification such as \citet{ShenTruex2020}, \citet{Kuran1995}, and \citet{NoelleNeumann1993} explains why visible expression is a particularly noisy indicator of genuine belief under political pressure. The implication is that the absence of overt dissent, or the presence of mild positive engagement, cannot be read straightforwardly as support.

        What remains unresolved is how these expressive pressures interact with WeChat's unusual blend of private and semi-public spaces. A platform centered on known social ties may generate different reputational risks than a more open public platform. This is precisely why a WeChat project should not simply import findings from Twitter, Facebook, or generic ``social media'' studies. The more productive move is to use this literature to theorize layered expression under conditions of audience uncertainty and asymmetric visibility.

        \section{Debates and Points of Tension}

        The first major debate concerns the meaning of engagement. Some studies implicitly treat following, reading, sharing, or reacting to official content as evidence of persuasion or political alignment. Others suggest that the same behaviors may instead reflect instrumental information use, routine service seeking, or risk-managed compliance. This is not a semantic disagreement. It is a measurement dispute, and it matters most on WeChat, where the line between service use and political communication is unusually thin.

        A second debate concerns the mechanism of state influence. One strand emphasizes censorship and coercive control \citep{tai_2014_journal_of_east_asian_studies,sun_2022_politics_society}. Another emphasizes soft propaganda, framing, audience acquisition, and content packaging \citep{repnikova_2018_jcc,han_2025_cq,LuPan2021}. The disagreement is partly theoretical, but it is also empirical: studies look at different outcomes, from deletions and discourse management to attitudes and compliance. For your project, the key synthesis is that these mechanisms are complementary rather than mutually exclusive. Service content, soft propaganda, and hard control can coexist within the same WeChat channel.

        A third debate concerns external validity across platforms. Many of the clearest causal designs in the broader media-effects literature use Facebook, Twitter, or generalized online environments. Those studies are useful for identifying mechanisms of exposure, misinformation, and network effects, but they do not solve the WeChat problem. WeChat integrates messaging, official broadcasting, payments, and everyday administrative interaction in a way that is institutionally denser and socially less public than the platforms on which most causal evidence is built.

        A fourth debate concerns level of analysis. Some papers study central propaganda strategy, others local government communication, and others platform-wide censorship or citizen behavior. Yet a WeChat project sits at the intersection of all three. Official accounts are run by specific agencies, embedded in a broader censorship regime, and interpreted by users within concrete social networks. One reason the literature feels fragmented is that these levels are often studied separately. A strong WeChat-centered project can contribute by connecting them: institutional communication strategy, platform design, and citizen behavioral response.

        \section{Annotated Paper-by-Paper Review}
        \subsection{Bertot et al. 2012: The impact of polices on government...}
\textbf{Citation key}: \texttt{bertot\_2012\_government\_information\_quarterly}.\\
\textbf{Full citation}: John Carlo Bertot, Paul T. Jaeger, and Derek Hansen 2012. "The impact of polices on government social media usage: Issues, challenges, and recommendations." Government Information Quarterly.

\paragraph{Summary} This paper studies the impact of polices on government social media usage: issues, challenges, and recommendations. Government agencies are increasingly using social media to connect with those they serve.

\paragraph{Context} However, interacting via social media introduces new challenges related to privacy, security, data management, accessibility, social inclusion, governance, and other information policy issues.

\paragraph{Main Argument} Core mechanism maps onto the cluster on platform affordances and public-private boundaries, but the exact statement should be verified from the full text if heavily used.

\paragraph{Independent Variable(s)} polices

\paragraph{Dependent Variable(s)} government social media usage: Issues

\paragraph{Methods and Design} Descriptive or unclear from available parse. Observational or descriptive design; verify details from full text if heavily used. Data / sample: However, interacting via social media introduces new challenges related to privacy, security, data management, accessibility, social inclusion, governance, and other information policy issues. These tools vary...

\paragraph{Findings} No clean finding sentence detected; verify from full text if central.

\paragraph{Contribution} Clarifies how state agencies use digital channels for communication and governance rather than treating online politics as purely oppositional.

\paragraph{Limitations} Design details need closer verification before the paper is used for strong causal claims.

\paragraph{Relevance to WeChat / My Project} Indirect benchmark: useful for comparison on platform effects or official communication, but needs careful translation to China's WeChat-centered environment.

\paragraph{Related Debate} How interface design changes the visibility and meaning of digital expression


\subsection{Magro 2012: A Review of Social Media Use...}
\textbf{Citation key}: \texttt{magro\_2012\_administrative\_sciences}.\\
\textbf{Full citation}: Michael J. Magro 2012. "A Review of Social Media Use in E-Government." Administrative Sciences.

\paragraph{Summary} This study examines a review of social media use in e-government using descriptive or unclear from available parse evidence and reports that the ability to easily find similar projects would be a great advantage to those that follow others.

\paragraph{Context} In the past few years, e-government has been a topic of much interest among those excited about the advent of Web 2.0 technologies. This paper reviews the recent literature concerning Web 2.0, social media, social networking, and how it has been used in...

\paragraph{Main Argument} The paper's core mechanism is that the ability to easily find similar projects would be a great advantage to those that follow others.

\paragraph{Independent Variable(s)} it has been used in the public sector. Key observations include literature themes such as the evolution of social...

\paragraph{Dependent Variable(s)} is still needed in government culture

\paragraph{Methods and Design} Descriptive or unclear from available parse. Observational or descriptive design; verify details from full text if heavily used. Data / sample: Key observations include literature themes such as the evolution of social media case studies in the literature, the progress of social media policies and strategies over time, and social media use in disaster...

\paragraph{Findings} The ability to easily find similar projects would be a great advantage to those that follow others.

\paragraph{Contribution} Provides an external benchmark for mechanisms that may or may not travel to China and WeChat.

\paragraph{Limitations} Design details need closer verification before the paper is used for strong causal claims.

\paragraph{Relevance to WeChat / My Project} Indirect benchmark: useful for comparison on platform effects or official communication, but needs careful translation to China's WeChat-centered environment.

\paragraph{Related Debate} Whether official digital communication mainly informs, disciplines, mobilizes, or improves responsiveness


\subsection{Mergel 2013: A framework for interpreting social media...}
\textbf{Citation key}: \texttt{mergel\_2013\_government\_information\_quarterly}.\\
\textbf{Full citation}: Ines Mergel 2013. "A framework for interpreting social media interactions in the public sector." Government Information Quarterly.

\paragraph{Summary} This paper studies a framework for interpreting social media interactions in the public sector. Social media applications are extending the information and communication technology landscape in the public sector and are used to increase government transparency, participation and collaboration in the U.S. federal...

\paragraph{Context} Many government agencies are experimenting with the use of social media, however very few actively measure the impact of their digital interactions.

\paragraph{Main Argument} Core mechanism maps onto the cluster on interpersonal discussion, networks, and opinion formation, but the exact statement should be verified from the full text if heavily used.

\paragraph{Independent Variable(s)} their digital interactions. This article builds

\paragraph{Dependent Variable(s)} insights from social media directors in the U

\paragraph{Methods and Design} interviews / focus groups. Observational or descriptive design; verify details from full text if heavily used. Data / sample: Case selection and context In order to understand the information needs of government social media professionals when they use social media applications for their government agency, a qualitative approach was used to...

\paragraph{Findings} No clean finding sentence detected; verify from full text if central.

\paragraph{Contribution} Provides an external benchmark for mechanisms that may or may not travel to China and WeChat.

\paragraph{Limitations} Design details need closer verification before the paper is used for strong causal claims.

\paragraph{Relevance to WeChat / My Project} Indirect benchmark: useful for comparison on platform effects or official communication, but needs careful translation to China's WeChat-centered environment.

\paragraph{Related Debate} How social ties and discussion networks mediate willingness to speak, share, and update opinions


\subsection{Reuter and Szakonyi 2013: Online Social Media and Political Awareness...}
\textbf{Citation key}: \texttt{reuter\_2013\_bjps}.\\
\textbf{Full citation}: Ora John Reuter and David Szakonyi 2013. "Online Social Media and Political Awareness in Authoritarian Regimes." British Journal of Political Science.

\paragraph{Summary} This study examines online social media and political awareness in authoritarian regimes using survey evidence and reports that survey data from the 2011 Russian parliamentary elections show that usage of Twitter and Facebook, which were politicized by opposition elites, significantly increased respondents'...

\paragraph{Context} \_S0007123413000203 How to cite this article: Ora John Reuter and David Szakonyi (2015).

\paragraph{Main Argument} The paper's core mechanism is that survey data from the 2011 Russian parliamentary elections show that usage of Twitter and Facebook, which were politicized by opposition elites, significantly increased respondents' perceptions of electoral fraud, while...

\paragraph{Independent Variable(s)} Not stated in standard IV/DV terms; inferred from title/abstract if needed

\paragraph{Dependent Variable(s)} Not stated in standard IV/DV terms; verify if central to argument

\paragraph{Methods and Design} survey. Survey-based observational analysis. Data / sample: Survey data from the 2011 Russian parliamentary elections show that usage of Twitter and Facebook, which were politicized by opposition elites, significantly increased respondents' perceptions of electoral fraud,...

\paragraph{Findings} Survey data from the 2011 Russian parliamentary elections show that usage of Twitter and Facebook, which were politicized by opposition elites, significantly increased respondents' perceptions of electoral fraud, while usage of Russia's domestic social networking platforms, VKontakte and Odnoklassniki, which were not politicized by opposition activists, had no effect on perceptions of fraud.

\paragraph{Contribution} Provides an external benchmark for mechanisms that may or may not travel to China and WeChat.

\paragraph{Limitations} Design details need closer verification before the paper is used for strong causal claims.

\paragraph{Relevance to WeChat / My Project} Indirect benchmark: useful for comparison on platform effects or official communication, but needs careful translation to China's WeChat-centered environment.

\paragraph{Related Debate} How social ties and discussion networks mediate willingness to speak, share, and update opinions


\subsection{Tai 2014: China's Media Censorship A Dynamic and...}
\textbf{Citation key}: \texttt{tai\_2014\_journal\_of\_east\_asian\_studies}.\\
\textbf{Full citation}: Qiuqing Tai 2014. "China's Media Censorship: A Dynamic and Diversified Regime." Journal of East Asian Studies.

\paragraph{Summary} This paper studies china's media censorship: a dynamic and diversified regime. \_id=1315882 (accessed May 1,2013).

\paragraph{Context} \_id=1315882 (accessed May 1,2013). Gody, Ahmed.

\paragraph{Main Argument} Core mechanism maps onto the cluster on censorship, propaganda, and state influence, but the exact statement should be verified from the full text if heavily used.

\paragraph{Independent Variable(s)} Not stated in standard IV/DV terms; inferred from title/abstract if needed

\paragraph{Dependent Variable(s)} Not stated in standard IV/DV terms; verify if central to argument

\paragraph{Methods and Design} field experiment. field experiment. Data / sample: Unclear from available parse; verify from full text if used centrally

\paragraph{Findings} No clean finding sentence detected; verify from full text if central.

\paragraph{Contribution} Offers comparatively strong leverage on causal mechanisms in digital political communication.

\paragraph{Limitations} Internal validity is strong, but external validity may depend on whether experimental treatments resemble real platform use.

\paragraph{Relevance to WeChat / My Project} Conceptually useful for a WeChat project because it speaks to the Chinese information environment, especially the cluster on censorship, propaganda, and state influence, even if the platform is not WeChat.

\paragraph{Related Debate} Whether digital control works through persuasion, selective amplification, covert influence, or audience management


\subsection{BOND and MESSING 2015: Quantifying Social Media s Political Space...}
\textbf{Citation key}: \texttt{bond\_2015\_apsr}.\\
\textbf{Full citation}: ROBERT BOND and SOLOMON MESSING 2015. "Quantifying Social Media's Political Space: Estimating Ideology from Publicly Revealed Preferences on Facebook." American Political Science Review.

\paragraph{Summary} This study examines quantifying social media's political space: estimating ideology from publicly revealed preferences on facebook using descriptive or unclear from available parse evidence and reports that 1 February 2015 c American Political Science Association 2015  doi:10.1017/S0003055414000525 Quantifying...

\paragraph{Context} American Political Science Review Vol. 109, No.

\paragraph{Main Argument} The paper's core mechanism is that 1 February 2015 c American Political Science Association 2015  doi:10.1017/S0003055414000525 Quantifying Social Media's Political Space: Estimating Ideology from Publicly Revealed Preferences on Facebook ROBERT BOND Ohio...

\paragraph{Independent Variable(s)} Not stated in standard IV/DV terms; inferred from title/abstract if needed

\paragraph{Dependent Variable(s)} Not stated in standard IV/DV terms; verify if central to argument

\paragraph{Methods and Design} Descriptive or unclear from available parse. Observational or descriptive design; verify details from full text if heavily used. Data / sample: 1 February 2015 c American Political Science Association 2015  doi:10.1017/S0003055414000525 Quantifying Social Media's Political Space: Estimating Ideology from Publicly Revealed Preferences on Facebook ROBERT BOND...

\paragraph{Findings} 1 February 2015 c American Political Science Association 2015  doi:10.1017/S0003055414000525 Quantifying Social Media's Political Space: Estimating Ideology from Publicly Revealed Preferences on Facebook ROBERT BOND Ohio State University SOLOMON MESSING Facebook Data Science W e demonstrate that social media data represent a useful resource for testing models of legislative and individual-level political behavior and attitudes.

\paragraph{Contribution} Provides an external benchmark for mechanisms that may or may not travel to China and WeChat.

\paragraph{Limitations} Design details need closer verification before the paper is used for strong causal claims.

\paragraph{Relevance to WeChat / My Project} Indirect benchmark: useful for comparison on platform effects or official communication, but needs careful translation to China's WeChat-centered environment.

\paragraph{Related Debate} How findings from facebook and comparative / general settings travel to a WeChat-centered China project


\subsection{Gunitsky 2015: Corrupting the Cyber-Commons Social Media as...}
\textbf{Citation key}: \texttt{gunitsky\_2015\_perspectives\_on\_politics}.\\
\textbf{Full citation}: Seva Gunitsky 2015. "Corrupting the Cyber-Commons: Social Media as a Tool of Autocratic Stability." Perspectives on Politics.

\paragraph{Summary} This study examines corrupting the cyber-commons: social media as a tool of autocratic stability using descriptive or unclear from available parse evidence and reports that most importantly, social media may decrease the incentives for hybrid and autocratic regimes to hold elections in order to reveal hidden...

\paragraph{Context} 51663057 Gehlbach, Scott. 2012.

\paragraph{Main Argument} The paper's core mechanism is that most importantly, social media may decrease the incentives for hybrid and autocratic regimes to hold elections in order to reveal hidden information and bolster regime legitimacy.

\paragraph{Independent Variable(s)} Not stated in standard IV/DV terms; inferred from title/abstract if needed

\paragraph{Dependent Variable(s)} Not stated in standard IV/DV terms; verify if central to argument

\paragraph{Methods and Design} Descriptive or unclear from available parse. Observational or descriptive design; verify details from full text if heavily used. Data / sample: "Authoritarianism Online: What Can We Learn from Internet Data in Nondemocracies?" PS: Political Science and Politics 46(2): 262--70.

\paragraph{Findings} Most importantly, social media may decrease the incentives for hybrid and autocratic regimes to hold elections in order to reveal hidden information and bolster regime legitimacy.

\paragraph{Contribution} Provides an external benchmark for mechanisms that may or may not travel to China and WeChat.

\paragraph{Limitations} Design details need closer verification before the paper is used for strong causal claims.

\paragraph{Relevance to WeChat / My Project} Indirect benchmark: useful for comparison on platform effects or official communication, but needs careful translation to China's WeChat-centered environment.

\paragraph{Related Debate} Whether digital control works through persuasion, selective amplification, covert influence, or audience management


\subsection{Priks 2015: The Effects of Surveillance Cameras on...}
\textbf{Citation key}: \texttt{priks\_2015\_ej}.\\
\textbf{Full citation}: Mikael Priks 2015. "The Effects of Surveillance Cameras on Crime: Evidence from the Stockholm Subway." The Economic Journal.

\paragraph{Summary} This paper studies the effects of surveillance cameras on crime: evidence from the stockholm subway. of the cameras reduced crime by approximately 25\% at stations in the city centre.

\paragraph{Context} of the cameras reduced crime by approximately 25\% at stations in the city centre. The types of crimes deterred by cameras are planned crime, that is, pickpocketing and robbery.

\paragraph{Main Argument} Core mechanism maps onto the cluster on comparative digital media benchmarks, but the exact statement should be verified from the full text if heavily used.

\paragraph{Independent Variable(s)} Surveillance Cameras

\paragraph{Dependent Variable(s)} Crime: Evidence from the Stockholm Subway

\paragraph{Methods and Design} Descriptive or unclear from available parse. Observational or descriptive design; verify details from full text if heavily used. Data / sample: I also had access to data on crime in the areas adjacent to the subway stations.

\paragraph{Findings} No clean finding sentence detected; verify from full text if central.

\paragraph{Contribution} Provides an external benchmark for mechanisms that may or may not travel to China and WeChat.

\paragraph{Limitations} Mechanisms may not transfer cleanly to WeChat's hybrid private-public architecture without further adaptation.

\paragraph{Relevance to WeChat / My Project} Peripheral for a WeChat project; best used as a comparative benchmark rather than core theoretical evidence.

\paragraph{Related Debate} How findings from not platform-specific and comparative / general settings travel to a WeChat-centered China project


\subsection{Shadmehr and Bernhardt 2015: State Censorship}
\textbf{Citation key}: \texttt{shadmehr\_2015\_aejmicro}.\\
\textbf{Full citation}: Mehdi Shadmehr and Dan Bernhardt 2015. "State Censorship." American Economic Journal: Microeconomics.

\paragraph{Summary} This paper studies state censorship. from strategic di nation decisions by the media (see e.g., Duggan and Martinelli (2011 ) or th papers in our literature review).

\paragraph{Context} Thus, we assume that when the media ob and the ruler does not censor, the media publicly reveals 75 However, the r censor the media at a cost c > 0, preventing the media from conveying 7 representative citizen, in which case the citizen does not receive...

\paragraph{Main Argument} Core mechanism maps onto the cluster on censorship, propaganda, and state influence, but the exact statement should be verified from the full text if heavily used.

\paragraph{Independent Variable(s)} Not stated in standard IV/DV terms; inferred from title/abstract if needed

\paragraph{Dependent Variable(s)} Not stated in standard IV/DV terms; verify if central to argument

\paragraph{Methods and Design} Descriptive or unclear from available parse. Observational or descriptive design; verify details from full text if heavily used. Data / sample: Thus, we assume that when the media ob and the ruler does not censor, the media publicly reveals 75 However, the r censor the media at a cost c > 0, preventing the media from conveying 7 representative citizen, in...

\paragraph{Findings} No clean finding sentence detected; verify from full text if central.

\paragraph{Contribution} Provides an external benchmark for mechanisms that may or may not travel to China and WeChat.

\paragraph{Limitations} Mechanisms may not transfer cleanly to WeChat's hybrid private-public architecture without further adaptation.

\paragraph{Relevance to WeChat / My Project} Peripheral for a WeChat project; best used as a comparative benchmark rather than core theoretical evidence.

\paragraph{Related Debate} Whether digital control works through persuasion, selective amplification, covert influence, or audience management


\subsection{Tai 2016: Western Media Exposure and Chinese Immigrants...}
\textbf{Citation key}: \texttt{tai\_2016\_polcomm}.\\
\textbf{Full citation}: Qiuqing Tai 2016. "Western Media Exposure and Chinese Immigrants' Political Perceptions." Political Communication.

\paragraph{Summary} This paper studies western media exposure and chinese immigrants' political perceptions. \_id= 2105407 Shirk, S. (2010).

\paragraph{Context} \_id= 2105407 Shirk, S. (2010). Changing media, changing China.

\paragraph{Main Argument} Core mechanism maps onto the cluster on comparative digital media benchmarks, but the exact statement should be verified from the full text if heavily used.

\paragraph{Independent Variable(s)} Not stated in standard IV/DV terms; inferred from title/abstract if needed

\paragraph{Dependent Variable(s)} Not stated in standard IV/DV terms; verify if central to argument

\paragraph{Methods and Design} survey experiment. survey experiment. Data / sample: Retrieved from http://www.nytimes.com/2008/06/04/world/asia/04china.html To test my hypotheses, I conducted a survey experiment among Chinese student immigrants in the United States. One problem that has hampered...

\paragraph{Findings} No clean finding sentence detected; verify from full text if central.

\paragraph{Contribution} Offers comparatively strong leverage on causal mechanisms in digital political communication.

\paragraph{Limitations} Internal validity is strong, but external validity may depend on whether experimental treatments resemble real platform use.

\paragraph{Relevance to WeChat / My Project} Conceptually useful for a WeChat project because it speaks to the Chinese information environment, especially the cluster on comparative digital media benchmarks, even if the platform is not WeChat.

\paragraph{Related Debate} How findings from internet / online media and china settings travel to a WeChat-centered China project


\subsection{Unknown 2017: Annual Review of Political Science}
\textbf{Citation key}: \texttt{annurev\_polisci\_051120\_103422\_2017\_arps}.\\
\textbf{Full citation}: 2017. "Annual Review of Political Science." Annual Review of Political Science.

\paragraph{Summary} This study examines annual review of political science using descriptive or unclear from available parse evidence and reports that this article tries to bridge these divides, to show how knowledge from different fields may be complementary, and to point to shortcomings and blind spots in existing research. By...

\paragraph{Context} See credit lines of images or other third-party material in this article for license information Do media influence policy making?

\paragraph{Main Argument} The paper's core mechanism is that this article tries to bridge these divides, to show how knowledge from different fields may be complementary, and to point to shortcomings and blind spots in existing research. By bringing different strands together, I...

\paragraph{Independent Variable(s)} knowledge from different fields may be complementary, and to point to shortcomings and blind spots in existing...

\paragraph{Dependent Variable(s)} in the media landscape are likely to affect the media--policy making nexus

\paragraph{Methods and Design} Descriptive or unclear from available parse. Observational or descriptive design; verify details from full text if heavily used. Data / sample: Unclear from available parse; verify from full text if used centrally

\paragraph{Findings} This article tries to bridge these divides, to show how knowledge from different fields may be complementary, and to point to shortcomings and blind spots in existing research. By bringing different strands together, I show that media, old and new, are the main arena for the battle over the scope of policy conflict.

\paragraph{Contribution} Provides an external benchmark for mechanisms that may or may not travel to China and WeChat.

\paragraph{Limitations} Mechanisms may not transfer cleanly to WeChat's hybrid private-public architecture without further adaptation.

\paragraph{Relevance to WeChat / My Project} Peripheral for a WeChat project; best used as a comparative benchmark rather than core theoretical evidence.

\paragraph{Related Debate} How findings from not platform-specific and france settings travel to a WeChat-centered China project


\subsection{Chen and Xu 2017: Information Manipulation and Reform in Authoritarian...}
\textbf{Citation key}: \texttt{chen\_2017\_psrm}.\\
\textbf{Full citation}: Jidong Chen and Yiqing Xu 2017. "Information Manipulation and Reform in Authoritarian Regimes." Political Science Research and Methods.

\paragraph{Summary} This study examines information manipulation and reform in authoritarian regimes using descriptive or unclear from available parse evidence and reports that aSSUMPTION 2: 9i 2 I 14 f1; 2g; Ai R; ti ; tj  <= k; Bi R; ti  <= k; Bi R; ti  <= k; 8ti ; tj ; ti 2 ft; tg; 8i 2 I 14 f1; 2g; Ai Q; ti 14 t; tj  <= k; Bi...

\paragraph{Context} than the one in this paper. The latter can be seen as a special case of former.

\paragraph{Main Argument} The paper's core mechanism is that aSSUMPTION 2: 9i 2 I 14 f1; 2g; Ai R; ti ; tj  <= k; Bi R; ti  <= k; Bi R; ti  <= k; 8ti ; tj ; ti 2 ft; tg; 8i 2 I 14 f1; 2g; Ai Q; ti 14 t; tj  <= k; Bi Q; ti 14 t <= k; 8tj 2 ft; tg: ASSUMPTION 3: A 14 Ai R; ti 14...

\paragraph{Independent Variable(s)} Not stated in standard IV/DV terms; inferred from title/abstract if needed

\paragraph{Dependent Variable(s)} Not stated in standard IV/DV terms; verify if central to argument

\paragraph{Methods and Design} Descriptive or unclear from available parse. Observational or descriptive design; verify details from full text if heavily used. Data / sample: The latter can be seen as a special case of former. Moreover, we define Ai x; ti ; tj  14 Vi11 x; ti ; tj   Vi01 x; ti ; tj ; (A1) Bi x; ti  14 Vi10 x; ti   Vi00 x; ti : (A2)...

\paragraph{Findings} ASSUMPTION 2: 9i 2 I 14 f1; 2g; Ai R; ti ; tj  <= k; Bi R; ti  <= k; Bi R; ti  <= k; 8ti ; tj ; ti 2 ft; tg; 8i 2 I 14 f1; 2g; Ai Q; ti 14 t; tj  <= k; Bi Q; ti 14 t <= k; 8tj 2 ft; tg: ASSUMPTION 3: A 14 Ai R; ti 14 t; tj 14 t; B 14 Bi R; ti 14 t do not depend on i. k >= A > B > k.14 Proof of Lemma 1 (a) Conditions Ai R; ti ; tj  <= k; Bi R; ti  <= k in Assumption 2 suggest that it is a dominant strategy for citizen i not to protest under the reform policy R.

\paragraph{Contribution} Helps identify which mechanisms in the Chinese information environment are platform-specific and which generalize beyond a single app.

\paragraph{Limitations} Design details need closer verification before the paper is used for strong causal claims.

\paragraph{Relevance to WeChat / My Project} Conceptually useful for a WeChat project because it speaks to the Chinese information environment, especially the cluster on censorship, propaganda, and state influence, even if the platform is not WeChat.

\paragraph{Related Debate} Whether digital control works through persuasion, selective amplification, covert influence, or audience management


\subsection{Jones and Mattiacci 2017: A Manifesto in 140 Characters or...}
\textbf{Citation key}: \texttt{jones\_2017\_bjps}.\\
\textbf{Full citation}: Benjamin T. Jones and Eleonora Mattiacci 2017. "A Manifesto, in 140 Characters or Fewer: Social Media as a Tool of Rebel Diplomacy." British Journal of Political Science.

\paragraph{Summary} This paper studies a manifesto, in 140 characters or fewer: social media as a tool of rebel diplomacy. B.J.Pol.S.

\paragraph{Context} B.J.Pol.S. 49, 739--761 Copyright  Cambridge University Press, 2017 doi:10.1017/S0007123416000612 First published online 12 April 2017 A Manifesto, in 140 Characters or Fewer: Social Media as a Tool of Rebel Diplomacy BENJAMIN T.

\paragraph{Main Argument} Core mechanism maps onto the cluster on comparative digital media benchmarks, but the exact statement should be verified from the full text if heavily used.

\paragraph{Independent Variable(s)} Not stated in standard IV/DV terms; inferred from title/abstract if needed

\paragraph{Dependent Variable(s)} Not stated in standard IV/DV terms; verify if central to argument

\paragraph{Methods and Design} Descriptive or unclear from available parse. Observational or descriptive design; verify details from full text if heavily used. Data / sample: Its unique social media features make it easier for potential supporters in foreign states (both elites and the masses) to access rebels' information and narratives than would be the case with other communication...

\paragraph{Findings} No clean finding sentence detected; verify from full text if central.

\paragraph{Contribution} Provides an external benchmark for mechanisms that may or may not travel to China and WeChat.

\paragraph{Limitations} Design details need closer verification before the paper is used for strong causal claims.

\paragraph{Relevance to WeChat / My Project} Indirect benchmark: useful for comparison on platform effects or official communication, but needs careful translation to China's WeChat-centered environment.

\paragraph{Related Debate} How findings from twitter / x and comparative / general settings travel to a WeChat-centered China project


\subsection{King et al. 2017: Great Firewall of China and the...}
\textbf{Citation key}: \texttt{king\_2017\_apsr}.\\
\textbf{Full citation}: Gary King, Jennifer Pan, and Margaret E. Roberts 2017. "Great Firewall of China), and the scholarly literature now offers considerable evidence on how and why." American Political Science Review.

\paragraph{Summary} This study examines great firewall of china), and the scholarly literature now offers considerable evidence on how and why using book / theory evidence and reports that "Leaked Emails Show Chinese Regime Employs 500,000 Internet Trolls." Epoch Times. http://j.mp/twomill.

\paragraph{Context} id=2738325. Naher, Anatol-Fiete, and Ivar Krumpal.

\paragraph{Main Argument} The paper's core mechanism is that "Leaked Emails Show Chinese Regime Employs 500,000 Internet Trolls." Epoch Times. http://j.mp/twomill.

\paragraph{Independent Variable(s)} Forgiving Wording and Question Context

\paragraph{Dependent Variable(s)} Social Desirability Bias

\paragraph{Methods and Design} book / theory. Theoretical or historical argument rather than a single empirical design. Data / sample: "Politics, Rumors, and Ambiguity: Tracking Censorship on WeChat's Public Accounts Platform." Munk School of Global Affairs. https://citizenlab.org/2015/07/ tracking-censorship-on-wechat-public-accounts-platform/....

\paragraph{Findings} "Leaked Emails Show Chinese Regime Employs 500,000 Internet Trolls." Epoch Times. http://j.mp/twomill.

\paragraph{Contribution} Directly extends the WeChat-centered literature by specifying platform-specific behavior rather than treating social media as generic.

\paragraph{Limitations} Useful for theory-building, but not designed to isolate a single platform-level causal effect.

\paragraph{Relevance to WeChat / My Project} Directly relevant: this paper studies WeChat behavior or institutions and can be used to theorize visibility, audience composition, and official-account interaction.

\paragraph{Related Debate} Whether semi-private WeChat interactions should be read as information use, endorsement, or strategically managed visibility


\subsection{Medaglia and Zhu 2017: Public deliberation on government-managed social media...}
\textbf{Citation key}: \texttt{medaglia\_2017\_government\_information\_quarterly}.\\
\textbf{Full citation}: Rony Medaglia and Demi Zhu 2017. "Public deliberation on government-managed social media: A study on Weibo users in China." Government Information Quarterly.

\paragraph{Summary} This paper studies public deliberation on government-managed social media: a study on weibo users in china. The evolution of social media platforms provides governments across the world with the potential to achieve objectives of improved information, service provision, citizen engagement, and legitimacy, as well...

\paragraph{Context} However, given the very nature of social media, its primary affordance is not just as a channel for government-citizen interactions, but also as enabler of user-to-user interactions. (Chen, Xu, Cao, \& Zhang, 2016; Treem \& Leonardi, 2012).

\paragraph{Main Argument} Core mechanism maps onto the cluster on platform affordances and public-private boundaries, but the exact statement should be verified from the full text if heavily used.

\paragraph{Independent Variable(s)} Not stated in standard IV/DV terms; inferred from title/abstract if needed

\paragraph{Dependent Variable(s)} Not stated in standard IV/DV terms; verify if central to argument

\paragraph{Methods and Design} Descriptive or unclear from available parse. Observational or descriptive design; verify details from full text if heavily used. Data / sample: The evolution of social media platforms provides governments across the world with the potential to achieve objectives of improved information, service provision, citizen engagement, and legitimacy, as well as to...

\paragraph{Findings} No clean finding sentence detected; verify from full text if central.

\paragraph{Contribution} Helps identify which mechanisms in the Chinese information environment are platform-specific and which generalize beyond a single app.

\paragraph{Limitations} Design details need closer verification before the paper is used for strong causal claims.

\paragraph{Relevance to WeChat / My Project} Conceptually useful for a WeChat project because it speaks to the Chinese information environment, especially the cluster on platform affordances and public-private boundaries, even if the platform is not WeChat.

\paragraph{Related Debate} How interface design changes the visibility and meaning of digital expression


\subsection{Medaglia and Zheng 2017: Mapping government social media research and...}
\textbf{Citation key}: \texttt{medaglia\_2017\_government\_information\_quarterly\_2}.\\
\textbf{Full citation}: Rony Medaglia and Lei Zheng 2017. "Mapping government social media research and moving it forward: A framework and a research agenda." Government Information Quarterly.

\paragraph{Summary} This paper studies mapping government social media research and moving it forward: a framework and a research agenda. In the public sector, social media initiatives are booming, with increasing agreement among managers on the importance of using social media platforms to interact with citizens.

\paragraph{Context} In the public sector, social media initiatives are booming, with increasing agreement among managers on the importance of using social media platforms to interact with citizens. Such initiatives are taken in response to demands from citizens who, as...

\paragraph{Main Argument} Core mechanism maps onto the cluster on participation, responsiveness, and political behavior, but the exact statement should be verified from the full text if heavily used.

\paragraph{Independent Variable(s)} Not stated in standard IV/DV terms; inferred from title/abstract if needed

\paragraph{Dependent Variable(s)} Not stated in standard IV/DV terms; verify if central to argument

\paragraph{Methods and Design} Descriptive or unclear from available parse. Observational or descriptive design; verify details from full text if heavily used. Data / sample: Such initiatives are taken in response to demands from citizens who, as experienced social media users, have increased and matured expectations towards public agencies in terms of responsiveness, information...

\paragraph{Findings} No clean finding sentence detected; verify from full text if central.

\paragraph{Contribution} Clarifies how state agencies use digital channels for communication and governance rather than treating online politics as purely oppositional.

\paragraph{Limitations} Design details need closer verification before the paper is used for strong causal claims.

\paragraph{Relevance to WeChat / My Project} Conceptually useful for a WeChat project because it speaks to the Chinese information environment, especially the cluster on participation, responsiveness, and political behavior, even if the platform is not WeChat.

\paragraph{Related Debate} Whether official digital communication mainly informs, disciplines, mobilizes, or improves responsiveness


\subsection{DePaula et al. 2018: Toward a typology of government social...}
\textbf{Citation key}: \texttt{depaula\_2018\_government\_information\_quarterly}.\\
\textbf{Full citation}: Nic DePaula, Ersin Dincelli, and Teresa M. Harrison 2018. "Toward a typology of government social media communication: Democratic goals, symbolic acts and self-presentation." Government Information Quarterly.

\paragraph{Summary} This paper studies toward a typology of government social media communication: democratic goals, symbolic acts and self-presentation. Since the inception of the Internet, social media has generated profound changes in how individuals and organizations communicate and exchange information.

\paragraph{Context} Since the inception of the Internet, social media has generated profound changes in how individuals and organizations communicate and exchange information. Research on the use of social media by governments first argued these tools would provide government...

\paragraph{Main Argument} Core mechanism maps onto the cluster on participation, responsiveness, and political behavior, but the exact statement should be verified from the full text if heavily used.

\paragraph{Independent Variable(s)} Not stated in standard IV/DV terms; inferred from title/abstract if needed

\paragraph{Dependent Variable(s)} Not stated in standard IV/DV terms; verify if central to argument

\paragraph{Methods and Design} content analysis. Content analysis. Data / sample: To answer RQ1 regarding the applicability of our typology to a sample of government Facebook posts and RQ2 regarding differences of uses across categories and government departments, we carried out a content analysis...

\paragraph{Findings} No clean finding sentence detected; verify from full text if central.

\paragraph{Contribution} Clarifies how state agencies use digital channels for communication and governance rather than treating online politics as purely oppositional.

\paragraph{Limitations} Inference depends on coding choices and text classification; behavior and belief are only indirectly observed.

\paragraph{Relevance to WeChat / My Project} Indirect benchmark: useful for comparison on platform effects or official communication, but needs careful translation to China's WeChat-centered environment.

\paragraph{Related Debate} Whether official digital communication mainly informs, disciplines, mobilizes, or improves responsiveness


\subsection{Meserve and Pemstein 2018: Google Politics The Political Determinants of...}
\textbf{Citation key}: \texttt{meserve\_2018\_psrm}.\\
\textbf{Full citation}: Stephen A. Meserve and Daniel Pemstein 2018. "Google Politics: The Political Determinants of Internet Censorship in Democracies." Political Science Research and Methods.

\paragraph{Summary} This paper studies google politics: the political determinants of internet censorship in democracies. =1682130.

\paragraph{Context} =1682130. Kaufmann, Daniel, Aart Kraay, and Massimo Mastruzzi.

\paragraph{Main Argument} Core mechanism maps onto the cluster on censorship, propaganda, and state influence, but the exact statement should be verified from the full text if heavily used.

\paragraph{Independent Variable(s)} Not stated in standard IV/DV terms; inferred from title/abstract if needed

\paragraph{Dependent Variable(s)} Not stated in standard IV/DV terms; verify if central to argument

\paragraph{Methods and Design} Descriptive or unclear from available parse. Observational or descriptive design; verify details from full text if heavily used. Data / sample: Unclear from available parse; verify from full text if used centrally

\paragraph{Findings} No clean finding sentence detected; verify from full text if central.

\paragraph{Contribution} Provides an external benchmark for mechanisms that may or may not travel to China and WeChat.

\paragraph{Limitations} Design details need closer verification before the paper is used for strong causal claims.

\paragraph{Relevance to WeChat / My Project} Peripheral for a WeChat project; best used as a comparative benchmark rather than core theoretical evidence.

\paragraph{Related Debate} Whether digital control works through persuasion, selective amplification, covert influence, or audience management


\subsection{Repnikova and Fang 2018: Authoritarian Participatory Persuasion 2 0 Netizens...}
\textbf{Citation key}: \texttt{repnikova\_2018\_jcc}.\\
\textbf{Full citation}: Maria Repnikova and Kecheng Fang 2018. "Authoritarian Participatory Persuasion 2.0: Netizens as Thought Work Collaborators in China." Journal of Contemporary China.

\paragraph{Summary} This paper studies authoritarian participatory persuasion 2.0: netizens as thought work collaborators in china. In the reform era, management of information by the Chinese Communist Party has been continuously moving away from explicit, crude tactics of the past toward more subtle and orderly mechanisms of the present.

\paragraph{Context} In the reform era, management of information by the Chinese Communist Party has been continuously moving away from explicit, crude tactics of the past toward more subtle and orderly mechanisms of the present. This study examines one facet of this...

\paragraph{Main Argument} Core mechanism maps onto the cluster on platform affordances and public-private boundaries, but the exact statement should be verified from the full text if heavily used.

\paragraph{Independent Variable(s)} Not stated in standard IV/DV terms; inferred from title/abstract if needed

\paragraph{Dependent Variable(s)} Not stated in standard IV/DV terms; verify if central to argument

\paragraph{Methods and Design} Descriptive or unclear from available parse. Observational or descriptive design; verify details from full text if heavily used. Data / sample: Unclear from available parse; verify from full text if used centrally

\paragraph{Findings} No clean finding sentence detected; verify from full text if central.

\paragraph{Contribution} Helps identify which mechanisms in the Chinese information environment are platform-specific and which generalize beyond a single app.

\paragraph{Limitations} Design details need closer verification before the paper is used for strong causal claims.

\paragraph{Relevance to WeChat / My Project} Conceptually useful for a WeChat project because it speaks to the Chinese information environment, especially the cluster on platform affordances and public-private boundaries, even if the platform is not WeChat.

\paragraph{Related Debate} How interface design changes the visibility and meaning of digital expression


\subsection{Barbera et al. 2019: Who Leads Who Follows Measuring Issue...}
\textbf{Citation key}: \texttt{barbera\_2019\_apsr}.\\
\textbf{Full citation}: Pablo Barbera, Andreu Casas, Jonathan Nagler, Patrick J. Egan, Richard Bonneau, John T. Jost, and Joshua A. Tucker 2019. "Who Leads? Who Follows? Measuring Issue Attention and Agenda Setting by Legislators and the Mass Public Using Social Media Data." American Political Science Review.

\paragraph{Summary} This paper studies who leads? who follows? measuring issue attention and agenda setting by legislators and the mass public using social.... =3395307.

\paragraph{Context} =3395307. Caughey, Devin, and Christopher Warshaw.

\paragraph{Main Argument} Core mechanism maps onto the cluster on participation, responsiveness, and political behavior, but the exact statement should be verified from the full text if heavily used.

\paragraph{Independent Variable(s)} News media Exposure

\paragraph{Dependent Variable(s)} Political Knowledge and Participation

\paragraph{Methods and Design} topic modeling. Topic modeling on platform texts. Data / sample: Survey Evidence from Local Politicians in the United States." Working paper. Egan,... that classifies tweets sent by legislators and citizens into topics, we use vector autoregression models to explore whose...

\paragraph{Findings} No clean finding sentence detected; verify from full text if central.

\paragraph{Contribution} Provides an external benchmark for mechanisms that may or may not travel to China and WeChat.

\paragraph{Limitations} Inference depends on coding choices and text classification; behavior and belief are only indirectly observed.

\paragraph{Relevance to WeChat / My Project} Indirect benchmark: useful for comparison on platform effects or official communication, but needs careful translation to China's WeChat-centered environment.

\paragraph{Related Debate} Whether official digital communication mainly informs, disciplines, mobilizes, or improves responsiveness


\subsection{Baum and Potter 2019: Media Public Opinion and Foreign Policy...}
\textbf{Citation key}: \texttt{baum\_2019\_jop}.\\
\textbf{Full citation}: Matthew A. Baum and Philip B. K. Potter 2019. "Media, Public Opinion, and Foreign Policy in the Age of Social Media." The Journal of Politics.

\paragraph{Summary} This paper studies media, public opinion, and foreign policy in the age of social media. FOREIGN POLICY IN THE AGE OF TRUMP SYMPOSIUM Media, Public Opinion, and Foreign Policy in the Age of Social Media Matthew A.

\paragraph{Context} FOREIGN POLICY IN THE AGE OF TRUMP SYMPOSIUM Media, Public Opinion, and Foreign Policy in the Age of Social Media Matthew A. Baum, Harvard University Philip B.

\paragraph{Main Argument} Core mechanism maps onto the cluster on interpersonal discussion, networks, and opinion formation, but the exact statement should be verified from the full text if heavily used.

\paragraph{Independent Variable(s)} Not stated in standard IV/DV terms; inferred from title/abstract if needed

\paragraph{Dependent Variable(s)} Not stated in standard IV/DV terms; verify if central to argument

\paragraph{Methods and Design} Descriptive or unclear from available parse. Observational or descriptive design; verify details from full text if heavily used. Data / sample: Unclear from available parse; verify from full text if used centrally

\paragraph{Findings} No clean finding sentence detected; verify from full text if central.

\paragraph{Contribution} Provides an external benchmark for mechanisms that may or may not travel to China and WeChat.

\paragraph{Limitations} Design details need closer verification before the paper is used for strong causal claims.

\paragraph{Relevance to WeChat / My Project} Indirect benchmark: useful for comparison on platform effects or official communication, but needs careful translation to China's WeChat-centered environment.

\paragraph{Related Debate} How social ties and discussion networks mediate willingness to speak, share, and update opinions


\subsection{Deibert 2019: Three Painful Truths About Social Media}
\textbf{Citation key}: \texttt{deibert\_2019\_journal\_of\_democracy}.\\
\textbf{Full citation}: Ronald J. Deibert 2019. "Three Painful Truths About Social Media." Journal of Democracy.

\paragraph{Summary} This paper studies three painful truths about social media. The Road to Digital Unfreedom: Three Painful Truths About Social Media Ronald J.

\paragraph{Context} 25-39 (Article) Published by Johns Hopkins University Press DOI: https://doi.org/10.1353/jod.2019.0002 For additional information about this article https://muse.jhu.edu/article/713720 [76.176.202.166] Project MUSE (2026-03-27 14:51 GMT) University of...

\paragraph{Main Argument} Core mechanism maps onto the cluster on comparative digital media benchmarks, but the exact statement should be verified from the full text if heavily used.

\paragraph{Independent Variable(s)} Not stated in standard IV/DV terms; inferred from title/abstract if needed

\paragraph{Dependent Variable(s)} Not stated in standard IV/DV terms; verify if central to argument

\paragraph{Methods and Design} book / theory. Theoretical or historical argument rather than a single empirical design. Data / sample: Unclear from available parse; verify from full text if used centrally

\paragraph{Findings} No clean finding sentence detected; verify from full text if central.

\paragraph{Contribution} Provides an external benchmark for mechanisms that may or may not travel to China and WeChat.

\paragraph{Limitations} Useful for theory-building, but not designed to isolate a single platform-level causal effect.

\paragraph{Relevance to WeChat / My Project} Indirect benchmark: useful for comparison on platform effects or official communication, but needs careful translation to China's WeChat-centered environment.

\paragraph{Related Debate} How findings from facebook and comparative / general settings travel to a WeChat-centered China project


\subsection{Guriev and Treisman 2019: Informational Autocrats}
\textbf{Citation key}: \texttt{guriev\_2019\_journal\_of\_economic\_perspectives}.\\
\textbf{Full citation}: Sergei Guriev and Daniel Treisman 2019. "Informational Autocrats." Journal of Economic Perspectives.

\paragraph{Summary} This paper studies informational autocrats. =1449034.

\paragraph{Context} =1449034. Cruz, Cesi, Philip Keefer, and Carlos Scartascini.

\paragraph{Main Argument} Core mechanism maps onto the cluster on censorship, propaganda, and state influence, but the exact statement should be verified from the full text if heavily used.

\paragraph{Independent Variable(s)} Not stated in standard IV/DV terms; inferred from title/abstract if needed

\paragraph{Dependent Variable(s)} Not stated in standard IV/DV terms; verify if central to argument

\paragraph{Methods and Design} panel / observational analysis. Observational or descriptive design; verify details from full text if heavily used. Data / sample: "Why Resource-Poor Dictators Allow Freer Media: A Theory and Evidence from Panel Data." American Political Science Review 103 (4): 645--68. Gehlbach,... of rule, which they can combine with various messages.

\paragraph{Findings} No clean finding sentence detected; verify from full text if central.

\paragraph{Contribution} Provides an external benchmark for mechanisms that may or may not travel to China and WeChat.

\paragraph{Limitations} Mechanisms may not transfer cleanly to WeChat's hybrid private-public architecture without further adaptation.

\paragraph{Relevance to WeChat / My Project} Peripheral for a WeChat project; best used as a comparative benchmark rather than core theoretical evidence.

\paragraph{Related Debate} Whether digital control works through persuasion, selective amplification, covert influence, or audience management


\subsection{Huang et al. 2019: Media Protest Diffusion and Authoritarian Resilience}
\textbf{Citation key}: \texttt{huang\_2019\_psrm}.\\
\textbf{Full citation}: Haifeng Huang, Serra Boranbay-Akan, and Ling Huang 2019. "Media, Protest Diffusion, and Authoritarian Resilience." Political Science Research and Methods.

\paragraph{Summary} This paper studies media, protest diffusion, and authoritarian resilience. , is that local protests will remain localized unless the national government ignores popular grievances.

\paragraph{Context} If the government fears, however, that there is a substantial risk that the initially localized protests may turn into a broader based, cross-regional movement challenging the regime, the government's attitudes toward the media will be more conflicted.

\paragraph{Main Argument} Core mechanism maps onto the cluster on comparative digital media benchmarks, but the exact statement should be verified from the full text if heavily used.

\paragraph{Independent Variable(s)} Not stated in standard IV/DV terms; inferred from title/abstract if needed

\paragraph{Dependent Variable(s)} Not stated in standard IV/DV terms; verify if central to argument

\paragraph{Methods and Design} Descriptive or unclear from available parse. Observational or descriptive design; verify details from full text if heavily used. Data / sample: Unclear from available parse; verify from full text if used centrally

\paragraph{Findings} No clean finding sentence detected; verify from full text if central.

\paragraph{Contribution} Helps identify which mechanisms in the Chinese information environment are platform-specific and which generalize beyond a single app.

\paragraph{Limitations} Design details need closer verification before the paper is used for strong causal claims.

\paragraph{Relevance to WeChat / My Project} Conceptually useful for a WeChat project because it speaks to the Chinese information environment, especially the cluster on comparative digital media benchmarks, even if the platform is not WeChat.

\paragraph{Related Debate} How findings from weibo and china settings travel to a WeChat-centered China project


\subsection{Levy and Mattsson 2019: The Effects of Social Movements Evidence...}
\textbf{Citation key}: \texttt{levy\_2019\_wp}.\\
\textbf{Full citation}: Roee Levy and Martin Mattsson 2019. "The Effects of Social Movements: Evidence from \#MeToo." SSRN Electronic Journal.

\paragraph{Summary} This study examines the effects of social movements: evidence from \#metoo using survey evidence and reports that using a differencein-differences strategy comparing sex crimes and non-sex crimes before and after the MeToo movement started, we find that the movement increased the number of sex crimes reported by...

\paragraph{Context} Social movements are associated with large societal changes, but evidence of their causal effects is limited. We study the effect of the MeToo movement on an important personal decision---reporting a sex crime to the police.

\paragraph{Main Argument} The paper's core mechanism is that using a differencein-differences strategy comparing sex crimes and non-sex crimes before and after the MeToo movement started, we find that the movement increased the number of sex crimes reported by approximately 10\% and...

\paragraph{Independent Variable(s)} Social Movements: Evidence from \#MeToo. Social movements are associated with large societal changes, but evidence of...

\paragraph{Dependent Variable(s)} an important personal decision---reporting a sex crime to the police

\paragraph{Methods and Design} survey. Survey-based observational analysis. Data / sample: Using a differencein-differences strategy comparing sex crimes and non-sex crimes before and after the MeToo movement started, we find that the movement increased the number of sex crimes reported by approximately...

\paragraph{Findings} Using a differencein-differences strategy comparing sex crimes and non-sex crimes before and after the MeToo movement started, we find that the movement increased the number of sex crimes reported by approximately 10\% and that the effect persisted until the end of our data, 15-27 months after the movement started. Using detailed US data, we show that the MeToo movement not only increased reporting but also increased arrests for sexual assaults; and that in contrast to a common criticism of the movement, the effects are similar across racial and socioeconomic groups.

\paragraph{Contribution} Provides an external benchmark for mechanisms that may or may not travel to China and WeChat.

\paragraph{Limitations} Mechanisms may not transfer cleanly to WeChat's hybrid private-public architecture without further adaptation.

\paragraph{Relevance to WeChat / My Project} Peripheral for a WeChat project; best used as a comparative benchmark rather than core theoretical evidence.

\paragraph{Related Debate} How findings from not platform-specific and comparative / general settings travel to a WeChat-centered China project


\subsection{Munger 2019: Social Media Political Science and Democracy}
\textbf{Citation key}: \texttt{munger\_2019\_jop}.\\
\textbf{Full citation}: Kevin Munger 2019. "Social Media, Political Science, and Democracy." The Journal of Politics.

\paragraph{Summary} This paper studies social media, political science, and democracy. REVIEW ESSAY Social Media, Political Science, and Democracy Kevin Munger, Pennsylvania State University Frenemies: How Social Media Polarizes America.

\paragraph{Context} REVIEW ESSAY Social Media, Political Science, and Democracy Kevin Munger, Pennsylvania State University Frenemies: How Social Media Polarizes America. By Jaime Settle.

\paragraph{Main Argument} Core mechanism maps onto the cluster on comparative digital media benchmarks, but the exact statement should be verified from the full text if heavily used.

\paragraph{Independent Variable(s)} Not stated in standard IV/DV terms; inferred from title/abstract if needed

\paragraph{Dependent Variable(s)} Not stated in standard IV/DV terms; verify if central to argument

\paragraph{Methods and Design} book / theory. Theoretical or historical argument rather than a single empirical design. Data / sample: Cambridge: Cambridge University Press, 2018. and produce peer-reviewed journal articles to have a chance at a tenure-track job.

\paragraph{Findings} No clean finding sentence detected; verify from full text if central.

\paragraph{Contribution} Provides an external benchmark for mechanisms that may or may not travel to China and WeChat.

\paragraph{Limitations} Useful for theory-building, but not designed to isolate a single platform-level causal effect.

\paragraph{Relevance to WeChat / My Project} Indirect benchmark: useful for comparison on platform effects or official communication, but needs careful translation to China's WeChat-centered environment.

\paragraph{Related Debate} How findings from facebook and united states settings travel to a WeChat-centered China project


\subsection{Munger et al. 2019: Elites Tweet to Get Feet Off...}
\textbf{Citation key}: \texttt{munger\_2019\_psrm}.\\
\textbf{Full citation}: Kevin Munger, Richard Bonneau, Jonathan Nagler, and Joshua A. Tucker 2019. "Elites Tweet to Get Feet Off the Streets: Measuring Regime Social Media Strategies During Protest." Political Science Research and Methods.

\paragraph{Summary} This paper studies elites tweet to get feet off the streets: measuring regime social media strategies during protest. =2529769, accessed 11 August 2015.

\paragraph{Context} =2529769, accessed 11 August 2015. Ciccariello-Maher, Georeg.

\paragraph{Main Argument} Core mechanism maps onto the cluster on comparative digital media benchmarks, but the exact statement should be verified from the full text if heavily used.

\paragraph{Independent Variable(s)} Not stated in standard IV/DV terms; inferred from title/abstract if needed

\paragraph{Dependent Variable(s)} Not stated in standard IV/DV terms; verify if central to argument

\paragraph{Methods and Design} topic modeling. Topic modeling on platform texts. Data / sample: Available at https://www.foreignaffairs. com/articles/venezuela/2013-01-04/chavismo-after-ch-vez, accessed 1 September 2015. https://doi.org/10.1017/psrm.2018.3 Published online by Cambridge University Press Elites...

\paragraph{Findings} No clean finding sentence detected; verify from full text if central.

\paragraph{Contribution} Provides an external benchmark for mechanisms that may or may not travel to China and WeChat.

\paragraph{Limitations} Inference depends on coding choices and text classification; behavior and belief are only indirectly observed.

\paragraph{Relevance to WeChat / My Project} Indirect benchmark: useful for comparison on platform effects or official communication, but needs careful translation to China's WeChat-centered environment.

\paragraph{Related Debate} How findings from twitter / x and comparative / general settings travel to a WeChat-centered China project


\subsection{Allcott et al. 2020: The Welfare Effects of Social Media}
\textbf{Citation key}: \texttt{allcott\_2020\_aer}.\\
\textbf{Full citation}: Hunt Allcott, Luca Braghieri, Sarah Eichmeyer, and Matthew Gentzkow 2020. "The Welfare Effects of Social Media." American Economic Review.

\paragraph{Summary} This paper studies the welfare effects of social media. =3148805.

\paragraph{Context} =3148805. Brynjolfsson, Erik, Felix Eggers, and Avinash Gannamaneni.

\paragraph{Main Argument} Core mechanism maps onto the cluster on information exposure, selective exposure, and persuasion, but the exact statement should be verified from the full text if heavily used.

\paragraph{Independent Variable(s)} Social Media. =3148805. Brynjolfsson, Erik, Felix Eggers, and Avinash Gannamaneni. 2018. "Using Massive Online...

\paragraph{Dependent Variable(s)} the Internet

\paragraph{Methods and Design} survey. Survey-based observational analysis. Data / sample: "Social Capital on Facebook: Differentiating Uses and Users." CHI '11: Proceedings of the SIGCHI Conference on Human Factors in Computing Systems: 571--80. The experiment has eight parts: recruitment, pre-screen,...

\paragraph{Findings} No clean finding sentence detected; verify from full text if central.

\paragraph{Contribution} Provides an external benchmark for mechanisms that may or may not travel to China and WeChat.

\paragraph{Limitations} Design details need closer verification before the paper is used for strong causal claims.

\paragraph{Relevance to WeChat / My Project} Indirect benchmark: useful for comparison on platform effects or official communication, but needs careful translation to China's WeChat-centered environment.

\paragraph{Related Debate} Whether exposure changes attitudes directly or is filtered by selection, framing, and prior preferences


\subsection{Clarke and Kocak 2020: Garrett that more attention to context...}
\textbf{Citation key}: \texttt{clarke\_2020\_bjps}.\\
\textbf{Full citation}: Killian Clarke and Korhan Kocak 2020. "Garrett that more attention to context might help scholars to parse the potentially diverse effects." British Journal of Political Science.

\paragraph{Summary} This paper studies garrett that more attention to context might help scholars to parse the potentially diverse effects. Drawing on evidence from the 2011 Egyptian uprising, this article demonstrates how the use of two social media platforms -- Facebook and Twitter -- contributed to a discrete mobilizational outcome:...

\paragraph{Context} Drawing on evidence from the 2011 Egyptian uprising, this article demonstrates how the use of two social media platforms -- Facebook and Twitter -- contributed to a discrete mobilizational outcome: the staging of a successful first protest in a revolutionary...

\paragraph{Main Argument} Core mechanism maps onto the cluster on comparative digital media benchmarks, but the exact statement should be verified from the full text if heavily used.

\paragraph{Independent Variable(s)} Not stated in standard IV/DV terms; inferred from title/abstract if needed

\paragraph{Dependent Variable(s)} Not stated in standard IV/DV terms; verify if central to argument

\paragraph{Methods and Design} interviews / focus groups. Observational or descriptive design; verify details from full text if heavily used. Data / sample: It draws on qualitative and quantitative evidence, including interviews, social media data and surveys, to analyze three mechanisms that linked these platforms to the success of the January 25 protest: (1) protester...

\paragraph{Findings} No clean finding sentence detected; verify from full text if central.

\paragraph{Contribution} Provides an external benchmark for mechanisms that may or may not travel to China and WeChat.

\paragraph{Limitations} Design details need closer verification before the paper is used for strong causal claims.

\paragraph{Relevance to WeChat / My Project} Peripheral for a WeChat project; best used as a comparative benchmark rather than core theoretical evidence.

\paragraph{Related Debate} How findings from facebook and united kingdom settings travel to a WeChat-centered China project


\subsection{Enikolopov et al. 2020: Social Media and Protest Participation Evidence...}
\textbf{Citation key}: \texttt{enikolopov\_2020\_ecma}.\\
\textbf{Full citation}: Ruben Enikolopov, Alexey Makarin, and Maria Petrova 2020. "Social Media and Protest Participation: Evidence From Russia." Econometrica.

\paragraph{Summary} This paper studies social media and protest participation: evidence from russia. COLLECTIVE ACTION PROBLEM has traditionally been seen as one of the major barriers to achieving socially beneficial outcomes (e.g., Olson (1965), Hardin (1982), Ostrom (1990)).

\paragraph{Context} COLLECTIVE ACTION PROBLEM has traditionally been seen as one of the major barriers to achieving socially beneficial outcomes (e.g., Olson (1965), Hardin (1982), Ostrom (1990)). In addition to the classic issue of free-riding, a group's ability to overcome...

\paragraph{Main Argument} Core mechanism maps onto the cluster on interpersonal discussion, networks, and opinion formation, but the exact statement should be verified from the full text if heavily used.

\paragraph{Independent Variable(s)} Not stated in standard IV/DV terms; inferred from title/abstract if needed

\paragraph{Dependent Variable(s)} Not stated in standard IV/DV terms; verify if central to argument

\paragraph{Methods and Design} Descriptive or unclear from available parse. Observational or descriptive design; verify details from full text if heavily used. Data / sample: New horizontal information exchange technologies, such as Facebook and Twitter, allow users to converse directly without intermediaries at a very low cost, thus potentially enhancing the spread of information and...

\paragraph{Findings} No clean finding sentence detected; verify from full text if central.

\paragraph{Contribution} Provides an external benchmark for mechanisms that may or may not travel to China and WeChat.

\paragraph{Limitations} Design details need closer verification before the paper is used for strong causal claims.

\paragraph{Relevance to WeChat / My Project} Indirect benchmark: useful for comparison on platform effects or official communication, but needs careful translation to China's WeChat-centered environment.

\paragraph{Related Debate} How social ties and discussion networks mediate willingness to speak, share, and update opinions


\subsection{Gohdes 2020: Repression Technology Internet Accessibility and State...}
\textbf{Citation key}: \texttt{gohdes\_2020\_ajps}.\\
\textbf{Full citation}: Anita R. Gohdes 2020. "Repression Technology: Internet Accessibility and State Violence." American Journal of Political Science.

\paragraph{Summary} This study examines repression technology: internet accessibility and state violence using descriptive or unclear from available parse evidence and reports that i find that higher levels of Internet accessibility are associated with increases in targeted repression, whereas areas with limited access experience more...

\paragraph{Context} This article offers a first subnational analysis of the relationship between states' dynamic control of Internet access and their use of violent repression.

\paragraph{Main Argument} The paper's core mechanism is that i find that higher levels of Internet accessibility are associated with increases in targeted repression, whereas areas with limited access experience more indiscriminate campaigns of violence. Even before unrest becomes...

\paragraph{Independent Variable(s)} Not stated in standard IV/DV terms; inferred from title/abstract if needed

\paragraph{Dependent Variable(s)} Not stated in standard IV/DV terms; verify if central to argument

\paragraph{Methods and Design} Descriptive or unclear from available parse. Observational or descriptive design; verify details from full text if heavily used. Data / sample: I present new data on killings in the Syrian conflict that distinguish between targeted and untargeted events, using supervised text classification.

\paragraph{Findings} I find that higher levels of Internet accessibility are associated with increases in targeted repression, whereas areas with limited access experience more indiscriminate campaigns of violence. Even before unrest becomes visible, citizens who show signs of opposing viewpoints are now liable to be placed under close surveillance, long before their political preferences become publicly known.

\paragraph{Contribution} Provides an external benchmark for mechanisms that may or may not travel to China and WeChat.

\paragraph{Limitations} Design details need closer verification before the paper is used for strong causal claims.

\paragraph{Relevance to WeChat / My Project} Peripheral for a WeChat project; best used as a comparative benchmark rather than core theoretical evidence.

\paragraph{Related Debate} How social ties and discussion networks mediate willingness to speak, share, and update opinions


\subsection{Graham 2020: Self-Awareness of Political Knowledge}
\textbf{Citation key}: \texttt{graham\_2020\_polbeh}.\\
\textbf{Full citation}: Matthew H. Graham 2020. "Self-Awareness of Political Knowledge." Political Behavior.

\paragraph{Summary} This paper studies self-awareness of political knowledge. Despite widespread concern over false beliefs about politically-relevant facts, little is known about how strongly Americans believe their answers to poll questions.

\paragraph{Context} Despite widespread concern over false beliefs about politically-relevant facts, little is known about how strongly Americans believe their answers to poll questions. I propose a conceptual framework for characterizing survey responses about facts:...

\paragraph{Main Argument} I propose a conceptual framework for characterizing survey responses about facts: self-awareness, or how well people can assess their own knowledge.

\paragraph{Independent Variable(s)} Not stated in standard IV/DV terms; inferred from title/abstract if needed

\paragraph{Dependent Variable(s)} Not stated in standard IV/DV terms; verify if central to argument

\paragraph{Methods and Design} survey. Survey-based observational analysis. Data / sample: I propose a conceptual framework for characterizing survey responses about facts: self-awareness, or how well people can assess their own knowledge. Even on "unfavorable" facts that are inconvenient to the...

\paragraph{Findings} No clean finding sentence detected; verify from full text if central.

\paragraph{Contribution} Provides an external benchmark for mechanisms that may or may not travel to China and WeChat.

\paragraph{Limitations} Design details need closer verification before the paper is used for strong causal claims.

\paragraph{Relevance to WeChat / My Project} Peripheral for a WeChat project; best used as a comparative benchmark rather than core theoretical evidence.

\paragraph{Related Debate} How social ties and discussion networks mediate willingness to speak, share, and update opinions


\subsection{SOBOLEV et al. 2020: News and Geolocated Social Media Accurately...}
\textbf{Citation key}: \texttt{sobolev\_2020\_apsr}.\\
\textbf{Full citation}: ANTON SOBOLEV, M. KEITH CHEN, JUNGSEOCK JOO, and ZACHARY C. STEINERT-THRELKELD 2020. "News and Geolocated Social Media Accurately Measure Protest Size Variation." American Political Science Review.

\paragraph{Summary} This study examines news and geolocated social media accurately measure protest size variation using descriptive or unclear from available parse evidence and reports that n https://doi.org/10.1017/S0003055420000295 Published online by Cambridge University Press ews and social media suggest that news sources provide...

\paragraph{Context} N https://doi.org/10.1017/S0003055420000295 Published online by Cambridge University Press ews and social media

\paragraph{Main Argument} The paper's core mechanism is that n https://doi.org/10.1017/S0003055420000295 Published online by Cambridge University Press ews and social media suggest that news sources provide accurate estimates of protest size variation, as do social media text or...

\paragraph{Independent Variable(s)} Not stated in standard IV/DV terms; inferred from title/abstract if needed

\paragraph{Dependent Variable(s)} Not stated in standard IV/DV terms; verify if central to argument

\paragraph{Methods and Design} Descriptive or unclear from available parse. Observational or descriptive design; verify details from full text if heavily used. Data / sample: Ryan Fox Squire and Aaron Hoffman assisted with the SafeGraph data. This work is supported by NSF RIDIR \#1831848, the UCLA Big Data Transdisciplinary Seed Grant, and a UCLA California Center for Population Research...

\paragraph{Findings} N https://doi.org/10.1017/S0003055420000295 Published online by Cambridge University Press ews and social media suggest that news sources provide accurate estimates of protest size variation, as do social media text or images of protesters (Botta, Moat, 1349 Anton Sobolev et al. and Preis 2015). Transforming, and not reporting, raw data throws away information future researchers may find useful.

\paragraph{Contribution} Provides an external benchmark for mechanisms that may or may not travel to China and WeChat.

\paragraph{Limitations} Design details need closer verification before the paper is used for strong causal claims.

\paragraph{Relevance to WeChat / My Project} Indirect benchmark: useful for comparison on platform effects or official communication, but needs careful translation to China's WeChat-centered environment.

\paragraph{Related Debate} How findings from social media (general) and united states settings travel to a WeChat-centered China project


\subsection{Tai and Fu 2020: Specificity Conflict and Focal Point A...}
\textbf{Citation key}: \texttt{tai\_2020\_journal\_of\_communication}.\\
\textbf{Full citation}: Yun Tai and King-wa Fu 2020. "Specificity, Conflict, and Focal Point: A Systematic Investigation into Social Media Censorship in China." Journal of Communication.

\paragraph{Summary} This study examines specificity, conflict, and focal point: a systematic investigation into social media censorship in china using descriptive or unclear from available parse evidence and reports that with the effects of account attributes and article topics excluded, we find that article specificity raises the...

\paragraph{Context} Journal of Communication ISSN 0021-9916 Specificity, Conflict, and Focal Point: A Systematic Investigation into Social Media Censorship in China Yun Tai 1 2 1 \& King-wa Fu 2 Department of Mass Communication, Tamkang University, Taiwan Journalism and Media...

\paragraph{Main Argument} The paper's core mechanism is that with the effects of account attributes and article topics excluded, we find that article specificity raises the odds of being censored. [press release].

\paragraph{Independent Variable(s)} Not stated in standard IV/DV terms; inferred from title/abstract if needed

\paragraph{Dependent Variable(s)} Not stated in standard IV/DV terms; verify if central to argument

\paragraph{Methods and Design} Descriptive or unclear from available parse. Observational or descriptive design; verify details from full text if heavily used. Data / sample: This study interrogates postpublication censorship on Chinese social media by examining the differences between 2,280 pairs of censored WeChat articles and matched remaining articles. With the effects of account...

\paragraph{Findings} With the effects of account attributes and article topics excluded, we find that article specificity raises the odds of being censored. [press release].

\paragraph{Contribution} Directly extends the WeChat-centered literature by specifying platform-specific behavior rather than treating social media as generic.

\paragraph{Limitations} Design details need closer verification before the paper is used for strong causal claims.

\paragraph{Relevance to WeChat / My Project} Directly relevant: this paper studies WeChat behavior or institutions and can be used to theorize visibility, audience composition, and official-account interaction.

\paragraph{Related Debate} Whether semi-private WeChat interactions should be read as information use, endorsement, or strategically managed visibility


\subsection{Wang 2020: Book review Tencent The Political Economy...}
\textbf{Citation key}: \texttt{wang\_2020\_global\_media\_and\_communication}.\\
\textbf{Full citation}: Grace Yuehan Wang 2020. "Book review: Tencent: The Political Economy of China's Surging Internet Giant." Global Media and Communication.

\paragraph{Summary} This study examines book review: tencent: the political economy of china's surging internet giant using case study evidence and reports that 113 trace the histories of these three types of entities and to analyze the politicaleconomic contexts that have enabled and conditioned their development would demonstrate...

\paragraph{Context}  1 1 History and Context 9 2 Economic Profile 32 3 Political Profile 62 4 Cultural Profile 83 Conclusion 108 Bibliography114 Index121 Abbreviations CNNIC FDI ICT IM IPO IVAS MEI MIH MII MIIT MMOG MPT MVAS PCCW SEZ VAS VC VIE China Internet Network...

\paragraph{Main Argument} The paper's core mechanism is that 113 trace the histories of these three types of entities and to analyze the politicaleconomic contexts that have enabled and conditioned their development would demonstrate how the participation of venture capital...

\paragraph{Independent Variable(s)} Not stated in standard IV/DV terms; inferred from title/abstract if needed

\paragraph{Dependent Variable(s)} Not stated in standard IV/DV terms; verify if central to argument

\paragraph{Methods and Design} case study. Case-study design. Data / sample:  1 1 History and Context 9 2 Economic Profile 32 3 Political Profile 62 4 Cultural Profile 83 Conclusion 108 Bibliography114 Index121 Abbreviations CNNIC FDI ICT IM IPO IVAS MEI MIH MII MIIT MMOG MPT MVAS...

\paragraph{Findings} 113 trace the histories of these three types of entities and to analyze the politicaleconomic contexts that have enabled and conditioned their development would demonstrate how the participation of venture capital investments has further consolidated and accelerated the private actors' dominance over the provision system of global communication and information.

\paragraph{Contribution} Directly extends the WeChat-centered literature by specifying platform-specific behavior rather than treating social media as generic.

\paragraph{Limitations} Limited external validity and weaker causal leverage because the design is case-based and interpretive.

\paragraph{Relevance to WeChat / My Project} Directly relevant: this paper studies WeChat behavior or institutions and can be used to theorize visibility, audience composition, and official-account interaction.

\paragraph{Related Debate} Whether semi-private WeChat interactions should be read as information use, endorsement, or strategically managed visibility


\subsection{Eck et al. 2021: Evade and Deceive Citizen Responses to...}
\textbf{Citation key}: \texttt{eck\_2021\_jop}.\\
\textbf{Full citation}: Kristine Eck, Sophia Hatz, Charles Crabtree, and Atsushi Tago 2021. "Evade and Deceive? Citizen Responses to Surveillance." The Journal of Politics.

\paragraph{Summary} This study examines evade and deceive? citizen responses to surveillance using survey experiment evidence and reports that third, the findings indicate the utility of considering a wider array of evasive behaviors than has been previously considered in the literature, which has typically truncated its examination of...

\paragraph{Context} This article theorizes about different citizen responses to surveillance that fall on what we term the evasion-deception spectrum, including preference falsification, self-censorship, and opting out.

\paragraph{Main Argument} The paper's core mechanism is that third, the findings indicate the utility of considering a wider array of evasive behaviors than has been previously considered in the literature, which has typically truncated its examination of this spectrum to the most...

\paragraph{Independent Variable(s)} Author inference from title/abstract: Evade and Deceive? Citizen Responses to Surveillance

\paragraph{Dependent Variable(s)} Not stated in standard IV/DV terms

\paragraph{Methods and Design} survey experiment. survey experiment. Data / sample: Specifically, we design an online public opinion survey experiment that measures three types of responses to surveillance: preference falsification, self-censorship, and opting out. We incorporate an experimental...

\paragraph{Findings} Third, the findings indicate the utility of considering a wider array of evasive behaviors than has been previously considered in the literature, which has typically truncated its examination of this spectrum to the most overt forms.

\paragraph{Contribution} Offers comparatively strong leverage on causal mechanisms in digital political communication.

\paragraph{Limitations} Internal validity is strong, but external validity may depend on whether experimental treatments resemble real platform use.

\paragraph{Relevance to WeChat / My Project} Peripheral for a WeChat project; best used as a comparative benchmark rather than core theoretical evidence.

\paragraph{Related Debate} Whether digital control works through persuasion, selective amplification, covert influence, or audience management


\subsection{Foos et al. 2021: Does Social Media Promote Civic Activism...}
\textbf{Citation key}: \texttt{foos\_2021\_psrm}.\\
\textbf{Full citation}: Florian Foos, Lyubomir Kostadinov, Nikolay Marinov, and Frank Schimmelfennig 2021. "Does Social Media Promote Civic Activism? A Field Experiment with a Civic Campaign." Political Science Research and Methods.

\paragraph{Summary} This study examines does social media promote civic activism? a field experiment with a civic campaign using field experiment evidence and reports that as expected, we find that Bulgarian Facebook users who are active in pro-environmental groups, and those who decide to follow the campaign, are more highly educated...

\paragraph{Context} However, the extent to which social media campaigns simply recruit like-minded individuals as compared to exerting a causal impact on joiners' attitudes is difficult to disentangle.

\paragraph{Main Argument} The paper's core mechanism is that as expected, we find that Bulgarian Facebook users who are active in pro-environmental groups, and those who decide to follow the campaign, are more highly educated than those who decide to stay at the sidelines. In...

\paragraph{Independent Variable(s)} Author inference from title/abstract: Does Social Media Promote Civic Activism? A Field Experiment with a Civic Campaign

\paragraph{Dependent Variable(s)} Not stated in standard IV/DV terms

\paragraph{Methods and Design} field experiment. field experiment. Data / sample: As expected, we find that Bulgarian Facebook users who are active in pro-environmental groups, and those who decide to follow the campaign, are more highly educated than those who decide to stay at the sidelines. As...

\paragraph{Findings} As expected, we find that Bulgarian Facebook users who are active in pro-environmental groups, and those who decide to follow the campaign, are more highly educated than those who decide to stay at the sidelines. In contrast, we find little evidence that the campaign affected opinions, knowledge, or self-reported behavior.

\paragraph{Contribution} Offers comparatively strong leverage on causal mechanisms in digital political communication.

\paragraph{Limitations} Internal validity is strong, but external validity may depend on whether experimental treatments resemble real platform use.

\paragraph{Relevance to WeChat / My Project} Indirect benchmark: useful for comparison on platform effects or official communication, but needs careful translation to China's WeChat-centered environment.

\paragraph{Related Debate} How social ties and discussion networks mediate willingness to speak, share, and update opinions


\subsection{Graham and Coppock 2021: Asking About Attitude Change}
\textbf{Citation key}: \texttt{graham\_2021\_poq}.\\
\textbf{Full citation}: Matthew H Graham and Alexander Coppock 2021. "Asking About Attitude Change." Public Opinion Quarterly.

\paragraph{Summary} This study examines asking about attitude change using survey evidence and reports that we show that this type of question (the change format) exhibits poor measurement properties, in large part because subjects engage in response substitution. Using a series of experiments embedded in four studies, we show that...

\paragraph{Context} Surveys often ask respondents how information or events changed their attitudes. Does [information X] make you more or less supportive of [policy Y]?

\paragraph{Main Argument} The paper's core mechanism is that we show that this type of question (the change format) exhibits poor measurement properties, in large part because subjects engage in response substitution. Using a series of experiments embedded in four studies, we show...

\paragraph{Independent Variable(s)} sexual misconduct allegations

\paragraph{Dependent Variable(s)} support for Republican Roy Moore

\paragraph{Methods and Design} survey. Survey-based observational analysis. Data / sample: Surveys often ask respondents how information or events changed their attitudes. In advance of Alabama's 2017 special election for US Senator, polling firm JMC Analytics released a survey that sought to estimate the...

\paragraph{Findings} We show that this type of question (the change format) exhibits poor measurement properties, in large part because subjects engage in response substitution. Using a series of experiments embedded in four studies, we show that the counterfactual format greatly reduces bias relative to the change format.

\paragraph{Contribution} Provides an external benchmark for mechanisms that may or may not travel to China and WeChat.

\paragraph{Limitations} Mechanisms may not transfer cleanly to WeChat's hybrid private-public architecture without further adaptation.

\paragraph{Relevance to WeChat / My Project} Peripheral for a WeChat project; best used as a comparative benchmark rather than core theoretical evidence.

\paragraph{Related Debate} How findings from not platform-specific and united states settings travel to a WeChat-centered China project


\subsection{Guo et al. 2021: Why do citizens participate on government...}
\textbf{Citation key}: \texttt{guo\_2021\_information\_management}.\\
\textbf{Full citation}: Junpeng Guo, Na Liu, Yi Wu, and Chunxin Zhang 2021. "Why do citizens participate on government social media accounts during crises? A civic voluntarism perspective." Information \& Management.

\paragraph{Summary} This paper studies why do citizens participate on government social media accounts during crises? a civic voluntarism perspective. Various natural and man-made disasters, such as earthquakes, tsu namis, and chemical explosions, pose an ever-present challenge of crisis management.

\paragraph{Context} Various natural and man-made disasters, such as earthquakes, tsu namis, and chemical explosions, pose an ever-present challenge of crisis management. To cope with these crises, the dissemination of relevant information to all crisis response organizations...

\paragraph{Main Argument} Core mechanism maps onto the cluster on participation, responsiveness, and political behavior, but the exact statement should be verified from the full text if heavily used.

\paragraph{Independent Variable(s)} Author inference from title/abstract: Why do citizens participate on government social media accounts during crises? A civic voluntarism perspective

\paragraph{Dependent Variable(s)} Not stated in standard IV/DV terms

\paragraph{Methods and Design} Descriptive or unclear from available parse. Observational or descriptive design; verify details from full text if heavily used. Data / sample: Since that day, around 330,000 Sina Weibo users have been engaged in discussions about the event; this engagement resulted in more than 4.7 million relevant posts.2 This context provides an appropriate research...

\paragraph{Findings} No clean finding sentence detected; verify from full text if central.

\paragraph{Contribution} Clarifies how state agencies use digital channels for communication and governance rather than treating online politics as purely oppositional.

\paragraph{Limitations} Design details need closer verification before the paper is used for strong causal claims.

\paragraph{Relevance to WeChat / My Project} Conceptually useful for a WeChat project because it speaks to the Chinese information environment, especially the cluster on participation, responsiveness, and political behavior, even if the platform is not WeChat.

\paragraph{Related Debate} Whether official digital communication mainly informs, disciplines, mobilizes, or improves responsiveness


\subsection{Jiang et al. 2021: Evaluating Chinese government WeChat official accounts...}
\textbf{Citation key}: \texttt{jiang\_2021\_government\_information\_quarterly}.\\
\textbf{Full citation}: Tingting Jiang, Ying Wang, Tianqianjin Lin, and Lina Shangguan 2021. "Evaluating Chinese government WeChat official accounts in public service delivery: A user-centered approach." Government Information Quarterly.

\paragraph{Summary} This study examines evaluating chinese government wechat official accounts in public service delivery: a user-centered approach using clickstream analysis evidence and reports that the former is similar to online media, sending group messages to followers on a regular basis. /static/Jszsyrq /accident/query...

\paragraph{Context} WeChat Official Accounts Platform (https://mp.weixin.qq.com/) is a popular type of social media in China which enables individuals, en terprises, and organizations to reach target audiences through their official accounts. Subscription accounts and...

\paragraph{Main Argument} The paper's core mechanism is that the former is similar to online media, sending group messages to followers on a regular basis. /static/Jszsyrq /accident/query /accident/accidentdetail /policyinfo/query /policyinfo/policylist /trafficbreak/infoquery...

\paragraph{Independent Variable(s)} Not stated in standard IV/DV terms; inferred from title/abstract if needed

\paragraph{Dependent Variable(s)} Not stated in standard IV/DV terms; verify if central to argument

\paragraph{Methods and Design} clickstream analysis. Unobtrusive clickstream analysis plus complementary qualitative methods. Data / sample: WeChat Official Accounts Platform (https://mp.weixin.qq.com/) is a popular type of social media in China which enables individuals, en terprises, and organizations to reach target audiences through their official...

\paragraph{Findings} The former is similar to online media, sending group messages to followers on a regular basis. /static/Jszsyrq /accident/query /accident/accidentdetail /policyinfo/query /policyinfo/policylist /trafficbreak/infoquery /trafficbreak/index /trafficbreak/detail /lookplate/platesearch /lookplate/norepeatinput /xianxing/weihao/type/7 lp14 lp15 lp16 lp17 lp18 lp19 rp0 rp1 rp2 rp3 rp4 rp5 rp6 rp7 rp8 rp9 /illegalcode/querycodeinput /illegalcode/querybycode /electriccars/showCardInfo /electriccars/showcardinfo?hplx= /tuoche/index /tuoche/detail /menu2018/report /fault/indexnew /fault/success www.whjg.gov.cn:8089/whtp/clx/wx/tpba/notice www.whjg.gov.cn:8089/whtp/clx/wx/tpba/show jw.whjg.gov.cn:8090/jw/wx/index.html?lwfl=wx jw.whjg.gov.cn:8090/jw/wx/report.html?lwfl=wx /backwaterreport/index /backwaterreport/success /freezereport/index User (U) Appoint-ment (A) Lookup (L) LT9 LT10 Reporting (R) LT11 Lookup of electric bicycle license plate information Lookup of vehicle detention records -- RT1 -- Reporting malfunctioning traffic facilities RT2 Reporting a fake-plated vehicle RT3 Reporting disciplinary violation RT4 Reporting road water RT5 Reporting road ice (continued on next page) 13 T.

\paragraph{Contribution} Directly extends the WeChat-centered literature by specifying platform-specific behavior rather than treating social media as generic.

\paragraph{Limitations} Behavioral trace data are strong on observed use but weaker on latent beliefs, interpretation, and broader generalizability.

\paragraph{Relevance to WeChat / My Project} Directly relevant: this paper studies WeChat behavior or institutions and can be used to theorize visibility, audience composition, and official-account interaction.

\paragraph{Related Debate} Whether semi-private WeChat interactions should be read as information use, endorsement, or strategically managed visibility


\subsection{Levy 2021: Social Media News Consumption and Polarization...}
\textbf{Citation key}: \texttt{levy\_2021\_aer}.\\
\textbf{Full citation}: Ro'ee Levy 2021. "Social Media, News Consumption, and Polarization: Evidence from a Field Experiment." American Economic Review.

\paragraph{Summary} This paper studies social media, news consumption, and polarization: evidence from a field experiment. American Economic Review 2021, 111(3): 831--870 https://doi.org/10.1257/aer.20191777 Social Media, News Consumption, and Polarization: Evidence from a Field Experiment By Ro'ee Levy* Does the consumption of...

\paragraph{Context} American Economic Review 2021, 111(3): 831--870 https://doi.org/10.1257/aer.20191777 Social Media, News Consumption, and Polarization: Evidence from a Field Experiment By Ro'ee Levy* Does the consumption of ideologically congruent news on social media...

\paragraph{Main Argument} Core mechanism maps onto the cluster on interpersonal discussion, networks, and opinion formation, but the exact statement should be verified from the full text if heavily used.

\paragraph{Independent Variable(s)} social media news exposure by conducting a large field experiment randomly offering participants subscriptions to...

\paragraph{Dependent Variable(s)} Facebook

\paragraph{Methods and Design} field experiment. field experiment. Data / sample: I collect data on the causal chain of media effects: subscriptions to outlets, exposure to news on Facebook, visits to online news sites, and sharing of posts, as well as changes in political opinions and attitudes....

\paragraph{Findings} No clean finding sentence detected; verify from full text if central.

\paragraph{Contribution} Offers comparatively strong leverage on causal mechanisms in digital political communication.

\paragraph{Limitations} Internal validity is strong, but external validity may depend on whether experimental treatments resemble real platform use.

\paragraph{Relevance to WeChat / My Project} Indirect benchmark: useful for comparison on platform effects or official communication, but needs careful translation to China's WeChat-centered environment.

\paragraph{Related Debate} How social ties and discussion networks mediate willingness to speak, share, and update opinions


\subsection{Lichter et al. 2021: The Long-Term Costs of Government Surveillance...}
\textbf{Citation key}: \texttt{lichter\_2021\_jeea}.\\
\textbf{Full citation}: Andreas Lichter, Max Loffler, and Sebastian Siegloch 2021. "The Long-Term Costs of Government Surveillance: Insights from Stasi Spying in East Germany." Journal of the European Economic Association.

\paragraph{Summary} This study examines the long-term costs of government surveillance: insights from stasi spying in east germany using descriptive or unclear from available parse evidence and reports that we find that a higher spying density led to persistently lower levels of interpersonal and institutional trust in...

\paragraph{Context} We also find substantial and long-lasting economic effects of Stasi surveillance, resulting in lower income, higher exposure to unemployment, and lower self-employment. (JEL: H11, N34, N44, P20) Teaching Slides A set of Teaching Slides to accompany this...

\paragraph{Main Argument} The paper's core mechanism is that we find that a higher spying density led to persistently lower levels of interpersonal and institutional trust in post-reunification Germany. We also find substantial and long-lasting economic effects of Stasi...

\paragraph{Independent Variable(s)} government surveillance

\paragraph{Dependent Variable(s)} civic capital and economic performance

\paragraph{Methods and Design} Descriptive or unclear from available parse. Observational or descriptive design; verify details from full text if heavily used. Data / sample: We investigate the long-run effects of government surveillance on civic capital and economic performance, studying the case of the Stasi in East Germany. We also find substantial and long-lasting economic effects of...

\paragraph{Findings} We find that a higher spying density led to persistently lower levels of interpersonal and institutional trust in post-reunification Germany. We also find substantial and long-lasting economic effects of Stasi surveillance, resulting in lower income, higher exposure to unemployment, and lower self-employment. (JEL: H11, N34, N44, P20) Teaching Slides A set of Teaching Slides to accompany this article are available online as Supplementary Data.

\paragraph{Contribution} Provides an external benchmark for mechanisms that may or may not travel to China and WeChat.

\paragraph{Limitations} Design details need closer verification before the paper is used for strong causal claims.

\paragraph{Relevance to WeChat / My Project} Peripheral for a WeChat project; best used as a comparative benchmark rather than core theoretical evidence.

\paragraph{Related Debate} Which actors can credibly correct misinformation and under what informational conditions


\subsection{Muchlinski et al. 2021: We need to go deeper measuring...}
\textbf{Citation key}: \texttt{muchlinski\_2021\_psrm}.\\
\textbf{Full citation}: David Muchlinski, Xiao Yang, Sarah Birch, Craig Macdonald, and Iadh Ounis 2021. "We need to go deeper: measuring electoral violence using convolutional neural networks and social media." Political Science Research and Methods.

\paragraph{Summary} This study examines we need to go deeper: measuring electoral violence using convolutional neural networks and social media using descriptive or unclear from available parse evidence and reports that we demonstrate that our methodology is more than 30 percent more accurate in measuring electoral violence than...

\paragraph{Context} While measuring the temporal aspect of this phenomenon is straightforward, measuring whether occurrences of violence are truly related to elections is more difficult.

\paragraph{Main Argument} The paper's core mechanism is that we demonstrate that our methodology is more than 30 percent more accurate in measuring electoral violence than previously utilized models. Additionally, we show that our measures of electoral violence conform to...

\paragraph{Independent Variable(s)} Not stated in standard IV/DV terms; inferred from title/abstract if needed

\paragraph{Dependent Variable(s)} Not stated in standard IV/DV terms; verify if central to argument

\paragraph{Methods and Design} Descriptive or unclear from available parse. Observational or descriptive design; verify details from full text if heavily used. Data / sample: Finally, we demonstrate the validity of our data by developing a qualitative coding ontology. https://doi.org/10.1017/psrm.2020.32 Published online by Cambridge University Press (2021), 9, 122--139...

\paragraph{Findings} We demonstrate that our methodology is more than 30 percent more accurate in measuring electoral violence than previously utilized models. Additionally, we show that our measures of electoral violence conform to theoretical expectations of this conflict more so than those that exist in event datasets commonly utilized to measure electoral violence including ACLED, ICEWS, and SCAD.

\paragraph{Contribution} Provides an external benchmark for mechanisms that may or may not travel to China and WeChat.

\paragraph{Limitations} Design details need closer verification before the paper is used for strong causal claims.

\paragraph{Relevance to WeChat / My Project} Indirect benchmark: useful for comparison on platform effects or official communication, but needs careful translation to China's WeChat-centered environment.

\paragraph{Related Debate} How social ties and discussion networks mediate willingness to speak, share, and update opinions


\subsection{Muller and Schwarz 2021: FANNING THE FLAMES OF HATE SOCIAL...}
\textbf{Citation key}: \texttt{muller\_2021\_jeea}.\\
\textbf{Full citation}: Karsten Muller and Carlo Schwarz 2021. "FANNING THE FLAMES OF HATE: SOCIAL MEDIA AND HATE CRIME." Journal of the European Economic Association.

\paragraph{Summary} This study examines fanning the flames of hate: social media and hate crime using descriptive or unclear from available parse evidence and reports that we show that antirefugee sentiment on Facebook predicts crimes against refugees in otherwise similar municipalities with higher social media usage. Consistent with...

\paragraph{Context} Our results suggest that social media can act as a propagation mechanism for violent crimes by enabling the spread of extreme viewpoints. (JEL: D74, J15, Z10, D72, O35) Teaching Slides A set of Teaching Slides to accompany this article are available online...

\paragraph{Main Argument} The paper's core mechanism is that we show that antirefugee sentiment on Facebook predicts crimes against refugees in otherwise similar municipalities with higher social media usage. Consistent with a role for "echo chambers," we find that rightwing social...

\paragraph{Independent Variable(s)} Not stated in standard IV/DV terms; inferred from title/abstract if needed

\paragraph{Dependent Variable(s)} Not stated in standard IV/DV terms; verify if central to argument

\paragraph{Methods and Design} Descriptive or unclear from available parse. Observational or descriptive design; verify details from full text if heavily used. Data / sample: Consistent with a role for "echo chambers," we find that rightwing social media posts contain narrower and more loaded content than news reports. Our results suggest that social media can act as a propagation...

\paragraph{Findings} We show that antirefugee sentiment on Facebook predicts crimes against refugees in otherwise similar municipalities with higher social media usage. Consistent with a role for "echo chambers," we find that rightwing social media posts contain narrower and more loaded content than news reports.

\paragraph{Contribution} Provides an external benchmark for mechanisms that may or may not travel to China and WeChat.

\paragraph{Limitations} Design details need closer verification before the paper is used for strong causal claims.

\paragraph{Relevance to WeChat / My Project} Indirect benchmark: useful for comparison on platform effects or official communication, but needs careful translation to China's WeChat-centered environment.

\paragraph{Related Debate} How findings from facebook and united states settings travel to a WeChat-centered China project


\subsection{Neagli 2021: Grassroots Astroturf or Something in Between...}
\textbf{Citation key}: \texttt{neagli\_2021\_journal\_of\_current\_chinese\_affairs}.\\
\textbf{Full citation}: Jackson Paul Neagli 2021. "Grassroots, Astroturf, or Something in Between? Semi-Official WeChat Accounts as Covert Vectors of Party-State Influence in Contemporary China." Journal of Current Chinese Affairs.

\paragraph{Summary} This paper studies grassroots, astroturf, or something in between? semi-official wechat accounts as covert vectors of party-state.... This article explores a facet of the Chinese propaganda apparatus that has yet to receive sufficient academic attention: the murky ecosystem of "semi-official" party- state...

\paragraph{Context} This article explores a facet of the Chinese propaganda apparatus that has yet to receive sufficient academic attention: the murky ecosystem of "semi-official" party- state presences on Chinese social media.

\paragraph{Main Argument} Core mechanism maps onto the cluster on platform affordances and public-private boundaries, but the exact statement should be verified from the full text if heavily used.

\paragraph{Independent Variable(s)} Author inference from title/abstract: Grassroots, Astroturf, or Something in Between? Semi-Official WeChat Accounts as Covert Vectors of Party-State...

\paragraph{Dependent Variable(s)} Not stated in standard IV/DV terms

\paragraph{Methods and Design} discourse analysis. Comparative discourse analysis. Data / sample: With a particular focus on WeChat public accounts, this investigation responds to two critical research questions: first, what differentiates official party-state social media presences from semi-official...

\paragraph{Findings} No clean finding sentence detected; verify from full text if central.

\paragraph{Contribution} Directly extends the WeChat-centered literature by specifying platform-specific behavior rather than treating social media as generic.

\paragraph{Limitations} Limited external validity and weaker causal leverage because the design is case-based and interpretive.

\paragraph{Relevance to WeChat / My Project} Directly relevant: this paper studies WeChat behavior or institutions and can be used to theorize visibility, audience composition, and official-account interaction.

\paragraph{Related Debate} Whether semi-private WeChat interactions should be read as information use, endorsement, or strategically managed visibility


\subsection{Sances 2021: Missing the Target Using Surveys to...}
\textbf{Citation key}: \texttt{sances\_2021\_psrm}.\\
\textbf{Full citation}: Michael W. Sances 2021. "Missing the Target? Using Surveys to Validate Social." Political Science Research and Methods.

\paragraph{Summary} This study examines missing the target? using surveys to validate social using survey evidence and reports that using a series of six surveys and 20 targeted ads, I show these ads regularly fail to reach their targets.

\paragraph{Context} The extent to which these ads actually reach their targets, however, is unknown.

\paragraph{Main Argument} The paper's core mechanism is that using a series of six surveys and 20 targeted ads, I show these ads regularly fail to reach their targets.

\paragraph{Independent Variable(s)} Author inference from title/abstract: Missing the Target? Using Surveys to Validate Social

\paragraph{Dependent Variable(s)} Not stated in standard IV/DV terms

\paragraph{Methods and Design} survey. Survey-based observational analysis. Data / sample: Facebook ads are increasingly used by political scientists as a method of survey recruitment. The success rate ranges from 23 to 99 percent, and ads targeted toward groups defined by self-reported data and broader...

\paragraph{Findings} Using a series of six surveys and 20 targeted ads, I show these ads regularly fail to reach their targets.

\paragraph{Contribution} Provides an external benchmark for mechanisms that may or may not travel to China and WeChat.

\paragraph{Limitations} Design details need closer verification before the paper is used for strong causal claims.

\paragraph{Relevance to WeChat / My Project} Peripheral for a WeChat project; best used as a comparative benchmark rather than core theoretical evidence.

\paragraph{Related Debate} How findings from facebook and germany settings travel to a WeChat-centered China project


\subsection{Wukich 2021: Government Social Media Engagement Strategies and...}
\textbf{Citation key}: \texttt{wukich\_2021\_public\_performance\_management\_review}.\\
\textbf{Full citation}: Clayton Wukich 2021. "Government Social Media Engagement Strategies and Public Roles." Public Performance \& Management Review.

\paragraph{Summary} This paper studies government social media engagement strategies and public roles. Through social media, public officials share information with the people they serve.

\paragraph{Context} In all, this article clarifies engagement opportunities based on public roles and more explicitly tethers social media content to foundational concepts in public administration.

\paragraph{Main Argument} Core mechanism maps onto the cluster on platform affordances and public-private boundaries, but the exact statement should be verified from the full text if heavily used.

\paragraph{Independent Variable(s)} Not stated in standard IV/DV terms; inferred from title/abstract if needed

\paragraph{Dependent Variable(s)} Not stated in standard IV/DV terms; verify if central to argument

\paragraph{Methods and Design} content analysis. Content analysis. Data / sample: Hurricane Florence provides a critical case to examine how cities in three U.S. states leveraged public roles for engagement purposes. Content analysis of Facebook data from 62 cities reveals their customer and...

\paragraph{Findings} No clean finding sentence detected; verify from full text if central.

\paragraph{Contribution} Clarifies how state agencies use digital channels for communication and governance rather than treating online politics as purely oppositional.

\paragraph{Limitations} Inference depends on coding choices and text classification; behavior and belief are only indirectly observed.

\paragraph{Relevance to WeChat / My Project} Indirect benchmark: useful for comparison on platform effects or official communication, but needs careful translation to China's WeChat-centered environment.

\paragraph{Related Debate} How interface design changes the visibility and meaning of digital expression


\subsection{Xu 2021: To Repress or to Co opt...}
\textbf{Citation key}: \texttt{xu\_2021\_ajps}.\\
\textbf{Full citation}: Xu Xu 2021. "To Repress or to Co-opt? Authoritarian Control in the Age of Digital Surveillance." American Journal of Political Science.

\paragraph{Summary} This study examines to repress or to co-opt? authoritarian control in the age of digital surveillance using difference-in-differences evidence and reports that using a difference-in-differences design that exploits temporal variation in digital surveillance systems among Chinese counties, I find that surveillance...

\paragraph{Context} This article studies the consequences of digital surveillance in dictatorships.

\paragraph{Main Argument} The paper's core mechanism is that using a difference-in-differences design that exploits temporal variation in digital surveillance systems among Chinese counties, I find that surveillance increases local governments' public security expenditure and...

\paragraph{Independent Variable(s)} Author inference from title/abstract: To Repress or to Co-opt? Authoritarian Control in the Age of Digital Surveillance

\paragraph{Dependent Variable(s)} Not stated in standard IV/DV terms

\paragraph{Methods and Design} difference-in-differences. Difference-in-differences. Data / sample: Congress passed the USA Freedom Act on June 2, 2015, which significantly curtailed the government's sweeping surveillance of phone calls and data gathering; on May 25, 2018, the European Union passed the strictest...

\paragraph{Findings} Using a difference-in-differences design that exploits temporal variation in digital surveillance systems among Chinese counties, I find that surveillance increases local governments' public security expenditure and arrests of political activists but decreases public goods provision. My theory and evidence suggest that improvements in governments' information make citizens worse off in dictatorships.

\paragraph{Contribution} Helps identify which mechanisms in the Chinese information environment are platform-specific and which generalize beyond a single app.

\paragraph{Limitations} Mechanisms may not transfer cleanly to WeChat's hybrid private-public architecture without further adaptation.

\paragraph{Relevance to WeChat / My Project} Conceptually useful for a WeChat project because it speaks to the Chinese information environment, especially the cluster on interpersonal discussion, networks, and opinion formation, even if the platform is not WeChat.

\paragraph{Related Debate} How social ties and discussion networks mediate willingness to speak, share, and update opinions


\subsection{Beknazar-Yuzbashev and Stalinski 2022: Do social media ads matter for...}
\textbf{Citation key}: \texttt{beknazar\_yuzbashev\_2022\_jpubeco}.\\
\textbf{Full citation}: George Beknazar-Yuzbashev and Mateusz Stalinski 2022. "Do social media ads matter for political behavior? A field experiment." Journal of Public Economics.

\paragraph{Summary} This paper studies do social media ads matter for political behavior? a field experiment. We exploit Facebook's

\paragraph{Context} We exploit Facebook's

\paragraph{Main Argument} Core mechanism maps onto the cluster on comparative digital media benchmarks, but the exact statement should be verified from the full text if heavily used.

\paragraph{Independent Variable(s)} Author inference from title/abstract: Do social media ads matter for political behavior? A field experiment

\paragraph{Dependent Variable(s)} Not stated in standard IV/DV terms

\paragraph{Methods and Design} Descriptive or unclear from available parse. Observational or descriptive design; verify details from full text if heavily used. Data / sample: We exploit Facebook's of measuring exposure to political ads.23 We expected that the ad filter would not be easily accessible or well-known among Facebook users. A.3 for ad examples).26 We targeted users aged 18+,...

\paragraph{Findings} No clean finding sentence detected; verify from full text if central.

\paragraph{Contribution} Provides an external benchmark for mechanisms that may or may not travel to China and WeChat.

\paragraph{Limitations} Design details need closer verification before the paper is used for strong causal claims.

\paragraph{Relevance to WeChat / My Project} Indirect benchmark: useful for comparison on platform effects or official communication, but needs careful translation to China's WeChat-centered environment.

\paragraph{Related Debate} How findings from facebook and united states settings travel to a WeChat-centered China project


\subsection{Braghieri et al. 2022: Social Media and Mental Health}
\textbf{Citation key}: \texttt{braghieri\_2022\_aer}.\\
\textbf{Full citation}: Luca Braghieri, Ro'ee Levy, and Alexey Makarin 2022. "Social Media and Mental Health." American Economic Review.

\paragraph{Summary} This paper studies social media and mental health. of Facebook across US colleges.

\paragraph{Context} of Facebook across US colleges. Our analysis couples

\paragraph{Main Argument} Core mechanism maps onto the cluster on comparative digital media benchmarks, but the exact statement should be verified from the full text if heavily used.

\paragraph{Independent Variable(s)} Not stated in standard IV/DV terms; inferred from title/abstract if needed

\paragraph{Dependent Variable(s)} Not stated in standard IV/DV terms; verify if central to argument

\paragraph{Methods and Design} Descriptive or unclear from available parse. Observational or descriptive design; verify details from full text if heavily used. Data / sample: Unclear from available parse; verify from full text if used centrally

\paragraph{Findings} No clean finding sentence detected; verify from full text if central.

\paragraph{Contribution} Provides an external benchmark for mechanisms that may or may not travel to China and WeChat.

\paragraph{Limitations} Design details need closer verification before the paper is used for strong causal claims.

\paragraph{Relevance to WeChat / My Project} Indirect benchmark: useful for comparison on platform effects or official communication, but needs careful translation to China's WeChat-centered environment.

\paragraph{Related Debate} How findings from facebook and united states settings travel to a WeChat-centered China project


\subsection{Criado and Villodre 2022: Revisiting social media institutionalization in government...}
\textbf{Citation key}: \texttt{criado\_2022\_government\_information\_quarterly}.\\
\textbf{Full citation}: J. Ignacio Criado and Julian Villodre 2022. "Revisiting social media institutionalization in government. An empirical analysis of barriers." Government Information Quarterly.

\paragraph{Summary} This paper studies revisiting social media institutionalization in government. an empirical analysis of barriers. Institutionalization of technology denotes a process of formalization in organizations, also in the public sector.

\paragraph{Context} Institutionalization of technology denotes a process of formalization in organizations, also in the public sector. Despite institutionalization is being carried out with similar patterns to previous technologies (Red dick \& Norris, 2013), social media...

\paragraph{Main Argument} Core mechanism maps onto the cluster on comparative digital media benchmarks, but the exact statement should be verified from the full text if heavily used.

\paragraph{Independent Variable(s)} Not stated in standard IV/DV terms; inferred from title/abstract if needed

\paragraph{Dependent Variable(s)} Not stated in standard IV/DV terms; verify if central to argument

\paragraph{Methods and Design} case study. Case-study design. Data / sample: Social media are developed and maintained by third-party companies, and individuals can easily use these platforms. strategy combining different sources of data for triangulation purposes, including a survey on...

\paragraph{Findings} No clean finding sentence detected; verify from full text if central.

\paragraph{Contribution} Provides an external benchmark for mechanisms that may or may not travel to China and WeChat.

\paragraph{Limitations} Limited external validity and weaker causal leverage because the design is case-based and interpretive.

\paragraph{Relevance to WeChat / My Project} Indirect benchmark: useful for comparison on platform effects or official communication, but needs careful translation to China's WeChat-centered environment.

\paragraph{Related Debate} How findings from social media (general) and comparative / general settings travel to a WeChat-centered China project


\subsection{Danieli et al. 2022: Decomposing the Rise of the Populist...}
\textbf{Citation key}: \texttt{danieli\_2022\_wp}.\\
\textbf{Full citation}: Oren Danieli, Noam Gidron, Shinnosuke Kikuchi, and Roee Levy 2022. "Decomposing the Rise of the Populist Radical Right." SSRN Electronic Journal.

\paragraph{Summary} This paper studies decomposing the rise of the populist radical right. =4255937 The rise of populist radical right parties (PRRP) is one of the most significant political developments in recent decades.

\paragraph{Context} While a growing literature has identified various factors affecting PRRP support (Rodrik, 2018; Fetzer, 2019; Mudde, 2019; Manacorda et al., 2023), there is still no consensus on what is the primary factor for the extensive rise of these parties and which...

\paragraph{Main Argument} Core mechanism maps onto the cluster on interpersonal discussion, networks, and opinion formation, but the exact statement should be verified from the full text if heavily used.

\paragraph{Independent Variable(s)} Not stated in standard IV/DV terms; inferred from title/abstract if needed

\paragraph{Dependent Variable(s)} Not stated in standard IV/DV terms; verify if central to argument

\paragraph{Methods and Design} Descriptive or unclear from available parse. Observational or descriptive design; verify details from full text if heavily used. Data / sample: In certain countries, populist governments have resulted in lower GDP and significant erosion of democratic norms and institutions (McCoy and Somer, 2019; Funke et al., 2022).

\paragraph{Findings} No clean finding sentence detected; verify from full text if central.

\paragraph{Contribution} Provides an external benchmark for mechanisms that may or may not travel to China and WeChat.

\paragraph{Limitations} Mechanisms may not transfer cleanly to WeChat's hybrid private-public architecture without further adaptation.

\paragraph{Relevance to WeChat / My Project} Peripheral for a WeChat project; best used as a comparative benchmark rather than core theoretical evidence.

\paragraph{Related Debate} How social ties and discussion networks mediate willingness to speak, share, and update opinions


\subsection{Golovchenko 2022: Fighting Propaganda with Censorship A Study...}
\textbf{Citation key}: \texttt{golovchenko\_2022\_jop}.\\
\textbf{Full citation}: Yevgeniy Golovchenko 2022. "Fighting Propaganda with Censorship: A Study of the Ukrainian Ban on Russian Social Media." The Journal of Politics.

\paragraph{Summary} This paper studies fighting propaganda with censorship: a study of the ukrainian ban on russian social media. Fighting Propaganda with Censorship: A Study of the Ukrainian Ban on Russian Social Media Yevgeniy Golovchenko, University of Copenhagen Many states have become concerned with Russian cyberattacks and...

\paragraph{Context} Fighting Propaganda with Censorship: A Study of the Ukrainian Ban on Russian Social Media Yevgeniy Golovchenko, University of Copenhagen Many states have become concerned with Russian cyberattacks and online propaganda. The Ukrainian government responded...

\paragraph{Main Argument} Core mechanism maps onto the cluster on censorship, propaganda, and state influence, but the exact statement should be verified from the full text if heavily used.

\paragraph{Independent Variable(s)} Not stated in standard IV/DV terms; inferred from title/abstract if needed

\paragraph{Dependent Variable(s)} Not stated in standard IV/DV terms; verify if central to argument

\paragraph{Methods and Design} Descriptive or unclear from available parse. Observational or descriptive design; verify details from full text if heavily used. Data / sample: By exploiting a natural experiment in Ukraine, I find that the sudden censorship policy reduced activity on VKontakte, despite the fact that a vast majority of the users were legally and technically able to bypass the...

\paragraph{Findings} No clean finding sentence detected; verify from full text if central.

\paragraph{Contribution} Provides an external benchmark for mechanisms that may or may not travel to China and WeChat.

\paragraph{Limitations} Design details need closer verification before the paper is used for strong causal claims.

\paragraph{Relevance to WeChat / My Project} Indirect benchmark: useful for comparison on platform effects or official communication, but needs careful translation to China's WeChat-centered environment.

\paragraph{Related Debate} Whether digital control works through persuasion, selective amplification, covert influence, or audience management


\subsection{HAGER and KRAKOWSKI 2022: Does State Repression Spark Protests Evidence...}
\textbf{Citation key}: \texttt{hager\_2022\_apsr}.\\
\textbf{Full citation}: ANSELM HAGER and KRZYSZTOF KRAKOWSKI 2022. "Does State Repression Spark Protests? Evidence from Secret Police Surveillance in Communist Poland." American Political Science Review.

\paragraph{Summary} This paper studies does state repression spark protests? evidence from secret police surveillance in communist poland. https://doi.org/10.1017/S0003055421000770 Published online by Cambridge University Press A uthoritarian regimes around the globe use repression to secure their survival.

\paragraph{Context} Much has been written about violent forms of repression.1 We know comparatively little, however, about nonviolent forms of repression.

\paragraph{Main Argument} Core mechanism maps onto the cluster on comparative digital media benchmarks, but the exact statement should be verified from the full text if heavily used.

\paragraph{Independent Variable(s)} physical surveillance

\paragraph{Dependent Variable(s)} antiregime resistance (Davenport 2005)

\paragraph{Methods and Design} Descriptive or unclear from available parse. Observational or descriptive design; verify details from full text if heavily used. Data / sample: While there is growing interest in digital monitoring (Gohdes 2020; King, Pan, and Roberts 2013; Xu 2020), we lack systematic accounts that explore whether physical surveillance affects antiregime resistance...

\paragraph{Findings} No clean finding sentence detected; verify from full text if central.

\paragraph{Contribution} Provides an external benchmark for mechanisms that may or may not travel to China and WeChat.

\paragraph{Limitations} Design details need closer verification before the paper is used for strong causal claims.

\paragraph{Relevance to WeChat / My Project} Peripheral for a WeChat project; best used as a comparative benchmark rather than core theoretical evidence.

\paragraph{Related Debate} How findings from internet / online media and poland settings travel to a WeChat-centered China project


\subsection{Hansen 2022: The Perils of Estimating Disengagement Effects...}
\textbf{Citation key}: \texttt{hansen\_2022\_bjps}.\\
\textbf{Full citation}: Tanja Marie Hansen 2022. "The Perils of Estimating Disengagement Effects of Deadly Terrorist Attacks Utilizing Social Media Data." British Journal of Political Science.

\paragraph{Summary} This paper studies the perils of estimating disengagement effects of deadly terrorist attacks utilizing social media data. This comment discusses the impact of social media rule enforcement protocols on research on online data sources.

\paragraph{Context} This comment discusses the impact of social media rule enforcement protocols on research on online data sources. It argues that the conclusions of the article 'Do Islamic State's Deadly Attacks Disengage, Deter, or Mobilize Supporters?' concerning the...

\paragraph{Main Argument} Core mechanism maps onto the cluster on comparative digital media benchmarks, but the exact statement should be verified from the full text if heavily used.

\paragraph{Independent Variable(s)} Deadly Terrorist Attacks Utilizing Social Media Data. This comment discusses the impact of social media rule...

\paragraph{Dependent Variable(s)} research on online data sources

\paragraph{Methods and Design} Descriptive or unclear from available parse. Observational or descriptive design; verify details from full text if heavily used. Data / sample: This comment discusses the impact of social media rule enforcement protocols on research on online data sources.

\paragraph{Findings} No clean finding sentence detected; verify from full text if central.

\paragraph{Contribution} Provides an external benchmark for mechanisms that may or may not travel to China and WeChat.

\paragraph{Limitations} Design details need closer verification before the paper is used for strong causal claims.

\paragraph{Relevance to WeChat / My Project} Indirect benchmark: useful for comparison on platform effects or official communication, but needs careful translation to China's WeChat-centered environment.

\paragraph{Related Debate} How findings from twitter / x and united states settings travel to a WeChat-centered China project


\subsection{Li et al. 2022: Unpacking government social media messaging strategies...}
\textbf{Citation key}: \texttt{li\_2022\_policy\_internet}.\\
\textbf{Full citation}: Yiran Li, Yanto Chandra, and Yingying Fan 2022. "Unpacking government social media messaging strategies during the COVID-19 pandemic in China." Policy \& Internet.

\paragraph{Summary} This study examines unpacking government social media messaging strategies during the covid-19 pandemic in china using topic modeling evidence and reports that the results show that local government agencies predominantly used the first two strategies, whereas the central government mainly relied on the last two....

\paragraph{Context} However, doing so during an emergency--- particularly a pandemic---is often complex and challenging.

\paragraph{Main Argument} The paper's core mechanism is that the results show that local government agencies predominantly used the first two strategies, whereas the central government mainly relied on the last two. These strategies explain how various levels of government engaged...

\paragraph{Independent Variable(s)} Not stated in standard IV/DV terms; inferred from title/abstract if needed

\paragraph{Dependent Variable(s)} Not stated in standard IV/DV terms; verify if central to argument

\paragraph{Methods and Design} topic modeling. Topic modeling on platform texts. Data / sample: Analyzing government social media posts during the COVID-19 outbreak in Wuhan ("text as data"), we conduct topic modeling analysis and identify four strategies that characterize Chinese governments' responses to a...

\paragraph{Findings} The results show that local government agencies predominantly used the first two strategies, whereas the central government mainly relied on the last two. These strategies explain how various levels of government engaged in agile governance through their communication with citizens, highlight the coordination and control work undertaken by governments at all levels, and demonstrate how these methods shielded the central government from blame for the pandemic. .

\paragraph{Contribution} Clarifies how state agencies use digital channels for communication and governance rather than treating online politics as purely oppositional.

\paragraph{Limitations} Inference depends on coding choices and text classification; behavior and belief are only indirectly observed.

\paragraph{Relevance to WeChat / My Project} Conceptually useful for a WeChat project because it speaks to the Chinese information environment, especially the cluster on participation, responsiveness, and political behavior, even if the platform is not WeChat.

\paragraph{Related Debate} Whether official digital communication mainly informs, disciplines, mobilizes, or improves responsiveness


\subsection{Munger et al. 2022: Political Knowledge and Misinformation in the...}
\textbf{Citation key}: \texttt{munger\_2022\_bjps}.\\
\textbf{Full citation}: Kevin Munger, Patrick J. Egan, Jonathan Nagler, Jonathan Ronen, and Joshua Tucker 2022. "Political Knowledge and Misinformation in the Era of Social Media: Evidence From the 2015 UK Election." British Journal of Political Science.

\paragraph{Summary} This study examines political knowledge and misinformation in the era of social media: evidence from the 2015 uk election using survey evidence and reports that but in a troubling demonstration of campaigns' ability to manipulate knowledge, messages from the parties also shifted voters' assessments of the economy...

\paragraph{Context} This study measures changes in political knowledge among a panel of voters surveyed during the 2015 UK general election campaign while monitoring the political information to which they were exposed on the Twitter social media platform.

\paragraph{Main Argument} The paper's core mechanism is that but in a troubling demonstration of campaigns' ability to manipulate knowledge, messages from the parties also shifted voters' assessments of the economy and immigration in directions favorable to the parties' platforms,...

\paragraph{Independent Variable(s)} information exposure

\paragraph{Dependent Variable(s)} changes in political knowledge

\paragraph{Methods and Design} survey. Survey-based observational analysis. Data / sample: Twitter use led to higher levels of knowledge about politics and public affairs, as information from news media improved knowledge of politically relevant facts, and messages sent by political parties increased...

\paragraph{Findings} But in a troubling demonstration of campaigns' ability to manipulate knowledge, messages from the parties also shifted voters' assessments of the economy and immigration in directions favorable to the parties' platforms, leaving some voters with beliefs further from the truth at the end of the campaign than they were at its beginning. suggest that as social media plays an ever more prominent role in political life, its effects on political knowledge will in many ways reinforce those of traditional media.

\paragraph{Contribution} Provides an external benchmark for mechanisms that may or may not travel to China and WeChat.

\paragraph{Limitations} Design details need closer verification before the paper is used for strong causal claims.

\paragraph{Relevance to WeChat / My Project} Indirect benchmark: useful for comparison on platform effects or official communication, but needs careful translation to China's WeChat-centered environment.

\paragraph{Related Debate} Which actors can credibly correct misinformation and under what informational conditions


\subsection{Pan et al. 2022: How government-controlled media shifts policy attitudes...}
\textbf{Citation key}: \texttt{pan\_2022\_psrm}.\\
\textbf{Full citation}: Jennifer Pan, Zijie Shao, and Yiqing Xu 2022. "How government-controlled media shifts policy attitudes through framing." Political Science Research and Methods.

\paragraph{Summary} This study examines how government-controlled media shifts policy attitudes through framing using descriptive or unclear from available parse evidence and reports that by conducting an experiment that exposes respondents to government-controlled media---in the form of TV news segments---on issues where the regime...

\paragraph{Context} Research shows that government-controlled media is an effective tool for authoritarian regimes to shape public opinion. Does government-controlled media remain effective when it is required to support changes in positions that autocrats take on issues?

\paragraph{Main Argument} The paper's core mechanism is that by conducting an experiment that exposes respondents to government-controlled media---in the form of TV news segments---on issues where the regime substantially changed its policy positions, we find that by framing the same...

\paragraph{Independent Variable(s)} government-controlled media

\paragraph{Dependent Variable(s)} policy attitudes through framing

\paragraph{Methods and Design} Descriptive or unclear from available parse. Observational or descriptive design; verify details from full text if heavily used. Data / sample: By conducting an experiment that exposes respondents to government-controlled media---in the form of TV news segments---on issues where the regime substantially changed its policy positions, we find that by framing the...

\paragraph{Findings} By conducting an experiment that exposes respondents to government-controlled media---in the form of TV news segments---on issues where the regime substantially changed its policy positions, we find that by framing the same issue differently, government-controlled media moves respondents to adopt policy positions closer to the ones espoused by the regime regardless of individual predisposition.

\paragraph{Contribution} Helps identify which mechanisms in the Chinese information environment are platform-specific and which generalize beyond a single app.

\paragraph{Limitations} Design details need closer verification before the paper is used for strong causal claims.

\paragraph{Relevance to WeChat / My Project} Conceptually useful for a WeChat project because it speaks to the Chinese information environment, especially the cluster on information exposure, selective exposure, and persuasion, even if the platform is not WeChat.

\paragraph{Related Debate} Whether exposure changes attitudes directly or is filtered by selection, framing, and prior preferences


\subsection{Sun and Zhao 2022: Delegated Censorship The Dynamic Layered and...}
\textbf{Citation key}: \texttt{sun\_2022\_politics\_society}.\\
\textbf{Full citation}: Taiyi Sun and Quansheng Zhao 2022. "Delegated Censorship: The Dynamic, Layered, and Multistage Information Control Regime in China." Politics \& Society.

\paragraph{Summary} This study examines delegated censorship: the dynamic, layered, and multistage information control regime in china using interviews / focus groups evidence and reports that such practices indicate differentiated utilization of social and market forces by the state.

\paragraph{Context} This approach aims to maintain regime stability and legitimacy while minimizing costs.

\paragraph{Main Argument} The paper's core mechanism is that such practices indicate differentiated utilization of social and market forces by the state.

\paragraph{Independent Variable(s)} Not stated in standard IV/DV terms; inferred from title/abstract if needed

\paragraph{Dependent Variable(s)} Not stated in standard IV/DV terms; verify if central to argument

\paragraph{Methods and Design} interviews / focus groups. Observational or descriptive design; verify details from full text if heavily used. Data / sample: An online political publication, Global China (), was created by the authors, and the pattern and record of articles being censored was analyzed. Using results from A/B tests on the articles and interviews with...

\paragraph{Findings} Such practices indicate differentiated utilization of social and market forces by the state.

\paragraph{Contribution} Helps identify which mechanisms in the Chinese information environment are platform-specific and which generalize beyond a single app.

\paragraph{Limitations} Design details need closer verification before the paper is used for strong causal claims.

\paragraph{Relevance to WeChat / My Project} Conceptually useful for a WeChat project because it speaks to the Chinese information environment, especially the cluster on censorship, propaganda, and state influence, even if the platform is not WeChat.

\paragraph{Related Debate} Whether digital control works through persuasion, selective amplification, covert influence, or audience management


\subsection{Wukich 2022: Social media engagement forms in government...}
\textbf{Citation key}: \texttt{wukich\_2022\_government\_information\_quarterly}.\\
\textbf{Full citation}: Clayton Wukich 2022. "Social media engagement forms in government: A structure-content framework." Government Information Quarterly.

\paragraph{Summary} This paper studies social media engagement forms in government: a structure-content framework. Over the past 15 years, the public has largely embraced social networking applications such as Facebook and Twitter, and government agencies followed suit "to be where the people are" (Mergel, 2013, p.

\paragraph{Context} Proponents lauded the technologies' potential to enable mean ingful interaction (Linders, 2012; Noveck, 2009); however, the type of deliberation they envisioned has largely failed to materialize (Bryer, 2020; Mergel, 2012, 2013; Reddick \& Norris, 2013;...

\paragraph{Main Argument} Core mechanism maps onto the cluster on participation, responsiveness, and political behavior, but the exact statement should be verified from the full text if heavily used.

\paragraph{Independent Variable(s)} Not stated in standard IV/DV terms; inferred from title/abstract if needed

\paragraph{Dependent Variable(s)} Not stated in standard IV/DV terms; verify if central to argument

\paragraph{Methods and Design} case study. Case-study design. Data / sample: Case selection Social media content is highly context specific (Mergel, 2016). A time-bounded case study accounts for the specific context in which engagement takes place rather than the researcher relying on random...

\paragraph{Findings} No clean finding sentence detected; verify from full text if central.

\paragraph{Contribution} Provides an external benchmark for mechanisms that may or may not travel to China and WeChat.

\paragraph{Limitations} Limited external validity and weaker causal leverage because the design is case-based and interpretive.

\paragraph{Relevance to WeChat / My Project} Indirect benchmark: useful for comparison on platform effects or official communication, but needs careful translation to China's WeChat-centered environment.

\paragraph{Related Debate} Whether official digital communication mainly informs, disciplines, mobilizes, or improves responsiveness


\subsection{Xu et al. 2022: Information Control and Public Support for...}
\textbf{Citation key}: \texttt{xu\_2022\_jop}.\\
\textbf{Full citation}: Xu Xu, Genia Kostka, and Xun Cao 2022. "Information Control and Public Support for Social Credit Systems in China." The Journal of Politics.

\paragraph{Summary} This study examines information control and public support for social credit systems in china using descriptive or unclear from available parse evidence and reports that using a field survey experiment, we show that respondents are not more supportive of the SCS when receiving information about its order-maintenance...

\paragraph{Context} Information Control and Public Support for Social Credit Systems in China Xu Xu, Princeton University Genia Kostka, Freie Universitat Berlin Xun Cao, Pennsylvania State University Critics see China's social credit system (SCS) as a tool of surveillance and...

\paragraph{Main Argument} The paper's core mechanism is that using a field survey experiment, we show that respondents are not more supportive of the SCS when receiving information about its order-maintenance role but largely decrease their support when knowing its repressive...

\paragraph{Independent Variable(s)} Not stated in standard IV/DV terms; inferred from title/abstract if needed

\paragraph{Dependent Variable(s)} Not stated in standard IV/DV terms; verify if central to argument

\paragraph{Methods and Design} Descriptive or unclear from available parse. Observational or descriptive design; verify details from full text if heavily used. Data / sample: Unclear from available parse; verify from full text if used centrally

\paragraph{Findings} Using a field survey experiment, we show that respondents are not more supportive of the SCS when receiving information about its order-maintenance role but largely decrease their support when knowing its repressive function. We further conduct text analysis of state media reports 2244 / Social Credit Systems in China Xu Xu, Genia Kostka, and Xun Cao and show that the government portrays the SCS in a very positive way, with little mention of its repressive function.

\paragraph{Contribution} Helps identify which mechanisms in the Chinese information environment are platform-specific and which generalize beyond a single app.

\paragraph{Limitations} Mechanisms may not transfer cleanly to WeChat's hybrid private-public architecture without further adaptation.

\paragraph{Relevance to WeChat / My Project} Conceptually useful for a WeChat project because it speaks to the Chinese information environment, especially the cluster on censorship, propaganda, and state influence, even if the platform is not WeChat.

\paragraph{Related Debate} Whether digital control works through persuasion, selective amplification, covert influence, or audience management


\subsection{Ali and Qazi 2023: Countering misinformation on social media through...}
\textbf{Citation key}: \texttt{ali\_2023\_jde}.\\
\textbf{Full citation}: Ayesha Ali and Ihsan Ayyub Qazi 2023. "Countering misinformation on social media through educational interventions: Evidence from a randomized experiment in Pakistan." Journal of Development Economics.

\paragraph{Summary} This study examines countering misinformation on social media through educational interventions: evidence from a randomized experiment... using randomized experiment evidence and reports that we find no evidence that video-based general educational messages about misinformation have any statistically significant...

\paragraph{Context} However, when such messages are augmented with personalized feedback based on individuals' past engagement with fake news, we find an improvement of 3.3 percentage points or 4.5 percent in correctly identifying news (true or fake) relative to the control...

\paragraph{Main Argument} The paper's core mechanism is that we find no evidence that video-based general educational messages about misinformation have any statistically significant impact on the ability to correctly identify news items as true or fake. However, when such messages...

\paragraph{Independent Variable(s)} Not stated in standard IV/DV terms; inferred from title/abstract if needed

\paragraph{Dependent Variable(s)} Not stated in standard IV/DV terms; verify if central to argument

\paragraph{Methods and Design} randomized experiment. randomized experiment. Data / sample: Fake news is a growing problem in developing countries with potentially far-reaching consequences. We find no evidence that video-based general educational messages about misinformation have any statistically...

\paragraph{Findings} We find no evidence that video-based general educational messages about misinformation have any statistically significant impact on the ability to correctly identify news items as true or fake. However, when such messages are augmented with personalized feedback based on individuals' past engagement with fake news, we find an improvement of 3.3 percentage points or 4.5 percent in correctly identifying news (true or fake) relative to the control group.

\paragraph{Contribution} Provides an external benchmark for mechanisms that may or may not travel to China and WeChat.

\paragraph{Limitations} Design details need closer verification before the paper is used for strong causal claims.

\paragraph{Relevance to WeChat / My Project} Indirect benchmark: useful for comparison on platform effects or official communication, but needs careful translation to China's WeChat-centered environment.

\paragraph{Related Debate} Which actors can credibly correct misinformation and under what informational conditions


\subsection{Allie 2023: Facial Recognition Technology and Voter Turnout}
\textbf{Citation key}: \texttt{allie\_2023\_jop}.\\
\textbf{Full citation}: Feyaad Allie 2023. "Facial Recognition Technology and Voter Turnout." The Journal of Politics.

\paragraph{Summary} This paper studies facial recognition technology and voter turnout. id p3041197.

\paragraph{Context} id p3041197. Desai, Zuheir, and Alexander Lee.

\paragraph{Main Argument} Core mechanism maps onto the cluster on participation, responsiveness, and political behavior, but the exact statement should be verified from the full text if heavily used.

\paragraph{Independent Variable(s)} Electronic Voting in India." Political Science Research and Methods 9:398--413. Number 1 January 2023 / 333 Driscoll,...

\paragraph{Dependent Variable(s)} Bureaucrat Diversity and Election Bias in India

\paragraph{Methods and Design} Descriptive or unclear from available parse. Observational or descriptive design; verify details from full text if heavily used. Data / sample: Unclear from available parse; verify from full text if used centrally

\paragraph{Findings} No clean finding sentence detected; verify from full text if central.

\paragraph{Contribution} Provides an external benchmark for mechanisms that may or may not travel to China and WeChat.

\paragraph{Limitations} Mechanisms may not transfer cleanly to WeChat's hybrid private-public architecture without further adaptation.

\paragraph{Relevance to WeChat / My Project} Peripheral for a WeChat project; best used as a comparative benchmark rather than core theoretical evidence.

\paragraph{Related Debate} Whether official digital communication mainly informs, disciplines, mobilizes, or improves responsiveness


\subsection{Beam 2023: Social media as a recruitment and...}
\textbf{Citation key}: \texttt{beam\_2023\_jde}.\\
\textbf{Full citation}: Emily A. Beam 2023. "Social media as a recruitment and data collection tool: Experimental evidence on the relative effectiveness of web surveys and chatbots." Journal of Development Economics.

\paragraph{Summary} This study examines social media as a recruitment and data collection tool: experimental evidence on the relative effectiveness of web... using randomized experiment evidence and reports that the results suggest that chatbot responses match CATI responses on multiple dimensions of quality. I show that social...

\paragraph{Context} Online technologies enable lower-cost, rapid data collection, but concerns about sample composition biases and mode-specific measurement error impede their use. I conduct a randomized experiment in the Philippines to test the effectiveness of web-form and...

\paragraph{Main Argument} The paper's core mechanism is that the results suggest that chatbot responses match CATI responses on multiple dimensions of quality. I show that social media-based recruitment can be an attractive alternative for targeted sampling and that online surveys...

\paragraph{Independent Variable(s)} Not stated in standard IV/DV terms; inferred from title/abstract if needed

\paragraph{Dependent Variable(s)} Not stated in standard IV/DV terms; verify if central to argument

\paragraph{Methods and Design} randomized experiment. randomized experiment. Data / sample: Online technologies enable lower-cost, rapid data collection, but concerns about sample composition biases and mode-specific measurement error impede their use. Chatbot surveys yield higher response rates and...

\paragraph{Findings} The results suggest that chatbot responses match CATI responses on multiple dimensions of quality. I show that social media-based recruitment can be an attractive alternative for targeted sampling and that online surveys can be implemented effectively at a fraction of the cost of phone surveys.

\paragraph{Contribution} Provides an external benchmark for mechanisms that may or may not travel to China and WeChat.

\paragraph{Limitations} Design details need closer verification before the paper is used for strong causal claims.

\paragraph{Relevance to WeChat / My Project} Indirect benchmark: useful for comparison on platform effects or official communication, but needs careful translation to China's WeChat-centered environment.

\paragraph{Related Debate} How findings from social media (general) and united states settings travel to a WeChat-centered China project


\subsection{Cirone and Hobbs 2023: Asymmetric flooding as a tool for...}
\textbf{Citation key}: \texttt{cirone\_2023\_psrm}.\\
\textbf{Full citation}: Alexandra Cirone and William Hobbs 2023. "Asymmetric flooding as a tool for foreign influence on social media." Political Science Research and Methods.

\paragraph{Summary} This study examines asymmetric flooding as a tool for foreign influence on social media using content analysis evidence and reports that here, we document the use of apolitical content---content that could counteract mobilization efforts and escape detection in future campaigns. Using automated text analysis and hand...

\paragraph{Context} Research on Russian troll activity during the 2016 US presidential campaign largely focused on divisive partisan messaging. Here, we document the use of apolitical content---content that could counteract mobilization efforts and escape detection in future...

\paragraph{Main Argument} The paper's core mechanism is that here, we document the use of apolitical content---content that could counteract mobilization efforts and escape detection in future campaigns. Using automated text analysis and hand coding to construct a timeline of IRA...

\paragraph{Independent Variable(s)} Not stated in standard IV/DV terms; inferred from title/abstract if needed

\paragraph{Dependent Variable(s)} Not stated in standard IV/DV terms; verify if central to argument

\paragraph{Methods and Design} content analysis. Content analysis. Data / sample: Unclear from available parse; verify from full text if used centrally

\paragraph{Findings} Here, we document the use of apolitical content---content that could counteract mobilization efforts and escape detection in future campaigns. Using automated text analysis and hand coding to construct a timeline of IRA messaging on Twitter, we find left-leaning trolls posted large volumes of entertainment content in their artificial liberal community and shifted away from political content late in the campaign.

\paragraph{Contribution} Provides an external benchmark for mechanisms that may or may not travel to China and WeChat.

\paragraph{Limitations} Inference depends on coding choices and text classification; behavior and belief are only indirectly observed.

\paragraph{Relevance to WeChat / My Project} Indirect benchmark: useful for comparison on platform effects or official communication, but needs careful translation to China's WeChat-centered environment.

\paragraph{Related Debate} How findings from twitter / x and russia settings travel to a WeChat-centered China project


\subsection{Fu 2023: Propagandization of Relative Gratification How Chinese...}
\textbf{Citation key}: \texttt{fu\_2023\_polcomm}.\\
\textbf{Full citation}: King-Wa Fu 2023. "Propagandization of Relative Gratification: How Chinese State Media Portray the International Pandemic." Political Communication.

\paragraph{Summary} This paper studies propagandization of relative gratification: how chinese state media portray the international pandemic. .

\paragraph{Context} . The SAGE Handbook of propaganda.

\paragraph{Main Argument} Core mechanism maps onto the cluster on censorship, propaganda, and state influence, but the exact statement should be verified from the full text if heavily used.

\paragraph{Independent Variable(s)} Not stated in standard IV/DV terms; inferred from title/abstract if needed

\paragraph{Dependent Variable(s)} Not stated in standard IV/DV terms; verify if central to argument

\paragraph{Methods and Design} Descriptive or unclear from available parse. Observational or descriptive design; verify details from full text if heavily used. Data / sample: Unclear from available parse; verify from full text if used centrally

\paragraph{Findings} No clean finding sentence detected; verify from full text if central.

\paragraph{Contribution} Helps identify which mechanisms in the Chinese information environment are platform-specific and which generalize beyond a single app.

\paragraph{Limitations} Design details need closer verification before the paper is used for strong causal claims.

\paragraph{Relevance to WeChat / My Project} Conceptually useful for a WeChat project because it speaks to the Chinese information environment, especially the cluster on censorship, propaganda, and state influence, even if the platform is not WeChat.

\paragraph{Related Debate} Whether digital control works through persuasion, selective amplification, covert influence, or audience management


\subsection{GRAHAM 2023: Measuring Misperceptions}
\textbf{Citation key}: \texttt{graham\_2023\_apsr}.\\
\textbf{Full citation}: MATTHEW H. GRAHAM 2023. "Measuring Misperceptions?." American Political Science Review.

\paragraph{Summary} This paper studies measuring misperceptions?. American Political Science Review (2023) 117, 1, 80--102 doi:10.1017/S0003055422000387 Science Association.  The Author(s), 2022.

\paragraph{Context} American Political Science Review (2023) 117, 1, 80--102 doi:10.1017/S0003055422000387 Science Association.  The Author(s), 2022. Published by Cambridge University Press on behalf of the American Political Measuring Misperceptions?

\paragraph{Main Argument} Core mechanism maps onto the cluster on comparative digital media benchmarks, but the exact statement should be verified from the full text if heavily used.

\paragraph{Independent Variable(s)} Author inference from title/abstract: Measuring Misperceptions

\paragraph{Dependent Variable(s)} Not stated in standard IV/DV terms

\paragraph{Methods and Design} survey. Survey-based observational analysis. Data / sample: By assuming the burden of proof for its interpretation of survey responses, research can develop a more trustworthy basis for understanding the prevalence and consequences of political misperceptions.

\paragraph{Findings} No clean finding sentence detected; verify from full text if central.

\paragraph{Contribution} Provides an external benchmark for mechanisms that may or may not travel to China and WeChat.

\paragraph{Limitations} Mechanisms may not transfer cleanly to WeChat's hybrid private-public architecture without further adaptation.

\paragraph{Relevance to WeChat / My Project} Peripheral for a WeChat project; best used as a comparative benchmark rather than core theoretical evidence.

\paragraph{Related Debate} How findings from not platform-specific and united states settings travel to a WeChat-centered China project


\subsection{Jiang and Tang 2023: Effects of local government social media...}
\textbf{Citation key}: \texttt{jiang\_2023\_jpart}.\\
\textbf{Full citation}: Hanchen Jiang and Xiao Tang 2023. "Effects of local government social media use on citizen compliance during a crisis: Evidence from the COVID-19 crisis in China." Public Administration.

\paragraph{Summary} This paper studies effects of local government social media use on citizen compliance during a crisis: evidence from the covid-19.... Beijing, China Improving citizen compliance is a major goal of public admin- 2 istration, especially during crises.

\paragraph{Context} Beijing, China Improving citizen compliance is a major goal of public admin- 2 istration, especially during crises. Although social media are School of Public Policy and Management, Tsinghua University, Beijing, China 3 widely used by government agencies...

\paragraph{Main Argument} Core mechanism maps onto the cluster on participation, responsiveness, and political behavior, but the exact statement should be verified from the full text if heavily used.

\paragraph{Independent Variable(s)} local government social media use

\paragraph{Dependent Variable(s)} citizen compliance during a crisis: Evidence from the COVID-19 crisis in China

\paragraph{Methods and Design} panel / observational analysis. Observational or descriptive design; verify details from full text if heavily used. Data / sample: Based on an original daily panel dataset of 189 cit- Xiao Tang, Institute for Contemporary China Studies, Tsinghua University, Beijing 100084, ies in China during COVID-19, this study provides empirical China....

\paragraph{Findings} No clean finding sentence detected; verify from full text if central.

\paragraph{Contribution} Clarifies how state agencies use digital channels for communication and governance rather than treating online politics as purely oppositional.

\paragraph{Limitations} Design details need closer verification before the paper is used for strong causal claims.

\paragraph{Relevance to WeChat / My Project} Conceptually useful for a WeChat project because it speaks to the Chinese information environment, especially the cluster on participation, responsiveness, and political behavior, even if the platform is not WeChat.

\paragraph{Related Debate} Whether official digital communication mainly informs, disciplines, mobilizes, or improves responsiveness


\subsection{Kocak and Kbrs 2023: Bennett and Segerberg 2012 Shirky 2009...}
\textbf{Citation key}: \texttt{kocak\_2023\_bjps}.\\
\textbf{Full citation}: Korhan Kocak and Ozgur Kbrs 2023. "(Bennett and Segerberg 2012; Shirky 2009). In some contexts, this assertion prevailed: greater." British Journal of Political Science.

\paragraph{Summary} This paper studies (bennett and segerberg 2012; shirky 2009). in some contexts, this assertion prevailed: greater. As internet penetration rapidly expanded throughout the world, press freedom and government accountability improved in some countries but backslid in others.

\paragraph{Context} As internet penetration rapidly expanded throughout the world, press freedom and government accountability improved in some countries but backslid in others. We propose a formal model that provides a mechanism that explains the observed divergent paths of...

\paragraph{Main Argument} We propose a formal model that provides a mechanism that explains the observed divergent paths of countries.

\paragraph{Independent Variable(s)} Not stated in standard IV/DV terms; inferred from title/abstract if needed

\paragraph{Dependent Variable(s)} Not stated in standard IV/DV terms; verify if central to argument

\paragraph{Methods and Design} Descriptive or unclear from available parse. Observational or descriptive design; verify details from full text if heavily used. Data / sample: As internet penetration rapidly expanded throughout the world, press freedom and government accountability improved in some countries but backslid in others. We propose a formal model that provides a mechanism that...

\paragraph{Findings} No clean finding sentence detected; verify from full text if central.

\paragraph{Contribution} Provides an external benchmark for mechanisms that may or may not travel to China and WeChat.

\paragraph{Limitations} Design details need closer verification before the paper is used for strong causal claims.

\paragraph{Relevance to WeChat / My Project} Peripheral for a WeChat project; best used as a comparative benchmark rather than core theoretical evidence.

\paragraph{Related Debate} How findings from social media (general) and united kingdom settings travel to a WeChat-centered China project


\subsection{Litvin 2023: This Hearing Should Be Flipped Democratic...}
\textbf{Citation key}: \texttt{litvin\_2023\_apsr}.\\
\textbf{Full citation}: Boris Litvin 2023. ""This Hearing Should Be Flipped": Democratic Spectatorship, Social Media, and the Problem of Demagogic Candor." American Political Science Review.

\paragraph{Summary} This paper studies "this hearing should be flipped": democratic spectatorship, social media, and the problem of demagogic candor. American Political Science Review (2023) 117, 1, 153--167 doi:10.1017/S0003055422000508 Science Association.  The Author(s), 2022.

\paragraph{Context} American Political Science Review (2023) 117, 1, 153--167 doi:10.1017/S0003055422000508 Science Association.  The Author(s), 2022. Published by Cambridge University Press on behalf of the American Political "This Hearing Should Be Flipped": Democratic...

\paragraph{Main Argument} Core mechanism maps onto the cluster on comparative digital media benchmarks, but the exact statement should be verified from the full text if heavily used.

\paragraph{Independent Variable(s)} Not stated in standard IV/DV terms; inferred from title/abstract if needed

\paragraph{Dependent Variable(s)} Not stated in standard IV/DV terms; verify if central to argument

\paragraph{Methods and Design} Descriptive or unclear from available parse. Observational or descriptive design; verify details from full text if heavily used. Data / sample: Unclear from available parse; verify from full text if used centrally

\paragraph{Findings} No clean finding sentence detected; verify from full text if central.

\paragraph{Contribution} Provides an external benchmark for mechanisms that may or may not travel to China and WeChat.

\paragraph{Limitations} Design details need closer verification before the paper is used for strong causal claims.

\paragraph{Relevance to WeChat / My Project} Indirect benchmark: useful for comparison on platform effects or official communication, but needs careful translation to China's WeChat-centered environment.

\paragraph{Related Debate} How findings from social media (general) and united states settings travel to a WeChat-centered China project


\subsection{Lutscher 2023: Hot topics Denial-of-Service attacks on news...}
\textbf{Citation key}: \texttt{lutscher\_2023\_psrm}.\\
\textbf{Full citation}: Philipp M. Lutscher 2023. "Hot topics: Denial-of-Service attacks on news websites in autocracies." Political Science Research and Methods.

\paragraph{Summary} This study examines hot topics: denial-of-service attacks on news websites in autocracies using content analysis evidence and reports that statistical analyses show that news content correlates to DoS attacks. While the Internet has made it easier for the press to circumvent censorship, my study could show that...

\paragraph{Context} Most authoritarian countries censor the press. As a response, many opposition and independent news outlets have found refuge on the Internet.

\paragraph{Main Argument} The paper's core mechanism is that statistical analyses show that news content correlates to DoS attacks. While the Internet has made it easier for the press to circumvent censorship, my study could show that news outlets are still vulnerable to censorship...

\paragraph{Independent Variable(s)} Not stated in standard IV/DV terms; inferred from title/abstract if needed

\paragraph{Dependent Variable(s)} Not stated in standard IV/DV terms; verify if central to argument

\paragraph{Methods and Design} content analysis. Content analysis. Data / sample: Most authoritarian countries censor the press. Email: philipp.lutscher@stv.uio.no (Received 11 November 2020; revised 8 June 2021; accepted 29 August 2021; first published online 6 December 2021) Abstract Most...

\paragraph{Findings} Statistical analyses show that news content correlates to DoS attacks. While the Internet has made it easier for the press to circumvent censorship, my study could show that news outlets are still vulnerable to censorship in the digital age.

\paragraph{Contribution} Provides an external benchmark for mechanisms that may or may not travel to China and WeChat.

\paragraph{Limitations} Inference depends on coding choices and text classification; behavior and belief are only indirectly observed.

\paragraph{Relevance to WeChat / My Project} Peripheral for a WeChat project; best used as a comparative benchmark rather than core theoretical evidence.

\paragraph{Related Debate} Whether digital control works through persuasion, selective amplification, covert influence, or audience management


\subsection{Ye and Zhao 2023: I know it's sensitive Internet censorship...}
\textbf{Citation key}: \texttt{ye\_2023\_discourse\_context\_media}.\\
\textbf{Full citation}: WeiMing Ye and Luming Zhao 2023. ""I know it's sensitive": Internet censorship, recoding, and the sensitive word culture in China." Discourse, Context \& Media.

\paragraph{Summary} This paper studies "i know it's sensitive": internet censorship, recoding, and the sensitive word culture in china. China has witnessed the synergistic development of civil society and the Internet (Zhu and Cheng, 2011; Fitzgerald et al., 2022).

\paragraph{Context} However, Internet censorship by governments around the world to monitor illegal or undesirable content to varying degrees is widespread (Burnett and Feamster, 2013; Zittrain et al., 2017; Chang and Lin, 2020).

\paragraph{Main Argument} Core mechanism maps onto the cluster on censorship, propaganda, and state influence, but the exact statement should be verified from the full text if heavily used.

\paragraph{Independent Variable(s)} Not stated in standard IV/DV terms; inferred from title/abstract if needed

\paragraph{Dependent Variable(s)} Not stated in standard IV/DV terms; verify if central to argument

\paragraph{Methods and Design} Descriptive or unclear from available parse. Observational or descriptive design; verify details from full text if heavily used. Data / sample: The spread of information and communication technologies (ICTs) has bridged the digital divide (Sparks, 2013) and facilitated access to quality and affordable information services for 1 billion Chinese Internet...

\paragraph{Findings} No clean finding sentence detected; verify from full text if central.

\paragraph{Contribution} Helps identify which mechanisms in the Chinese information environment are platform-specific and which generalize beyond a single app.

\paragraph{Limitations} Design details need closer verification before the paper is used for strong causal claims.

\paragraph{Relevance to WeChat / My Project} Conceptually useful for a WeChat project because it speaks to the Chinese information environment, especially the cluster on censorship, propaganda, and state influence, even if the platform is not WeChat.

\paragraph{Related Debate} Whether digital control works through persuasion, selective amplification, covert influence, or audience management


\subsection{Yuan et al. 2023: Government Digital Transformation Understanding the Role...}
\textbf{Citation key}: \texttt{yuan\_2023\_government\_information\_quarterly}.\\
\textbf{Full citation}: Yun-Peng Yuan, Yogesh K. Dwivedi, Garry Wei-Han Tan, Tat-Huei Cham, Keng-Boon Ooi, Eugene Cheng-Xi Aw, and Wendy Currie 2023. "Government Digital Transformation: Understanding the Role of Government Social Media." Government Information Quarterly.

\paragraph{Summary} This paper studies government digital transformation: understanding the role of government social media. The imminent growth of digital innovations and heightened citizens' expectations have compelled governments to engage in digital trans formation.

\paragraph{Context} The imminent growth of digital innovations and heightened citizens' expectations have compelled governments to engage in digital trans formation. The undertaking of digital initiatives by governments is ex pected to increase government operation...

\paragraph{Main Argument} Core mechanism maps onto the cluster on participation, responsiveness, and political behavior, but the exact statement should be verified from the full text if heavily used.

\paragraph{Independent Variable(s)} Not stated in standard IV/DV terms; inferred from title/abstract if needed

\paragraph{Dependent Variable(s)} Not stated in standard IV/DV terms; verify if central to argument

\paragraph{Methods and Design} Descriptive or unclear from available parse. Observational or descriptive design; verify details from full text if heavily used. Data / sample: The fifth section presents the data analysis, followed by the

\paragraph{Findings} No clean finding sentence detected; verify from full text if central.

\paragraph{Contribution} Clarifies how state agencies use digital channels for communication and governance rather than treating online politics as purely oppositional.

\paragraph{Limitations} Design details need closer verification before the paper is used for strong causal claims.

\paragraph{Relevance to WeChat / My Project} Indirect benchmark: useful for comparison on platform effects or official communication, but needs careful translation to China's WeChat-centered environment.

\paragraph{Related Debate} Whether official digital communication mainly informs, disciplines, mobilizes, or improves responsiveness


\subsection{Binstok et al. 2024: Policies or Prejudices An Analysis of...}
\textbf{Citation key}: \texttt{binstok\_2024\_jeea}.\\
\textbf{Full citation}: Noam Binstok, Eric D Gould, and Todd Kaplan 2024. "Policies or Prejudices? An Analysis of Antisemitic and Anti-Israel Views on Social Media and Social Surveys." Journal of the European Economic Association.

\paragraph{Summary} This study examines policies or prejudices? an analysis of antisemitic and anti-israel views on social media and social surveys using survey evidence and reports that a causal interpretation of this finding is supported by an IV analysis motivated by Voigtlander and Voth who show that Nazi indoctrination during the...

\paragraph{Context} This paper examines the extent to which personal biases affect political views, in the context of how antisemitism influences opinions about Israel. Two empirical analyses are conducted.

\paragraph{Main Argument} The paper's core mechanism is that a causal interpretation of this finding is supported by an IV analysis motivated by Voigtlander and Voth who show that Nazi indoctrination during the WWII period had a lifelong impact on antisemitic views.

\paragraph{Independent Variable(s)} Author inference from title/abstract: Policies or Prejudices? An Analysis of Antisemitic and Anti-Israel Views on Social Media and Social Surveys

\paragraph{Dependent Variable(s)} Not stated in standard IV/DV terms

\paragraph{Methods and Design} survey. Survey-based observational analysis. Data / sample: The second empirical analysis uses the 2016 wave of the German Social Survey, which reveals a strong and robust relationship between several commonly used measures of antisemitic beliefs and holding anti-Israel...

\paragraph{Findings} A causal interpretation of this finding is supported by an IV analysis motivated by Voigtlander and Voth who show that Nazi indoctrination during the WWII period had a lifelong impact on antisemitic views.

\paragraph{Contribution} Provides an external benchmark for mechanisms that may or may not travel to China and WeChat.

\paragraph{Limitations} Design details need closer verification before the paper is used for strong causal claims.

\paragraph{Relevance to WeChat / My Project} Indirect benchmark: useful for comparison on platform effects or official communication, but needs careful translation to China's WeChat-centered environment.

\paragraph{Related Debate} How social ties and discussion networks mediate willingness to speak, share, and update opinions


\subsection{Boussalis et al. 2024: Rally Round the Mask Congressional Social...}
\textbf{Citation key}: \texttt{boussalis\_2024\_jop}.\\
\textbf{Full citation}: Constantine Boussalis, Travis G. Coan, and Mirya R. Holman 2024. "Rally 'Round the Mask: Congressional Social Media Images and Masking during COVID-19." The Journal of Politics.

\paragraph{Summary} This study examines rally 'round the mask: congressional social media images and masking during covid-19 using descriptive or unclear from available parse evidence and reports that we show that the appearance of masks varies across MOC and time, pointing to the importance of evaluating variation in images to...

\paragraph{Context} SHORT ARTICLE Rally 'Round the Mask: Congressional Social Media Images and Masking during COVID-19 Constantine Boussalis, Trinity College Dublin Travis G. Coan, University of Exeter and Exeter Q-Step Centre Mirya R.

\paragraph{Main Argument} The paper's core mechanism is that we show that the appearance of masks varies across MOC and time, pointing to the importance of evaluating variation in images to understand political communication.

\paragraph{Independent Variable(s)} Not stated in standard IV/DV terms; inferred from title/abstract if needed

\paragraph{Dependent Variable(s)} Not stated in standard IV/DV terms; verify if central to argument

\paragraph{Methods and Design} Descriptive or unclear from available parse. Observational or descriptive design; verify details from full text if heavily used. Data / sample: Messages from elites can shape public views, increase trust in the government, and offer opportunities to transform distress into national pride (Hetherington and Nelson 2003).

\paragraph{Findings} We show that the appearance of masks varies across MOC and time, pointing to the importance of evaluating variation in images to understand political communication.

\paragraph{Contribution} Provides an external benchmark for mechanisms that may or may not travel to China and WeChat.

\paragraph{Limitations} Design details need closer verification before the paper is used for strong causal claims.

\paragraph{Relevance to WeChat / My Project} Indirect benchmark: useful for comparison on platform effects or official communication, but needs careful translation to China's WeChat-centered environment.

\paragraph{Related Debate} How findings from facebook and united states settings travel to a WeChat-centered China project


\subsection{Chen and Li 2024: Accountability from cyberspace Scandal exposure on...}
\textbf{Citation key}: \texttt{chen\_2024\_psrm}.\\
\textbf{Full citation}: Shuo Chen and Yiran Li 2024. "Accountability from cyberspace? Scandal exposure on the Internet and official governance in China." Political Science Research and Methods.

\paragraph{Summary} This study examines accountability from cyberspace? scandal exposure on the internet and official governance in china using descriptive or unclear from available parse evidence and reports that we find that the government employs clear strategies: higher levels of online discussion lead to quicker government...

\paragraph{Context} This article explores the effects of social media on government accountability under authoritarian regimes.

\paragraph{Main Argument} The paper's core mechanism is that we find that the government employs clear strategies: higher levels of online discussion lead to quicker government responses and more severe punishment of the officials involved. Our findings highlight the accountability...

\paragraph{Independent Variable(s)} social media

\paragraph{Dependent Variable(s)} government accountability under authoritarian regimes

\paragraph{Methods and Design} Descriptive or unclear from available parse. Observational or descriptive design; verify details from full text if heavily used. Data / sample: We use a unique dataset containing records of scandals discussed on microblogs in China to systematically study their effects on the government response process and officials' disciplining. We use a unique dataset...

\paragraph{Findings} We find that the government employs clear strategies: higher levels of online discussion lead to quicker government responses and more severe punishment of the officials involved. Our findings highlight the accountability mechanism facilitated by social media and the power of social media empowerment. https://doi.org/10.1017/psrm.2023.13 Published online by Cambridge University Press Series 2000-08.

\paragraph{Contribution} Helps identify which mechanisms in the Chinese information environment are platform-specific and which generalize beyond a single app.

\paragraph{Limitations} Design details need closer verification before the paper is used for strong causal claims.

\paragraph{Relevance to WeChat / My Project} Conceptually useful for a WeChat project because it speaks to the Chinese information environment, especially the cluster on interpersonal discussion, networks, and opinion formation, even if the platform is not WeChat.

\paragraph{Related Debate} How social ties and discussion networks mediate willingness to speak, share, and update opinions


\subsection{Cobb-Clark et al. 2024: Surveillance and Self-Control}
\textbf{Citation key}: \texttt{cobb\_clark\_2024\_ej}.\\
\textbf{Full citation}: Deborah A Cobb-Clark, Sarah C Dahmann, Daniel A Kamhofer, and Hannah Schildberg-Horisch 2024. "Surveillance and Self-Control." The Economic Journal.

\paragraph{Summary} This paper studies surveillance and self-control. The Economic Journal, 134 (May), 1666--1682 https://doi.org/10.1093/ej/uead111  The Author(s) 2023.

\paragraph{Context} The Economic Journal, 134 (May), 1666--1682 https://doi.org/10.1093/ej/uead111  The Author(s) 2023. Published by Oxford University Press on behalf of Royal Economic Society.

\paragraph{Main Argument} Core mechanism maps onto the cluster on comparative digital media benchmarks, but the exact statement should be verified from the full text if heavily used.

\paragraph{Independent Variable(s)} Not stated in standard IV/DV terms; inferred from title/abstract if needed

\paragraph{Dependent Variable(s)} Not stated in standard IV/DV terms; verify if central to argument

\paragraph{Methods and Design} Descriptive or unclear from available parse. Observational or descriptive design; verify details from full text if heavily used. Data / sample: Unclear from available parse; verify from full text if used centrally

\paragraph{Findings} No clean finding sentence detected; verify from full text if central.

\paragraph{Contribution} Provides an external benchmark for mechanisms that may or may not travel to China and WeChat.

\paragraph{Limitations} Mechanisms may not transfer cleanly to WeChat's hybrid private-public architecture without further adaptation.

\paragraph{Relevance to WeChat / My Project} Peripheral for a WeChat project; best used as a comparative benchmark rather than core theoretical evidence.

\paragraph{Related Debate} How findings from not platform-specific and germany settings travel to a WeChat-centered China project


\subsection{Enriquez et al. 2024: Mass Political Information on Social Media...}
\textbf{Citation key}: \texttt{enriquez\_2024\_jeea}.\\
\textbf{Full citation}: Jose Ramon Enriquez, Horacio Larreguy, John Marshall, and Alberto Simpser 2024. "Mass Political Information on Social Media: Facebook Ads, Electorate Saturation, and Electoral Accountability in Mexico." Journal of the European Economic Association.

\paragraph{Summary} This study examines mass political information on social media: facebook ads, electorate saturation, and electoral accountability in mexico using descriptive or unclear from available parse evidence and reports that these findings demonstrate how mass media can ignite social interactions to promote political...

\paragraph{Context} Social media's capacity to quickly and inexpensively reach large audiences almost simultaneously has the potential to promote electoral accountability. Beyond increasing direct exposure to information, high saturation campaigns---which target substantial...

\paragraph{Main Argument} The paper's core mechanism is that these findings demonstrate how mass media can ignite social interactions to promote political accountability. (JEL: D72, O12) The editor in charge of this paper was Imran Rasul.

\paragraph{Independent Variable(s)} Not stated in standard IV/DV terms; inferred from title/abstract if needed

\paragraph{Dependent Variable(s)} Not stated in standard IV/DV terms; verify if central to argument

\paragraph{Methods and Design} Descriptive or unclear from available parse. Observational or descriptive design; verify details from full text if heavily used. Data / sample: We standardize the SCI, which is available for all but one municipality in our sample, to facilitate interpretation.

\paragraph{Findings} These findings demonstrate how mass media can ignite social interactions to promote political accountability. (JEL: D72, O12) The editor in charge of this paper was Imran Rasul.

\paragraph{Contribution} Provides an external benchmark for mechanisms that may or may not travel to China and WeChat.

\paragraph{Limitations} Design details need closer verification before the paper is used for strong causal claims.

\paragraph{Relevance to WeChat / My Project} Indirect benchmark: useful for comparison on platform effects or official communication, but needs careful translation to China's WeChat-centered environment.

\paragraph{Related Debate} How interface design changes the visibility and meaning of digital expression


\subsection{Forestal 2024: european douchebag 2012b European douchebag is...}
\textbf{Citation key}: \texttt{forestal\_2024\_apsr}.\\
\textbf{Full citation}: Jennifer Forestal 2024. "(european\_douchebag 2012b). European\_douchebag is not unique in his experience." American Political Science Review.

\paragraph{Summary} This paper studies (european\_douchebag 2012b). european\_douchebag is not unique in his experience. , Reddit also hosted a productive shaming of european\_douchebag---at least initially.

\paragraph{Context} , Reddit also hosted a productive shaming of european\_douchebag---at least initially. After posting a photo of a Sikh woman to r/funny, european\_douchebag was shamed by his fellow Redditors.

\paragraph{Main Argument} Core mechanism maps onto the cluster on interpersonal discussion, networks, and opinion formation, but the exact statement should be verified from the full text if heavily used.

\paragraph{Independent Variable(s)} Not stated in standard IV/DV terms; inferred from title/abstract if needed

\paragraph{Dependent Variable(s)} Not stated in standard IV/DV terms; verify if central to argument

\paragraph{Methods and Design} Descriptive or unclear from available parse. Observational or descriptive design; verify details from full text if heavily used. Data / sample: Unclear from available parse; verify from full text if used centrally

\paragraph{Findings} No clean finding sentence detected; verify from full text if central.

\paragraph{Contribution} Provides an external benchmark for mechanisms that may or may not travel to China and WeChat.

\paragraph{Limitations} Design details need closer verification before the paper is used for strong causal claims.

\paragraph{Relevance to WeChat / My Project} Peripheral for a WeChat project; best used as a comparative benchmark rather than core theoretical evidence.

\paragraph{Related Debate} How social ties and discussion networks mediate willingness to speak, share, and update opinions


\subsection{Fujiwara et al. 2024: The Effect of Social Media on...}
\textbf{Citation key}: \texttt{fujiwara\_2024\_jeea}.\\
\textbf{Full citation}: Thomas Fujiwara, Karsten Muller, and Carlo Schwarz 2024. "The Effect of Social Media on Elections: Evidence from The United States." Journal of the European Economic Association.

\paragraph{Summary} This study examines the effect of social media on elections: evidence from the united states using survey evidence and reports that we show that this variation is unrelated to observable county characteristics and electoral outcomes before the launch of Twitter. Our results indicate that Twitter lowered the...

\paragraph{Context} Evidence from survey data, primary elections, and text analysis of millions of tweets suggests that Twitter's relatively liberal content may have persuaded voters with moderate views to vote against Donald Trump. (JEL: L82, D72) Teaching Slides A set of...

\paragraph{Main Argument} The paper's core mechanism is that we show that this variation is unrelated to observable county characteristics and electoral outcomes before the launch of Twitter. Our results indicate that Twitter lowered the Republican vote share in the 2016 and 2020...

\paragraph{Independent Variable(s)} Social Media

\paragraph{Dependent Variable(s)} Elections: Evidence from The United States

\paragraph{Methods and Design} survey. Survey-based observational analysis. Data / sample: We use variation in the number of Twitter users across counties induced by early adopters at the 2007 South by Southwest festival, a key event in Twitter's rise to popularity. Evidence from survey data, primary...

\paragraph{Findings} We show that this variation is unrelated to observable county characteristics and electoral outcomes before the launch of Twitter. Our results indicate that Twitter lowered the Republican vote share in the 2016 and 2020 presidential elections, but had limited effects on Congressional elections and previous presidential elections.

\paragraph{Contribution} Provides an external benchmark for mechanisms that may or may not travel to China and WeChat.

\paragraph{Limitations} Design details need closer verification before the paper is used for strong causal claims.

\paragraph{Relevance to WeChat / My Project} Indirect benchmark: useful for comparison on platform effects or official communication, but needs careful translation to China's WeChat-centered environment.

\paragraph{Related Debate} How interface design changes the visibility and meaning of digital expression


\subsection{Gaughan 2024: Estimating Ideal Points of British MPs...}
\textbf{Citation key}: \texttt{gaughan\_2024\_bjps}.\\
\textbf{Full citation}: Conor Gaughan 2024. "Estimating Ideal Points of British MPs Through Their Social Media Followership." British Journal of Political Science.

\paragraph{Summary} This paper studies estimating ideal points of british mps through their social media followership. Ideal points of MPs in the UK House of Commons (HoC) are characteristically difficult to ascertain due to tight party discipline and strategic voting by opposition members.

\paragraph{Context} Ideal points of MPs in the UK House of Commons (HoC) are characteristically difficult to ascertain due to tight party discipline and strategic voting by opposition members. This research note generates left/ right ideal point estimates for 591 British MPs...

\paragraph{Main Argument} Core mechanism maps onto the cluster on interpersonal discussion, networks, and opinion formation, but the exact statement should be verified from the full text if heavily used.

\paragraph{Independent Variable(s)} Not stated in standard IV/DV terms; inferred from title/abstract if needed

\paragraph{Dependent Variable(s)} Not stated in standard IV/DV terms; verify if central to argument

\paragraph{Methods and Design} survey. Survey-based observational analysis. Data / sample: Specifically, estimates are derived by conducting correspondence analysis (CA) on MP Twitter (X) follower networks, which are subsequently validated against an expert survey, confirming that these estimates have a...

\paragraph{Findings} No clean finding sentence detected; verify from full text if central.

\paragraph{Contribution} Provides an external benchmark for mechanisms that may or may not travel to China and WeChat.

\paragraph{Limitations} Design details need closer verification before the paper is used for strong causal claims.

\paragraph{Relevance to WeChat / My Project} Indirect benchmark: useful for comparison on platform effects or official communication, but needs careful translation to China's WeChat-centered environment.

\paragraph{Related Debate} How social ties and discussion networks mediate willingness to speak, share, and update opinions


\subsection{Pradel et al. 2024: W hen is speech on social...}
\textbf{Citation key}: \texttt{pradel\_2024\_apsr}.\\
\textbf{Full citation}: Franziska Pradel, Jan Zilinsky, Spyros Kosmidis, and Yannis Theocharis 2024. "W hen is speech on social media toxic enough to warrant content moderation? Platforms impose." American Political Science Review.

\paragraph{Summary} This paper studies w hen is speech on social media toxic enough to warrant content moderation? platforms impose. P https://doi.org/10.1017/S000305542300134X Published online by Cambridge University Press hilosophers have argued that civility as a normative ideal is a requirement for democratic discourse.

\paragraph{Context} P https://doi.org/10.1017/S000305542300134X Published online by Cambridge University Press hilosophers have argued that civility as a normative ideal is a requirement for democratic discourse. Apart from a type of behavior that displays good manners,...

\paragraph{Main Argument} Core mechanism maps onto the cluster on comparative digital media benchmarks, but the exact statement should be verified from the full text if heavily used.

\paragraph{Independent Variable(s)} Author inference from title/abstract: W hen is speech on social media toxic enough to warrant content moderation? Platforms impose

\paragraph{Dependent Variable(s)} Not stated in standard IV/DV terms

\paragraph{Methods and Design} Descriptive or unclear from available parse. Observational or descriptive design; verify details from full text if heavily used. Data / sample: Apart from a type of behavior that displays good manners, character, courtesy, and selfcontrol, civility in normative visions of theorists like Habermas and Rawls is deemed a virtue because it promotes respect and...

\paragraph{Findings} No clean finding sentence detected; verify from full text if central.

\paragraph{Contribution} Provides an external benchmark for mechanisms that may or may not travel to China and WeChat.

\paragraph{Limitations} Design details need closer verification before the paper is used for strong causal claims.

\paragraph{Relevance to WeChat / My Project} Indirect benchmark: useful for comparison on platform effects or official communication, but needs careful translation to China's WeChat-centered environment.

\paragraph{Related Debate} How findings from social media (general) and united states settings travel to a WeChat-centered China project


\subsection{Qin et al. 2024: Social Media and Collective Action in...}
\textbf{Citation key}: \texttt{qin\_2024\_ecma}.\\
\textbf{Full citation}: Bei Qin, David Stromberg, and Yanhui Wu 2024. "Social Media and Collective Action in China." Econometrica.

\paragraph{Summary} This study examines social media and collective action in china using descriptive or unclear from available parse evidence and reports that rECENT STUDIES show that communication through social media is powerful in organizing grassroots protests (Tucker, Nagler, Mac Duffee Metzger, Barbera, Penfold-Brown, and...

\paragraph{Context} However, it remains an open question whether this extends to a more strictly controlled media environment.

\paragraph{Main Argument} The paper's core mechanism is that rECENT STUDIES show that communication through social media is powerful in organizing grassroots protests (Tucker, Nagler, Mac Duffee Metzger, Barbera, Penfold-Brown, and Bonneau (2016), Steinert-Threlkeld (2017),...

\paragraph{Independent Variable(s)} Not stated in standard IV/DV terms; inferred from title/abstract if needed

\paragraph{Dependent Variable(s)} Not stated in standard IV/DV terms; verify if central to argument

\paragraph{Methods and Design} Descriptive or unclear from available parse. Observational or descriptive design; verify details from full text if heavily used. Data / sample: Answering this question is important, as an increasing number of governments are tightening their control over social media to suppress protests (Howard (2011), Woolley and Howard (2019)). to Chinese social media...

\paragraph{Findings} RECENT STUDIES show that communication through social media is powerful in organizing grassroots protests (Tucker, Nagler, Mac Duffee Metzger, Barbera, Penfold-Brown, and Bonneau (2016), Steinert-Threlkeld (2017), Acemoglu, Hassan, and Tahoun (2018); Enikolopov, Makarin, and Petrova (2010); Fergusson and Molina (2020)).

\paragraph{Contribution} Helps identify which mechanisms in the Chinese information environment are platform-specific and which generalize beyond a single app.

\paragraph{Limitations} Design details need closer verification before the paper is used for strong causal claims.

\paragraph{Relevance to WeChat / My Project} Conceptually useful for a WeChat project because it speaks to the Chinese information environment, especially the cluster on censorship, propaganda, and state influence, even if the platform is not WeChat.

\paragraph{Related Debate} Whether digital control works through persuasion, selective amplification, covert influence, or audience management


\subsection{Thomson 2024: The bureaucratic politics of authoritarian repression...}
\textbf{Citation key}: \texttt{thomson\_2024\_psrm}.\\
\textbf{Full citation}: Henry Thomson 2024. "The bureaucratic politics of authoritarian repression: intra-agency reform and surveillance capacity in communist Poland." Political Science Research and Methods.

\paragraph{Summary} This study examines the bureaucratic politics of authoritarian repression: intra-agency reform and surveillance capacity in communist poland using difference-in-differences evidence and reports that difference-in-differences models find that when security headquarters were duplicated through an administrative...

\paragraph{Context} Bureaucratic reorganizations within security institutions cause significant variation in their behavior, however.

\paragraph{Main Argument} The paper's core mechanism is that difference-in-differences models find that when security headquarters were duplicated through an administrative reform, the proliferation of higherlevel posts within the service caused a large and statistically...

\paragraph{Independent Variable(s)} intra-agency reforms

\paragraph{Dependent Variable(s)} surveillance capacity

\paragraph{Methods and Design} difference-in-differences. Difference-in-differences. Data / sample: Difference-in-differences models find that when security headquarters were duplicated through an administrative reform, the proliferation of higherlevel posts within the service caused a large and statistically...

\paragraph{Findings} Difference-in-differences models find that when security headquarters were duplicated through an administrative reform, the proliferation of higherlevel posts within the service caused a large and statistically significant increase in the number of informants it employed. Following the research design of this paper, I suggest that studying the effects of exogenous shocks to authoritarian bureaucracies is a promising strategy for analyzing intra-agency structures and reforms.

\paragraph{Contribution} Provides an external benchmark for mechanisms that may or may not travel to China and WeChat.

\paragraph{Limitations} Design details need closer verification before the paper is used for strong causal claims.

\paragraph{Relevance to WeChat / My Project} Peripheral for a WeChat project; best used as a comparative benchmark rather than core theoretical evidence.

\paragraph{Related Debate} How findings from internet / online media and poland settings travel to a WeChat-centered China project


\subsection{Zhao and Zhang 2024: Visual propaganda in chinese central and...}
\textbf{Citation key}: \texttt{zhao\_2024\_humanities\_and\_social\_sciences\_communications}.\\
\textbf{Full citation}: Jiaye Zhao and Dechun Zhang 2024. "Visual propaganda in chinese central and local news agencies: a douyin case study." Humanities and Social Sciences Communications.

\paragraph{Summary} This study examines visual propaganda in chinese central and local news agencies: a douyin case study using descriptive or unclear from available parse evidence and reports that open Access This article is licensed under a Creative Commons Attribution 4.0 International License, which permits use, sharing,...

\paragraph{Context} ngagement with social media has increased, leading governments to expand their presence on social media platforms (Lu and Pan 2021; Zhang et al. 2020; Zhang and Xu 2023).

\paragraph{Main Argument} The paper's core mechanism is that open Access This article is licensed under a Creative Commons Attribution 4.0 International License, which permits use, sharing, adaptation, distribution and reproduction in any medium or format, as long as you give...

\paragraph{Independent Variable(s)} Not stated in standard IV/DV terms; inferred from title/abstract if needed

\paragraph{Dependent Variable(s)} Not stated in standard IV/DV terms; verify if central to argument

\paragraph{Methods and Design} Descriptive or unclear from available parse. Observational or descriptive design; verify details from full text if heavily used. Data / sample: Unclear from available parse; verify from full text if used centrally

\paragraph{Findings} Open Access This article is licensed under a Creative Commons Attribution 4.0 International License, which permits use, sharing, adaptation, distribution and reproduction in any medium or format, as long as you give appropriate credit to the original author(s) and the source, provide a link to the Creative Commons licence, and indicate if changes were made.

\paragraph{Contribution} Helps identify which mechanisms in the Chinese information environment are platform-specific and which generalize beyond a single app.

\paragraph{Limitations} Design details need closer verification before the paper is used for strong causal claims.

\paragraph{Relevance to WeChat / My Project} Conceptually useful for a WeChat project because it speaks to the Chinese information environment, especially the cluster on censorship, propaganda, and state influence, even if the platform is not WeChat.

\paragraph{Related Debate} Whether digital control works through persuasion, selective amplification, covert influence, or audience management


\subsection{Berger et al. 2025: Debunking fake news on social media...}
\textbf{Citation key}: \texttt{berger\_2025\_jpubeco}.\\
\textbf{Full citation}: Lara Marie Berger, Anna Kerkhof, Felix Mindl, and Johannes Munster 2025. "Debunking "fake news" on social media: Immediate and short-term effects of fact-checking and media literacy interventions." Journal of Public Economics.

\paragraph{Summary} This study examines debunking "fake news" on social media: immediate and short-term effects of fact-checking and media literacy... using survey experiment evidence and reports that we show that fact-checking primarily affects the specific fake news it directly addresses, whereas media literacy helps to distinguish...

\paragraph{Context} A plausible mechanism is that media literacy enables participants to critically evaluate social media postings, while fact-checking fails to enhance their skills as much.

\paragraph{Main Argument} The paper's core mechanism is that we show that fact-checking primarily affects the specific fake news it directly addresses, whereas media literacy helps to distinguish between false and correct information more generally, both immediately and around two...

\paragraph{Independent Variable(s)} Not stated in standard IV/DV terms; inferred from title/abstract if needed

\paragraph{Dependent Variable(s)} Not stated in standard IV/DV terms; verify if central to argument

\paragraph{Methods and Design} survey experiment. survey experiment. Data / sample: We conduct a randomized survey experiment to compare the immediate and short-term effects of fact-checking to a brief media literacy intervention.

\paragraph{Findings} We show that fact-checking primarily affects the specific fake news it directly addresses, whereas media literacy helps to distinguish between false and correct information more generally, both immediately and around two weeks after the intervention. Third and relatedly, while we demonstrate that media literacy interventions could help users to better distinguish between false and correct information that they encounter online, we remain agnostic about the concrete implementation of such interventions on behalf of social media platforms.

\paragraph{Contribution} Offers comparatively strong leverage on causal mechanisms in digital political communication.

\paragraph{Limitations} Internal validity is strong, but external validity may depend on whether experimental treatments resemble real platform use.

\paragraph{Relevance to WeChat / My Project} Indirect benchmark: useful for comparison on platform effects or official communication, but needs careful translation to China's WeChat-centered environment.

\paragraph{Related Debate} How findings from facebook and germany settings travel to a WeChat-centered China project


\subsection{Bollen et al. 2025: The value of dignity appeals evidence...}
\textbf{Citation key}: \texttt{bollen\_2025\_psrm}.\\
\textbf{Full citation}: Paige Bollen, Will Kymlicka, Evan Lieberman, and Blair Read 2025. "The value of dignity appeals: evidence from a social media experiment." Political Science Research and Methods.

\paragraph{Summary} This study examines the value of dignity appeals: evidence from a social media experiment using descriptive or unclear from available parse evidence and reports that consistent with preregistered hypotheses and specifications, we find that adding dignity appeals increases the likelihood of positive reactions to...

\paragraph{Context} In recent decades, activists and leaders of government and nongovernment organizations have increasingly and explicitly called for greater attention to human dignity in their efforts to promote pro-social relations. In this study, we investigate whether...

\paragraph{Main Argument} The paper's core mechanism is that consistent with preregistered hypotheses and specifications, we find that adding dignity appeals increases the likelihood of positive reactions to such ads, but only when the vulnerable are considered less "blameworthy"...

\paragraph{Independent Variable(s)} dignity appeals

\paragraph{Dependent Variable(s)} their likelihood of engaging with content concerning people facing homelessness or incarceration

\paragraph{Methods and Design} Descriptive or unclear from available parse. Observational or descriptive design; verify details from full text if heavily used. Data / sample: Using the digital advertising platform on Facebook, we randomly assign ads to over 90,000 adult American users to estimate the effects of dignity appeals on their likelihood of engaging with content concerning people...

\paragraph{Findings} Consistent with preregistered hypotheses and specifications, we find that adding dignity appeals increases the likelihood of positive reactions to such ads, but only when the vulnerable are considered less "blameworthy" for their situation. https://doi.org/10.1017/psrm.2025.10022 Published online by Cambridge University Press from the United States before and during the Covid-19 pandemic.

\paragraph{Contribution} Provides an external benchmark for mechanisms that may or may not travel to China and WeChat.

\paragraph{Limitations} Design details need closer verification before the paper is used for strong causal claims.

\paragraph{Relevance to WeChat / My Project} Indirect benchmark: useful for comparison on platform effects or official communication, but needs careful translation to China's WeChat-centered environment.

\paragraph{Related Debate} How findings from facebook and united states settings travel to a WeChat-centered China project


\subsection{Braghieri et al. 2025: Article-Level Slant and Polarization of News...}
\textbf{Citation key}: \texttt{braghieri\_2025\_wp}.\\
\textbf{Full citation}: Luca Braghieri, Sarah Eichmeyer, Ro'ee Levy, Markus Mobius, Jacob Steinhardt, and Ruiqi Zhong 2025. Article-Level Slant and Polarization of News Consumption on Social Media. Working paper.

\paragraph{Summary} This study examines article-level slant and polarization of news consumption on social media using panel / observational analysis evidence and reports that we find that polarized news consumption on Facebook, defined as the difference in the average slant of articles consumed by Republicans and Democrats on the...

\paragraph{Context} How polarized is news consumption on social media? We fine-tune a large-language-model to assign a content-based measure of slant to the near-universe of hard-news articles published online by the top 100 U.S. outlets in 2019.

\paragraph{Main Argument} The paper's core mechanism is that we find that polarized news consumption on Facebook, defined as the difference in the average slant of articles consumed by Republicans and Democrats on the platform, is arguably high. Analyzing the distribution of slant...

\paragraph{Independent Variable(s)} Not stated in standard IV/DV terms; inferred from title/abstract if needed

\paragraph{Dependent Variable(s)} Not stated in standard IV/DV terms; verify if central to argument

\paragraph{Methods and Design} panel / observational analysis. Observational or descriptive design; verify details from full text if heavily used. Data / sample: We fine-tune a large-language-model to assign a content-based measure of slant to the near-universe of hard-news articles published online by the top 100 U.S. outlets in 2019. We find that polarized news consumption...

\paragraph{Findings} We find that polarized news consumption on Facebook, defined as the difference in the average slant of articles consumed by Republicans and Democrats on the platform, is arguably high. Analyzing the distribution of slant on the production side, we find that 67\% of the article-level variation in slant is within rather than across outlets. * Braghieri: Bocconi University, CEPR, CESifo, and IGIER.

\paragraph{Contribution} Provides an external benchmark for mechanisms that may or may not travel to China and WeChat.

\paragraph{Limitations} Design details need closer verification before the paper is used for strong causal claims.

\paragraph{Relevance to WeChat / My Project} Indirect benchmark: useful for comparison on platform effects or official communication, but needs careful translation to China's WeChat-centered environment.

\paragraph{Related Debate} How social ties and discussion networks mediate willingness to speak, share, and update opinions


\subsection{Bursztyn et al. 2025: When Product Markets Become Collective Traps...}
\textbf{Citation key}: \texttt{bursztyn\_2025\_aer}.\\
\textbf{Full citation}: Leonardo Bursztyn, Benjamin Handel, Rafael Jimenez-Duran, and Christopher Roth 2025. "When Product Markets Become Collective Traps: The Case of Social Media." American Economic Review.

\paragraph{Summary} This paper studies when product markets become collective traps: the case of social media. .---We inform all respondents that we will conduct a deactivation study in which we will ask students at their university to deactivate their social media accounts for four weeks in exchange for monetary compensation.

\paragraph{Context} .---We inform all respondents that we will conduct a deactivation study in which we will ask students at their university to deactivate their social media accounts for four weeks in exchange for monetary compensation. To enhance the credibility of our...

\paragraph{Main Argument} Core mechanism maps onto the cluster on trust, credibility, rumors, and misinformation, but the exact statement should be verified from the full text if heavily used.

\paragraph{Independent Variable(s)} Not stated in standard IV/DV terms; inferred from title/abstract if needed

\paragraph{Dependent Variable(s)} Not stated in standard IV/DV terms; verify if central to argument

\paragraph{Methods and Design} Descriptive or unclear from available parse. Observational or descriptive design; verify details from full text if heavily used. Data / sample: .---We inform all respondents that we will conduct a deactivation study in which we will ask students at their university to deactivate their social media accounts for four weeks in exchange for monetary compensation....

\paragraph{Findings} No clean finding sentence detected; verify from full text if central.

\paragraph{Contribution} Provides an external benchmark for mechanisms that may or may not travel to China and WeChat.

\paragraph{Limitations} Design details need closer verification before the paper is used for strong causal claims.

\paragraph{Relevance to WeChat / My Project} Indirect benchmark: useful for comparison on platform effects or official communication, but needs careful translation to China's WeChat-centered environment.

\paragraph{Related Debate} Which actors can credibly correct misinformation and under what informational conditions


\subsection{Chau et al. 2025: Communication coordination and surveillance in the...}
\textbf{Citation key}: \texttt{chau\_2025\_ajps}.\\
\textbf{Full citation}: Tak-Huen Chau, Mai Hassan, and Andrew T. Little 2025. "Communication, coordination, and surveillance in the shadow of repression." American Journal of Political Science.

\paragraph{Summary} This study examines communication, coordination, and surveillance in the shadow of repression using descriptive or unclear from available parse evidence and reports that we study this trade-off in a model where protesters want to show up at the same time and place, but also want to avoid government forces.

\paragraph{Context} Communication technology helps protesters organize, but also allows the government to monitor and repress their actions. We study this trade-off in a model where protesters want to show up at the same time and place, but also want to avoid government forces.

\paragraph{Main Argument} The paper's core mechanism is that we study this trade-off in a model where protesters want to show up at the same time and place, but also want to avoid government forces.

\paragraph{Independent Variable(s)} Not stated in standard IV/DV terms; inferred from title/abstract if needed

\paragraph{Dependent Variable(s)} Not stated in standard IV/DV terms; verify if central to argument

\paragraph{Methods and Design} Descriptive or unclear from available parse. Observational or descriptive design; verify details from full text if heavily used. Data / sample: If leaders of a movement can send messages observed only by other protesters, they can successfully coordinate on a variety of sites and force the government to spread resources thin, helping the success of the...

\paragraph{Findings} We study this trade-off in a model where protesters want to show up at the same time and place, but also want to avoid government forces.

\paragraph{Contribution} Provides an external benchmark for mechanisms that may or may not travel to China and WeChat.

\paragraph{Limitations} Mechanisms may not transfer cleanly to WeChat's hybrid private-public architecture without further adaptation.

\paragraph{Relevance to WeChat / My Project} Peripheral for a WeChat project; best used as a comparative benchmark rather than core theoretical evidence.

\paragraph{Related Debate} How social ties and discussion networks mediate willingness to speak, share, and update opinions


\subsection{Chen et al. 2025: On the political consequences of local...}
\textbf{Citation key}: \texttt{chen\_2025\_psrm}.\\
\textbf{Full citation}: Jidong Chen, Yukun Wang, Ming-ang Zhang, and Xingyu Zhou 2025. "On the political consequences of local deliberative governance in China." Political Science Research and Methods.

\paragraph{Summary} This study examines on the political consequences of local deliberative governance in china using difference-in-differences evidence and reports that empirical results based on a staggered difference-in-differences design suggest that broadcasting such programs in China produces a prompt increase in citizens' trust...

\paragraph{Context} How can local governments in developing countries, constrained by limited resources, identify and respond to the most pressing public demands? This paper posits that public deliberative platforms, even those with controlled agendas, can be instrumental in...

\paragraph{Main Argument} The paper's core mechanism is that empirical results based on a staggered difference-in-differences design suggest that broadcasting such programs in China produces a prompt increase in citizens' trust in local officials. Our results demonstrate that...

\paragraph{Independent Variable(s)} Not stated in standard IV/DV terms; inferred from title/abstract if needed

\paragraph{Dependent Variable(s)} Not stated in standard IV/DV terms; verify if central to argument

\paragraph{Methods and Design} difference-in-differences. Difference-in-differences. Data / sample: How can local governments in developing countries, constrained by limited resources, identify and respond to the most pressing public demands? Our results demonstrate that public deliberation can yield noticeable...

\paragraph{Findings} Empirical results based on a staggered difference-in-differences design suggest that broadcasting such programs in China produces a prompt increase in citizens' trust in local officials. Our results demonstrate that public deliberation can yield noticeable outcomes in developing countries, even with controlled agendas and constrained resources. https://doi.org/10.1017/psrm.2025.10037 Published online by Cambridge University Press of how our work relates to the seminal study by He and Warren (2011) is provided in Section A8 of the

\paragraph{Contribution} Helps identify which mechanisms in the Chinese information environment are platform-specific and which generalize beyond a single app.

\paragraph{Limitations} Design details need closer verification before the paper is used for strong causal claims.

\paragraph{Relevance to WeChat / My Project} Conceptually useful for a WeChat project because it speaks to the Chinese information environment, especially the cluster on trust, credibility, rumors, and misinformation, even if the platform is not WeChat.

\paragraph{Related Debate} Which actors can credibly correct misinformation and under what informational conditions


\subsection{Eady et al. 2025: News Sharing on Social Media Mapping...}
\textbf{Citation key}: \texttt{eady\_2025\_pan}.\\
\textbf{Full citation}: Gregory Eady, Richard Bonneau, Joshua A. Tucker, and Jonathan Nagler 2025. "News Sharing on Social Media: Mapping the Ideology of News Media, Politicians, and the Mass Public." Political Analysis.

\paragraph{Summary} This study examines news sharing on social media: mapping the ideology of news media, politicians, and the mass public using descriptive or unclear from available parse evidence and reports that empirically, we show that (1) politicians who share ideologically polarized content share, by far, the most political...

\paragraph{Context} This article examines the information sharing behavior of U.S. politicians and the mass public by mapping the ideological sharing space of political news on social media.

\paragraph{Main Argument} The paper's core mechanism is that empirically, we show that (1) politicians who share ideologically polarized content share, by far, the most political news and commentary and (2) that the less competitive elections are, the more likely politicians are to...

\paragraph{Independent Variable(s)} Not stated in standard IV/DV terms; inferred from title/abstract if needed

\paragraph{Dependent Variable(s)} Not stated in standard IV/DV terms; verify if central to argument

\paragraph{Methods and Design} Descriptive or unclear from available parse. Observational or descriptive design; verify details from full text if heavily used. Data / sample: As data, we use the near-universal currency of online information exchange: web links. We introduce a methodological approach and software to unify the measurement of ideology across social media platforms by using...

\paragraph{Findings} Empirically, we show that (1) politicians who share ideologically polarized content share, by far, the most political news and commentary and (2) that the less competitive elections are, the more likely politicians are to share polarized information. These results demonstrate that news and commentary shared by politicians come from a highly unrepresentative set of ideologically extreme legislators and that decreases in election pressures (e.g., by gerrymandering) may encourage polarized sharing behavior. suggest that election competition may act as a constraint on politicians from sharing ideologically extreme news media.

\paragraph{Contribution} Provides an external benchmark for mechanisms that may or may not travel to China and WeChat.

\paragraph{Limitations} Design details need closer verification before the paper is used for strong causal claims.

\paragraph{Relevance to WeChat / My Project} Indirect benchmark: useful for comparison on platform effects or official communication, but needs careful translation to China's WeChat-centered environment.

\paragraph{Related Debate} How interface design changes the visibility and meaning of digital expression


\subsection{Fan et al. 2025: Bringing in horizontal strategic interactions Blame...}
\textbf{Citation key}: \texttt{fan\_2025\_government\_information\_quarterly}.\\
\textbf{Full citation}: Ziteng Fan, Yijia Jing, and Shaowei Chen 2025. "Bringing in horizontal strategic interactions: Blame avoidance and local governments' bandwagon strategy in prioritizing E-participation." Government Information Quarterly.

\paragraph{Summary} This paper studies bringing in horizontal strategic interactions: blame avoidance and local governments' bandwagon strategy in.... E-participation has gained widespread attention for its potential to facilitate citizen--state interactions and achieve inclusive governance.

\paragraph{Context} However, the effectiveness of e-participation often falls short of expec tations (Norris, 2010; Chadwick, 2011; Panagiotopoulos et al., 2012; Sandoval-Almazan \& Gil-Garcia, 2012; Janssen et al., 2015; Medaglia \& Zhu, 2017; Reddick and Norris, 2013;...

\paragraph{Main Argument} Core mechanism maps onto the cluster on comparative digital media benchmarks, but the exact statement should be verified from the full text if heavily used.

\paragraph{Independent Variable(s)} Not stated in standard IV/DV terms; inferred from title/abstract if needed

\paragraph{Dependent Variable(s)} Not stated in standard IV/DV terms; verify if central to argument

\paragraph{Methods and Design} survey. Survey-based observational analysis. Data / sample: According to an egovernment survey by the United Nations (2020), expanding e-partici pation opportunities in many countries has yet to result in wider or more meaningful public engagement. , and provincial official...

\paragraph{Findings} No clean finding sentence detected; verify from full text if central.

\paragraph{Contribution} Helps identify which mechanisms in the Chinese information environment are platform-specific and which generalize beyond a single app.

\paragraph{Limitations} Design details need closer verification before the paper is used for strong causal claims.

\paragraph{Relevance to WeChat / My Project} Conceptually useful for a WeChat project because it speaks to the Chinese information environment, especially the cluster on comparative digital media benchmarks, even if the platform is not WeChat.

\paragraph{Related Debate} How findings from internet / online media and china settings travel to a WeChat-centered China project


\subsection{Garbe et al. 2025: Who Wants to be Legible Digitalization...}
\textbf{Citation key}: \texttt{garbe\_2025\_cps}.\\
\textbf{Full citation}: Lisa Garbe, Nina McMurry, Alexandra Scacco, and Kelly Zhang 2025. "Who Wants to be Legible? Digitalization and Intergroup Inequality in Kenya." Comparative Political Studies.

\paragraph{Summary} This paper studies who wants to be legible? digitalization and intergroup inequality in kenya. Governments across the Global South have begun introducing biometric IDs (eIDs) in an attempt to improve citizen-state legibility.

\paragraph{Context} While such initiatives can improve government efficiency, they also raise important questions about citizen privacy, especially for groups with a history of mistrust in the state.

\paragraph{Main Argument} Core mechanism maps onto the cluster on participation, responsiveness, and political behavior, but the exact statement should be verified from the full text if heavily used.

\paragraph{Independent Variable(s)} Author inference from title/abstract: Who Wants to be Legible? Digitalization and Intergroup Inequality in Kenya

\paragraph{Dependent Variable(s)} Not stated in standard IV/DV terms

\paragraph{Methods and Design} Descriptive or unclear from available parse. Observational or descriptive design; verify details from full text if heavily used. Data / sample: In a conjoint experiment with 2,072 respondents from four Kenyan regions, we examine how perceptions of and willingness to register for eID under different policy conditions vary across politically dominant,...

\paragraph{Findings} No clean finding sentence detected; verify from full text if central.

\paragraph{Contribution} Provides an external benchmark for mechanisms that may or may not travel to China and WeChat.

\paragraph{Limitations} Mechanisms may not transfer cleanly to WeChat's hybrid private-public architecture without further adaptation.

\paragraph{Relevance to WeChat / My Project} Peripheral for a WeChat project; best used as a comparative benchmark rather than core theoretical evidence.

\paragraph{Related Debate} Whether official digital communication mainly informs, disciplines, mobilizes, or improves responsiveness


\subsection{Gohdes and Steinert-Threlkeld 2025: Civilian behavior on social media during...}
\textbf{Citation key}: \texttt{gohdes\_2025\_ajps}.\\
\textbf{Full citation}: Anita R. Gohdes and Zachary C. Steinert-Threlkeld 2025. "Civilian behavior on social media during civil war." American Journal of Political Science.

\paragraph{Summary} This study examines civilian behavior on social media during civil war using descriptive or unclear from available parse evidence and reports that results show that users in Aleppo post more positive and pro-Assad content, but only when self-disclosing their location. From isolated protests to countrywide uprisings...

\paragraph{Context} In active conflict, however, social media posts indicating political loyalties can pose severe risks to civilians.

\paragraph{Main Argument} The paper's core mechanism is that results show that users in Aleppo post more positive and pro-Assad content, but only when self-disclosing their location. From isolated protests to countrywide uprisings to organized armed conflict, civilians and...

\paragraph{Independent Variable(s)} civilians modify their online behavior as part of efforts to improve their security during conflict. After major

\paragraph{Dependent Variable(s)} in territorial control

\paragraph{Methods and Design} Descriptive or unclear from available parse. Observational or descriptive design; verify details from full text if heavily used. Data / sample: In active conflict, however, social media posts indicating political loyalties can pose severe risks to civilians. We match Aleppo-based Twitter users with users from other parts of Syria and use large language...

\paragraph{Findings} Results show that users in Aleppo post more positive and pro-Assad content, but only when self-disclosing their location. From isolated protests to countrywide uprisings to organized armed conflict, civilians and government actors routinely use social media to document and participate in such events.

\paragraph{Contribution} Provides an external benchmark for mechanisms that may or may not travel to China and WeChat.

\paragraph{Limitations} Design details need closer verification before the paper is used for strong causal claims.

\paragraph{Relevance to WeChat / My Project} Indirect benchmark: useful for comparison on platform effects or official communication, but needs careful translation to China's WeChat-centered environment.

\paragraph{Related Debate} How interface design changes the visibility and meaning of digital expression


\subsection{Graham 2025: Techniques for Measuring Political Beliefs}
\textbf{Citation key}: \texttt{graham\_2025\_wp}.\\
\textbf{Full citation}: Matthew H. Graham 2025. Techniques for Measuring Political Beliefs. Oxford University Press.

\paragraph{Summary} This paper studies techniques for measuring political beliefs. Political beliefs are generally measured using surveys.

\paragraph{Context} Political beliefs are generally measured using surveys. This chapter organizes prevailing techniques around four concerns about why respondents may not "believe" their answers.

\paragraph{Main Argument} Core mechanism maps onto the cluster on comparative digital media benchmarks, but the exact statement should be verified from the full text if heavily used.

\paragraph{Independent Variable(s)} Not stated in standard IV/DV terms; inferred from title/abstract if needed

\paragraph{Dependent Variable(s)} Not stated in standard IV/DV terms; verify if central to argument

\paragraph{Methods and Design} survey. Survey-based observational analysis. Data / sample: This chapter organizes prevailing techniques around four concerns about why respondents may not "believe" their answers. Researchers who worry that respondents are being insincere use techniques designed to encourage...

\paragraph{Findings} No clean finding sentence detected; verify from full text if central.

\paragraph{Contribution} Provides an external benchmark for mechanisms that may or may not travel to China and WeChat.

\paragraph{Limitations} Mechanisms may not transfer cleanly to WeChat's hybrid private-public architecture without further adaptation.

\paragraph{Relevance to WeChat / My Project} Peripheral for a WeChat project; best used as a comparative benchmark rather than core theoretical evidence.

\paragraph{Related Debate} How findings from not platform-specific and comparative / general settings travel to a WeChat-centered China project


\subsection{Han 2025: Propaganda State 2 0 in China}
\textbf{Citation key}: \texttt{han\_2025\_cq}.\\
\textbf{Full citation}: Rongbin Han 2025. ""Propaganda State 2.0" in China." The China Quarterly.

\paragraph{Summary} This paper studies "propaganda state 2.0" in china. By observing China's domestic media landscape and state policies during the COVID-19 pandemic, this paper proposes the concept of "propaganda state 2.0" as a framework for exploring autocratic state propaganda from a holistic perspective.

\paragraph{Context} Theoretically, the article highlights how state propaganda in contemporary China can shape, if not dictate, state policy, while serving as a more organic framework that bridges the "hard" and "soft" propaganda literatures.  ,...

\paragraph{Main Argument} Core mechanism maps onto the cluster on censorship, propaganda, and state influence, but the exact statement should be verified from the full text if heavily used.

\paragraph{Independent Variable(s)} Not stated in standard IV/DV terms; inferred from title/abstract if needed

\paragraph{Dependent Variable(s)} Not stated in standard IV/DV terms; verify if central to argument

\paragraph{Methods and Design} Descriptive or unclear from available parse. Observational or descriptive design; verify details from full text if heavily used. Data / sample: This system performs both conventional persuasion and indoctrination functions, as well as fulfilling the now underestimated mission of agitation, which can enhance the credibility of propaganda messages. Indeed,...

\paragraph{Findings} No clean finding sentence detected; verify from full text if central.

\paragraph{Contribution} Helps identify which mechanisms in the Chinese information environment are platform-specific and which generalize beyond a single app.

\paragraph{Limitations} Mechanisms may not transfer cleanly to WeChat's hybrid private-public architecture without further adaptation.

\paragraph{Relevance to WeChat / My Project} Conceptually useful for a WeChat project because it speaks to the Chinese information environment, especially the cluster on censorship, propaganda, and state influence, even if the platform is not WeChat.

\paragraph{Related Debate} Whether digital control works through persuasion, selective amplification, covert influence, or audience management


\subsection{Huang et al. 2025: CORRELATION NEGLECT ON SOCIAL MEDIA EFFECTS...}
\textbf{Citation key}: \texttt{huang\_2025\_ej}.\\
\textbf{Full citation}: Yihong Huang, Yixi Jiang, and Ziqi Lu 2025. "CORRELATION NEGLECT ON SOCIAL MEDIA: EFFECTS ON CIVIL SERVICE APPLICATIONS IN CHINA*." The Economic Journal.

\paragraph{Summary} This paper studies correlation neglect on social media: effects on civil service applications in china*. tasks, with further lab-based studies suggesting its field relevance by applying it to real-world decisions, such as school selection and portfolio management (Eyster and Weizsacker, 2016; Rees-Jones et al., 2022).

\paragraph{Context} However, our sampling of Weibo posts by current civil servants shows that 46\% of posts are negative, compared to 36\% of posts that are neutral, and 19\% of posts that are positive.3 This suggests that the public's favourable view of these jobs may not align...

\paragraph{Main Argument} Core mechanism maps onto the cluster on comparative digital media benchmarks, but the exact statement should be verified from the full text if heavily used.

\paragraph{Independent Variable(s)} Not stated in standard IV/DV terms; inferred from title/abstract if needed

\paragraph{Dependent Variable(s)} Not stated in standard IV/DV terms; verify if central to argument

\paragraph{Methods and Design} Descriptive or unclear from available parse. Observational or descriptive design; verify details from full text if heavily used. Data / sample: Most participants (75\%) in our sample are interested in the national exam. However, our sampling of Weibo posts by current civil servants shows that 46\% of posts are negative, compared to 36\% of posts that are...

\paragraph{Findings} No clean finding sentence detected; verify from full text if central.

\paragraph{Contribution} Helps identify which mechanisms in the Chinese information environment are platform-specific and which generalize beyond a single app.

\paragraph{Limitations} Design details need closer verification before the paper is used for strong causal claims.

\paragraph{Relevance to WeChat / My Project} Conceptually useful for a WeChat project because it speaks to the Chinese information environment, especially the cluster on comparative digital media benchmarks, even if the platform is not WeChat.

\paragraph{Related Debate} How findings from weibo and china settings travel to a WeChat-centered China project


\subsection{Karpa and Rochlitz 2025: Authoritarian Surveillance and Public Support for...}
\textbf{Citation key}: \texttt{karpa\_2025\_cps}.\\
\textbf{Full citation}: David Karpa and Michael Rochlitz 2025. "Authoritarian Surveillance and Public Support for Digital Governance Solutions." Comparative Political Studies.

\paragraph{Summary} This study examines authoritarian surveillance and public support for digital governance solutions using survey evidence and reports that recent survey data show a big variation in trust in the police across countries, with 86\% of respondents trusting the police in Estonia, 49\% in Germany, 48\% in the US, 39\% in...

\paragraph{Context} Importantly, while this effect has previously been documented for China, we find it irrespective of regime type for an autocracy, a hybrid regime and three democracies.

\paragraph{Main Argument} The paper's core mechanism is that recent survey data show a big variation in trust in the police across countries, with 86\% of respondents trusting the police in Estonia, 49\% in Germany, 48\% in the US, 39\% in Turkey, but only 20\% in Russia in 2022.

\paragraph{Independent Variable(s)} Not stated in standard IV/DV terms; inferred from title/abstract if needed

\paragraph{Dependent Variable(s)} Not stated in standard IV/DV terms; verify if central to argument

\paragraph{Methods and Design} survey. Survey-based observational analysis. Data / sample: We conduct survey experiments in Russia, Germany, Turkey, the United States, and Estonia, and find that awareness of potential misuse of digital governance tools by the government reduces support. Instead,...

\paragraph{Findings} Recent survey data show a big variation in trust in the police across countries, with 86\% of respondents trusting the police in Estonia, 49\% in Germany, 48\% in the US, 39\% in Turkey, but only 20\% in Russia in 2022.

\paragraph{Contribution} Provides an external benchmark for mechanisms that may or may not travel to China and WeChat.

\paragraph{Limitations} Mechanisms may not transfer cleanly to WeChat's hybrid private-public architecture without further adaptation.

\paragraph{Relevance to WeChat / My Project} Peripheral for a WeChat project; best used as a comparative benchmark rather than core theoretical evidence.

\paragraph{Related Debate} Whether official digital communication mainly informs, disciplines, mobilizes, or improves responsiveness


\subsection{Larreguy and Raffler 2025: Accountability in Developing Democracies The Impact...}
\textbf{Citation key}: \texttt{larreguy\_2025\_arps}.\\
\textbf{Full citation}: Horacio Larreguy and Pia J. Raffler 2025. "Accountability in Developing Democracies: The Impact of the Internet, Social Media, and Polarization." Annual Review of Political Science.

\paragraph{Summary} This study examines accountability in developing democracies: the impact of the internet, social media, and polarization using descriptive or unclear from available parse evidence and reports that rather than offering definitive conclusions, we highlight key trade tradeoffs and emphasize that the overall effects...

\paragraph{Context} See credit lines of images or other third-party material in this article for license information.

\paragraph{Main Argument} The paper's core mechanism is that rather than offering definitive conclusions, we highlight key trade tradeoffs and emphasize that the overall effects are contingent on the status quo.

\paragraph{Independent Variable(s)} the Internet, Social Media, and Polarization. https://doi.org/10.1146/annurev-polisci-033123015559 Copyright  2025...

\paragraph{Dependent Variable(s)} accountability in developing democracies through the lens of two nested principal--agent problems: the relationship...

\paragraph{Methods and Design} Descriptive or unclear from available parse. Observational or descriptive design; verify details from full text if heavily used. Data / sample: Unclear from available parse; verify from full text if used centrally

\paragraph{Findings} Rather than offering definitive conclusions, we highlight key trade tradeoffs and emphasize that the overall effects are contingent on the status quo.

\paragraph{Contribution} Provides an external benchmark for mechanisms that may or may not travel to China and WeChat.

\paragraph{Limitations} Design details need closer verification before the paper is used for strong causal claims.

\paragraph{Relevance to WeChat / My Project} Indirect benchmark: useful for comparison on platform effects or official communication, but needs careful translation to China's WeChat-centered environment.

\paragraph{Related Debate} Whether official digital communication mainly informs, disciplines, mobilizes, or improves responsiveness


\subsection{Liu 2025: How Allowing a Little Bit of...}
\textbf{Citation key}: \texttt{liu\_2025\_jeea}.\\
\textbf{Full citation}: Zhenqi (Jessie) Liu 2025. "How Allowing a Little Bit of Dissent Helps Control Social Media: Impact of Market Structure on Censorship Compliance." Journal of the European Economic Association.

\paragraph{Summary} This paper studies how allowing a little bit of dissent helps control social media: impact of market structure on censorship compliance. This paper studies the role of market structure in regulatory compliance through a unique empirical example: censorship via content removal by three major live-streaming platforms...

\paragraph{Context} This paper studies the role of market structure in regulatory compliance through a unique empirical example: censorship via content removal by three major live-streaming platforms in China. Based on 30 unexpected sensitive events, I first present...

\paragraph{Main Argument} Core mechanism maps onto the cluster on censorship, propaganda, and state influence, but the exact statement should be verified from the full text if heavily used.

\paragraph{Independent Variable(s)} Market Structure

\paragraph{Dependent Variable(s)} Censorship Compliance

\paragraph{Methods and Design} survey. Survey-based observational analysis. Data / sample: Beyond China, other countries have also been pushing for stricter regulations on social media platforms following data breach scandals and instances of terrorist groups using social media as a recruiting device...

\paragraph{Findings} No clean finding sentence detected; verify from full text if central.

\paragraph{Contribution} Helps identify which mechanisms in the Chinese information environment are platform-specific and which generalize beyond a single app.

\paragraph{Limitations} Design details need closer verification before the paper is used for strong causal claims.

\paragraph{Relevance to WeChat / My Project} Conceptually useful for a WeChat project because it speaks to the Chinese information environment, especially the cluster on censorship, propaganda, and state influence, even if the platform is not WeChat.

\paragraph{Related Debate} Whether digital control works through persuasion, selective amplification, covert influence, or audience management


\subsection{Lu et al. 2025: Decentralized propaganda in the era of...}
\textbf{Citation key}: \texttt{lu\_2025\_ajps}.\\
\textbf{Full citation}: Yingdan Lu, Jennifer Pan, Xu Xu, and Yiqing Xu 2025. "Decentralized propaganda in the era of digital media: The massive presence of the Chinese state on Douyin." American Journal of Political Science.

\paragraph{Summary} This study examines decentralized propaganda in the era of digital media: the massive presence of the chinese state on douyin using descriptive or unclear from available parse evidence and reports that we empirically demonstrate the existence of this new model in China by creating a novel data set of over five...

\paragraph{Context} The rise of social media in the digital era poses unprecedented challenges to authoritarian regimes that aim to influence public attitudes and behaviors. To address these challenges, we argue that authoritarian regimes have adopted a decentralized approach...

\paragraph{Main Argument} The paper's core mechanism is that we empirically demonstrate the existence of this new model in China by creating a novel data set of over five million videos from over 18,000 regime-affiliated accounts on Douyin, a popular social media platform in China.

\paragraph{Independent Variable(s)} Not stated in standard IV/DV terms; inferred from title/abstract if needed

\paragraph{Dependent Variable(s)} Not stated in standard IV/DV terms; verify if central to argument

\paragraph{Methods and Design} Descriptive or unclear from available parse. Observational or descriptive design; verify details from full text if heavily used. Data / sample: We empirically demonstrate the existence of this new model in China by creating a novel data set of over five million videos from over 18,000 regime-affiliated accounts on Douyin, a popular social media platform in...

\paragraph{Findings} We empirically demonstrate the existence of this new model in China by creating a novel data set of over five million videos from over 18,000 regime-affiliated accounts on Douyin, a popular social media platform in China.

\paragraph{Contribution} Helps identify which mechanisms in the Chinese information environment are platform-specific and which generalize beyond a single app.

\paragraph{Limitations} Design details need closer verification before the paper is used for strong causal claims.

\paragraph{Relevance to WeChat / My Project} Conceptually useful for a WeChat project because it speaks to the Chinese information environment, especially the cluster on censorship, propaganda, and state influence, even if the platform is not WeChat.

\paragraph{Related Debate} Whether digital control works through persuasion, selective amplification, covert influence, or audience management


\subsection{Pedersen and Seeberg 2025: Using Social Media to Respond to...}
\textbf{Citation key}: \texttt{pedersen\_2025\_jop}.\\
\textbf{Full citation}: Helene Helboe Pedersen and Henrik Bech Seeberg 2025. "Using Social Media to Respond to Negative Polls: Politicians' Issue Responsiveness on Facebook." The Journal of Politics.

\paragraph{Summary} This paper studies using social media to respond to negative polls: politicians' issue responsiveness on facebook. Using Social Media to Respond to Negative Polls: Politicians' Issue Responsiveness on Facebook Helene Helboe Pedersen, Aarhus University Henrik Bech Seeberg, Aarhus University Social media has changed...

\paragraph{Context} Using Social Media to Respond to Negative Polls: Politicians' Issue Responsiveness on Facebook Helene Helboe Pedersen, Aarhus University Henrik Bech Seeberg, Aarhus University Social media has changed the link between politicians and voters. An unsettled...

\paragraph{Main Argument} Core mechanism maps onto the cluster on participation, responsiveness, and political behavior, but the exact statement should be verified from the full text if heavily used.

\paragraph{Independent Variable(s)} Not stated in standard IV/DV terms; inferred from title/abstract if needed

\paragraph{Dependent Variable(s)} Not stated in standard IV/DV terms; verify if central to argument

\paragraph{Methods and Design} Descriptive or unclear from available parse. Observational or descriptive design; verify details from full text if heavily used. Data / sample: Unclear from available parse; verify from full text if used centrally

\paragraph{Findings} No clean finding sentence detected; verify from full text if central.

\paragraph{Contribution} Provides an external benchmark for mechanisms that may or may not travel to China and WeChat.

\paragraph{Limitations} Design details need closer verification before the paper is used for strong causal claims.

\paragraph{Relevance to WeChat / My Project} Indirect benchmark: useful for comparison on platform effects or official communication, but needs careful translation to China's WeChat-centered environment.

\paragraph{Related Debate} Whether official digital communication mainly informs, disciplines, mobilizes, or improves responsiveness


\subsection{Rosenzweig and Offer-Westort 2025: Conversations with a concern-addressing chatbot increase...}
\textbf{Citation key}: \texttt{rosenzweig\_2025\_jop}.\\
\textbf{Full citation}: Leah Rosenzweig and Molly Offer-Westort 2025. "Conversations with a concern-addressing chatbot increase COVID-19 vaccination intentions among social media users in Kenya and Nigeria." The Journal of Politics.

\paragraph{Summary} This study examines conversations with a concern-addressing chatbot increase covid-19 vaccination intentions among social media users in... using descriptive or unclear from available parse evidence and reports that we find that the optimized concern-addressing messaging increases COVID-19 vaccine intentions and...

\paragraph{Context} This study leverages a two-stage response-adaptive experimental design to explore how tailored digital messaging can influence public attitudes toward health interventions, with broader implications for political communication. In public health settings,...

\paragraph{Main Argument} The paper's core mechanism is that we find that the optimized concern-addressing messaging increases COVID-19 vaccine intentions and willingness by 4-5\% compared to a control condition, and by 3-4\% compared to a public service announcement.

\paragraph{Independent Variable(s)} Not stated in standard IV/DV terms; inferred from title/abstract if needed

\paragraph{Dependent Variable(s)} Not stated in standard IV/DV terms; verify if central to argument

\paragraph{Methods and Design} Descriptive or unclear from available parse. Observational or descriptive design; verify details from full text if heavily used. Data / sample: For this study, we recruited 22,052 Facebook users in Kenya and Nigeria to engage in conversations on Messenger about their concerns regarding COVID-19 vaccines. Overall, the personalized concern-addressing chatbot...

\paragraph{Findings} We find that the optimized concern-addressing messaging increases COVID-19 vaccine intentions and willingness by 4-5\% compared to a control condition, and by 3-4\% compared to a public service announcement.

\paragraph{Contribution} Provides an external benchmark for mechanisms that may or may not travel to China and WeChat.

\paragraph{Limitations} Design details need closer verification before the paper is used for strong causal claims.

\paragraph{Relevance to WeChat / My Project} Indirect benchmark: useful for comparison on platform effects or official communication, but needs careful translation to China's WeChat-centered environment.

\paragraph{Related Debate} Whether official digital communication mainly informs, disciplines, mobilizes, or improves responsiveness


\subsection{Shen 2025: The effects of forced versus selective...}
\textbf{Citation key}: \texttt{shen\_2025\_psrm}.\\
\textbf{Full citation}: Xiaoxiao Shen 2025. "The effects of forced versus selective exposure to propaganda in China." Political Science Research and Methods.

\paragraph{Summary} This study examines the effects of forced versus selective exposure to propaganda in china using survey evidence and reports that but the study addressed in this article found that citizens who prefer not to view propaganda news, when given a choice, actually demonstrate higher average treatment effects on...

\paragraph{Context} This article describes how and why propaganda affected recipients differently in two distinct situations, namely forced exposure and selective exposure, when they received propaganda during a series of six original survey experiments conducted in China.

\paragraph{Main Argument} The paper's core mechanism is that but the study addressed in this article found that citizens who prefer not to view propaganda news, when given a choice, actually demonstrate higher average treatment effects on pro-regime attitudes compared to those who...

\paragraph{Independent Variable(s)} forced versus selective exposure to propaganda in China. This article describes how and why propaganda affected...

\paragraph{Dependent Variable(s)} pro-regime attitudes compared to those who willingly read propaganda news (i

\paragraph{Methods and Design} survey. Survey-based observational analysis. Data / sample: This article describes how and why propaganda affected recipients differently in two distinct situations, namely forced exposure and selective exposure, when they received propaganda during a series of six original...

\paragraph{Findings} But the study addressed in this article found that citizens who prefer not to view propaganda news, when given a choice, actually demonstrate higher average treatment effects on pro-regime attitudes compared to those who willingly read propaganda news (i.e. where participants in the control group were assigned a reading of non-propaganda news).

\paragraph{Contribution} Helps identify which mechanisms in the Chinese information environment are platform-specific and which generalize beyond a single app.

\paragraph{Limitations} Design details need closer verification before the paper is used for strong causal claims.

\paragraph{Relevance to WeChat / My Project} Conceptually useful for a WeChat project because it speaks to the Chinese information environment, especially the cluster on censorship, propaganda, and state influence, even if the platform is not WeChat.

\paragraph{Related Debate} Whether digital control works through persuasion, selective amplification, covert influence, or audience management


\subsection{Vecchiato and Munger 2025: Introducing the Visual Conjoint with an...}
\textbf{Citation key}: \texttt{vecchiato\_2025\_jeps}.\\
\textbf{Full citation}: Alessandro Vecchiato and Kevin Munger 2025. "Introducing the Visual Conjoint, with an Application to Candidate Evaluation on Social Media." Journal of Experimental Political Science.

\paragraph{Summary} This study examines introducing the visual conjoint, with an application to candidate evaluation on social media using descriptive or unclear from available parse evidence and reports that we demonstrate that the visual conjoint more effectively encodes image-based information and social endorsement information.

\paragraph{Context} Conjoint experiments have enabled scholars to understand the preferences of citizens in a variety of political contexts. We propose a method to modify the standard text-only "box conjoint" to make the treatment higher in external validity with respect to a...

\paragraph{Main Argument} The paper's core mechanism is that we demonstrate that the visual conjoint more effectively encodes image-based information and social endorsement information.

\paragraph{Independent Variable(s)} Not stated in standard IV/DV terms; inferred from title/abstract if needed

\paragraph{Dependent Variable(s)} Not stated in standard IV/DV terms; verify if central to argument

\paragraph{Methods and Design} Descriptive or unclear from available parse. Observational or descriptive design; verify details from full text if heavily used. Data / sample: Citizens frequently encounter political information encoded as images and in particular in the form of politicians' social media posts and profiles. Citizens frequently encounter political information encoded as...

\paragraph{Findings} We demonstrate that the visual conjoint more effectively encodes image-based information and social endorsement information.

\paragraph{Contribution} Provides an external benchmark for mechanisms that may or may not travel to China and WeChat.

\paragraph{Limitations} Design details need closer verification before the paper is used for strong causal claims.

\paragraph{Relevance to WeChat / My Project} Indirect benchmark: useful for comparison on platform effects or official communication, but needs careful translation to China's WeChat-centered environment.

\paragraph{Related Debate} How findings from twitter / x and comparative / general settings travel to a WeChat-centered China project


\subsection{Walk et al. 2025: Social Media Narratives Across Platforms in...}
\textbf{Citation key}: \texttt{walk\_2025\_jop}.\\
\textbf{Full citation}: Erin Walk, Elizabeth Parker-Magyar, Kiran Garimella, Ahmet Akbiyik, and Fotini Christia 2025. "Social Media Narratives Across Platforms in Conflict: Evidence from Syria." The Journal of Politics.

\paragraph{Summary} This study examines social media narratives across platforms in conflict: evidence from syria using descriptive or unclear from available parse evidence and reports that our findings show that complementary if divergent discussions of violence remain central even amid a period of relative de-escalation. , the...

\paragraph{Context} Social Media Narratives Across Platforms in Conflict: Evidence from Syria Erin Walk, University of Pennsylvania Elizabeth Parker-Magyar, Harvard University Kiran Garimella, Rutgers University Ahmet Akbiyik, Harvard Kennedy School Fotini Christia,...

\paragraph{Main Argument} The paper's core mechanism is that our findings show that complementary if divergent discussions of violence remain central even amid a period of relative de-escalation. , the relationship between the message, the platform, and the audience is triangular.

\paragraph{Independent Variable(s)} Not stated in standard IV/DV terms; inferred from title/abstract if needed

\paragraph{Dependent Variable(s)} Not stated in standard IV/DV terms; verify if central to argument

\paragraph{Methods and Design} Descriptive or unclear from available parse. Observational or descriptive design; verify details from full text if heavily used. Data / sample: We constructed and analyzed comparable datasets of public messages and images from elite- and meso-level Syrian actors posting on three popular social media platforms. Our findings show that complementary if divergent...

\paragraph{Findings} Our findings show that complementary if divergent discussions of violence remain central even amid a period of relative de-escalation. , the relationship between the message, the platform, and the audience is triangular.

\paragraph{Contribution} Provides an external benchmark for mechanisms that may or may not travel to China and WeChat.

\paragraph{Limitations} Design details need closer verification before the paper is used for strong causal claims.

\paragraph{Relevance to WeChat / My Project} Indirect benchmark: useful for comparison on platform effects or official communication, but needs careful translation to China's WeChat-centered environment.

\paragraph{Related Debate} How social ties and discussion networks mediate willingness to speak, share, and update opinions


\subsection{Wendsjo et al. 2025: A Male Hostility Spiral Polarized Communication...}
\textbf{Citation key}: \texttt{wendsjo\_2025\_jop}.\\
\textbf{Full citation}: Albert Wendsjo, Hanna Back, and Andrej Kokkonen 2025. "A Male Hostility Spiral? Polarized Communication among Political Elites on Social Media." The Journal of Politics.

\paragraph{Summary} This study examines a male hostility spiral? polarized communication among political elites on social media using descriptive or unclear from available parse evidence and reports that specifically, we use machine learning to measure the tone of over 200,000 twitter interactions and find that male politicians are...

\paragraph{Context} However, there is a lack of research analyzing gendered patterns of polarizing communication among political elites on social media.

\paragraph{Main Argument} The paper's core mechanism is that specifically, we use machine learning to measure the tone of over 200,000 twitter interactions and find that male politicians are more likely to attack political opponents representing 'outgroups' and that male...

\paragraph{Independent Variable(s)} Author inference from title/abstract: A Male Hostility Spiral? Polarized Communication among Political Elites on Social Media

\paragraph{Dependent Variable(s)} Not stated in standard IV/DV terms

\paragraph{Methods and Design} Descriptive or unclear from available parse. Observational or descriptive design; verify details from full text if heavily used. Data / sample: We focus on filling this gap, analyzing how politicians communicate on social media in 24 western countries. Cases and data To test our hypotheses, we use tweets made by politicians during the period 2017-2019.

\paragraph{Findings} Specifically, we use machine learning to measure the tone of over 200,000 twitter interactions and find that male politicians are more likely to attack political opponents representing 'outgroups' and that male representatives receive more outgroup negativity. Short title: A Male Hostility Spiral? show that the relation between politician i and j converge to 0 if their different sentiments take out each other.

\paragraph{Contribution} Provides an external benchmark for mechanisms that may or may not travel to China and WeChat.

\paragraph{Limitations} Design details need closer verification before the paper is used for strong causal claims.

\paragraph{Relevance to WeChat / My Project} Indirect benchmark: useful for comparison on platform effects or official communication, but needs careful translation to China's WeChat-centered environment.

\paragraph{Related Debate} Which actors can credibly correct misinformation and under what informational conditions


\subsection{Enikolopov et al. 2026: Polarise and Rule Independent Media in...}
\textbf{Citation key}: \texttt{enikolopov\_2026\_ej}.\\
\textbf{Full citation}: Ruben Enikolopov, Michael Rochlitz, Koen Schoors, and Nikita Zakharov 2026. "Polarise and Rule: Independent Media in Autocracy and the Role of Social Media." The Economic Journal.

\paragraph{Summary} This study examines polarise and rule: independent media in autocracy and the role of social media using descriptive or unclear from available parse evidence and reports that in both experiments, we find that independent media foster polarisation. Our results highlight how social media can condition the effect of...

\paragraph{Context} However, this effect only applies to voters who rely on news from social media.

\paragraph{Main Argument} The paper's core mechanism is that in both experiments, we find that independent media foster polarisation. Our results highlight how social media can condition the effect of traditional media.

\paragraph{Independent Variable(s)} do independent media

\paragraph{Dependent Variable(s)} the support of the regime in an autocracy? We carried out two complementary field experiments in Russia at the city...

\paragraph{Methods and Design} Descriptive or unclear from available parse. Observational or descriptive design; verify details from full text if heavily used. Data / sample: ED of watching the channel using O 15 Due to a lack of data on online posts in 2011, we are not able to perform a placebo test as we did for the case of TV Rain viewership. 15 Downloaded from...

\paragraph{Findings} In both experiments, we find that independent media foster polarisation. Our results highlight how social media can condition the effect of traditional media.

\paragraph{Contribution} Provides an external benchmark for mechanisms that may or may not travel to China and WeChat.

\paragraph{Limitations} Design details need closer verification before the paper is used for strong causal claims.

\paragraph{Relevance to WeChat / My Project} Indirect benchmark: useful for comparison on platform effects or official communication, but needs careful translation to China's WeChat-centered environment.

\paragraph{Related Debate} Whether digital control works through persuasion, selective amplification, covert influence, or audience management


\subsection{Germano 2026: Ranking for engagement How social media...}
\textbf{Citation key}: \texttt{germano\_2026\_jpubeco}.\\
\textbf{Full citation}: Fabrizio Germano 2026. "Ranking for engagement: How social media algorithms fuel misinformation and polarization." Journal of Public Economics.

\paragraph{Summary} This study examines ranking for engagement: how social media algorithms fuel misinformation and polarization using descriptive or unclear from available parse evidence and reports that put simply, these algorithms decide what information to show to users and, importantly, also in what order to show it. are...

\paragraph{Context} Digital platforms play an increasingly central role in the consumption of political news (Aridor et al., 2024). A defining feature of these plat forms is the use of recommendation algorithms that rank content based on measures of user engagement---most...

\paragraph{Main Argument} The paper's core mechanism is that put simply, these algorithms decide what information to show to users and, importantly, also in what order to show it. are representative of the US voting-age population.

\paragraph{Independent Variable(s)} Not stated in standard IV/DV terms; inferred from title/abstract if needed

\paragraph{Dependent Variable(s)} Not stated in standard IV/DV terms; verify if central to argument

\paragraph{Methods and Design} Descriptive or unclear from available parse. Observational or descriptive design; verify details from full text if heavily used. Data / sample: Put simply, these algorithms decide what information to show to users and, importantly, also in what order to show it. Accordingly, we estimate the following econometric specification: 27 An illustrative case is the...

\paragraph{Findings} Put simply, these algorithms decide what information to show to users and, importantly, also in what order to show it. are representative of the US voting-age population.

\paragraph{Contribution} Provides an external benchmark for mechanisms that may or may not travel to China and WeChat.

\paragraph{Limitations} Design details need closer verification before the paper is used for strong causal claims.

\paragraph{Relevance to WeChat / My Project} Indirect benchmark: useful for comparison on platform effects or official communication, but needs careful translation to China's WeChat-centered environment.

\paragraph{Related Debate} How interface design changes the visibility and meaning of digital expression


\subsection{Hirshleifer et al. 2026: The spread of mis information A...}
\textbf{Citation key}: \texttt{hirshleifer\_2026\_jde}.\\
\textbf{Full citation}: Sarojini Hirshleifer, Mustafa Naseem, Agha Ali Raza, and Arman Rezaee 2026. "The spread of (mis)information: A social media experiment in Pakistan." Journal of Development Economics.

\paragraph{Summary} This paper studies the spread of (mis)information: a social media experiment in pakistan. rna lP This study examines how controlling misinformation on a social media platform in Pakistan affects users' exposure to both accurate and false information.

\paragraph{Context} rna lP This study examines how controlling misinformation on a social media platform in Pakistan affects users' exposure to both accurate and false information. It combines an intervention to disseminate official information about the COVID-19 pandemic...

\paragraph{Main Argument} Core mechanism maps onto the cluster on platform affordances and public-private boundaries, but the exact statement should be verified from the full text if heavily used.

\paragraph{Independent Variable(s)} fully controlling access to pandemic-related misinformation. The treatments rely

\paragraph{Dependent Variable(s)} a higher-intensity

\paragraph{Methods and Design} randomized experiment. randomized experiment. Data / sample: rna lP This study examines how controlling misinformation on a social media platform in Pakistan affects users' exposure to both accurate and false information. Fully controlling misinformation, as in the treatments,...

\paragraph{Findings} No clean finding sentence detected; verify from full text if central.

\paragraph{Contribution} Provides an external benchmark for mechanisms that may or may not travel to China and WeChat.

\paragraph{Limitations} Design details need closer verification before the paper is used for strong causal claims.

\paragraph{Relevance to WeChat / My Project} Indirect benchmark: useful for comparison on platform effects or official communication, but needs careful translation to China's WeChat-centered environment.

\paragraph{Related Debate} How interface design changes the visibility and meaning of digital expression


\subsection{Lai 2026: Can autocratic power influence the media...}
\textbf{Citation key}: \texttt{lai\_2026\_psrm}.\\
\textbf{Full citation}: Ruilin Lai 2026. "Can autocratic power influence the media in democracies? Evidence from China's expulsion of American journalists." Political Science Research and Methods.

\paragraph{Summary} This study examines can autocratic power influence the media in democracies? evidence from china's expulsion of american journalists using descriptive or unclear from available parse evidence and reports that exploiting China's sudden expulsion of American journalists in March 2020, I show that US news outlets...

\paragraph{Context} Can autocratic governments influence foreign media? This study finds that political pressures from autocrats can motivate the media in democracies to tone down their negativity.

\paragraph{Main Argument} The paper's core mechanism is that exploiting China's sudden expulsion of American journalists in March 2020, I show that US news outlets targeted by the expulsion adopted a more positive tone toward China in their subsequent coverage, compared to outlets...

\paragraph{Independent Variable(s)} Author inference from title/abstract: Can autocratic power influence the media in democracies? Evidence from China's expulsion of American journalists

\paragraph{Dependent Variable(s)} Not stated in standard IV/DV terms

\paragraph{Methods and Design} Descriptive or unclear from available parse. Observational or descriptive design; verify details from full text if heavily used. Data / sample: Unclear from available parse; verify from full text if used centrally

\paragraph{Findings} Exploiting China's sudden expulsion of American journalists in March 2020, I show that US news outlets targeted by the expulsion adopted a more positive tone toward China in their subsequent coverage, compared to outlets that were not targeted. The findings highlight the overt threats autocracies pose to media freedom, a fundamental pillar of democratic societies. https://doi.org/10.1017/psrm.2025.5 Published online by Cambridge University Press generalizable to other contexts, given that the estimates are specific to China's expulsion of American journalists?

\paragraph{Contribution} Provides an external benchmark for mechanisms that may or may not travel to China and WeChat.

\paragraph{Limitations} Design details need closer verification before the paper is used for strong causal claims.

\paragraph{Relevance to WeChat / My Project} Peripheral for a WeChat project; best used as a comparative benchmark rather than core theoretical evidence.

\paragraph{Related Debate} Whether digital control works through persuasion, selective amplification, covert influence, or audience management


\subsection{Rohner et al. 2026: Science under Threat A Natural Experiment...}
\textbf{Citation key}: \texttt{rohner\_2026\_wp}.\\
\textbf{Full citation}: Dominic Rohner, Oliver Vanden Eynde, and Philine Widmer 2026. Science under Threat? A Natural Experiment in Economics. Elsevier BV.

\paragraph{Summary} This study examines science under threat? a natural experiment in economics using descriptive or unclear from available parse evidence and reports that we find that the release of this list leads to a sharp reduction in the use of banned words in sensitive contexts among economists working at universities that rely...

\paragraph{Context} Academic freedom has come under growing strain across the world. To study whether and how academics react to political pressure, we exploit a natural experiment: the U.S. government's "blacklist" of undesirable words released in early 2025.

\paragraph{Main Argument} The paper's core mechanism is that we find that the release of this list leads to a sharp reduction in the use of banned words in sensitive contexts among economists working at universities that rely heavily on NSF funding.

\paragraph{Independent Variable(s)} Author inference from title/abstract: Science under Threat? A Natural Experiment in Economics

\paragraph{Dependent Variable(s)} Not stated in standard IV/DV terms

\paragraph{Methods and Design} Descriptive or unclear from available parse. Observational or descriptive design; verify details from full text if heavily used. Data / sample: Unclear from available parse; verify from full text if used centrally

\paragraph{Findings} We find that the release of this list leads to a sharp reduction in the use of banned words in sensitive contexts among economists working at universities that rely heavily on NSF funding.

\paragraph{Contribution} Provides an external benchmark for mechanisms that may or may not travel to China and WeChat.

\paragraph{Limitations} Mechanisms may not transfer cleanly to WeChat's hybrid private-public architecture without further adaptation.

\paragraph{Relevance to WeChat / My Project} Peripheral for a WeChat project; best used as a comparative benchmark rather than core theoretical evidence.

\paragraph{Related Debate} How findings from not platform-specific and united states settings travel to a WeChat-centered China project


\subsection{Roshchin 2026: Professorial Silence Academic Freedom Domination and...}
\textbf{Citation key}: \texttt{roshchin\_2026\_jop}.\\
\textbf{Full citation}: Evgeny Roshchin 2026. "Professorial Silence: Academic Freedom, Domination, and Self-Censorship in Contemporary Russia." The Journal of Politics.

\paragraph{Summary} This paper studies professorial silence: academic freedom, domination, and self-censorship in contemporary russia. of repressive wartime laws in 2022.

\paragraph{Context} Given the ubiquity of the practice in the Russian context, this article asks the question of what makes faculty in the country turn to self-censorship, restraining themselves as professors in truth-seeking and citizens in their public roles.

\paragraph{Main Argument} Core mechanism maps onto the cluster on censorship, propaganda, and state influence, but the exact statement should be verified from the full text if heavily used.

\paragraph{Independent Variable(s)} Not stated in standard IV/DV terms; inferred from title/abstract if needed

\paragraph{Dependent Variable(s)} Not stated in standard IV/DV terms; verify if central to argument

\paragraph{Methods and Design} Descriptive or unclear from available parse. Observational or descriptive design; verify details from full text if heavily used. Data / sample: Unclear from available parse; verify from full text if used centrally

\paragraph{Findings} No clean finding sentence detected; verify from full text if central.

\paragraph{Contribution} Provides an external benchmark for mechanisms that may or may not travel to China and WeChat.

\paragraph{Limitations} Mechanisms may not transfer cleanly to WeChat's hybrid private-public architecture without further adaptation.

\paragraph{Relevance to WeChat / My Project} Peripheral for a WeChat project; best used as a comparative benchmark rather than core theoretical evidence.

\paragraph{Related Debate} Whether digital control works through persuasion, selective amplification, covert influence, or audience management


\subsection{Schmidt-Feuerheerd 2026: Collaborative authoritarianism and its unintended consequences}
\textbf{Citation key}: \texttt{schmidt\_feuerheerd\_2026\_democratization}.\\
\textbf{Full citation}: Bruno Schmidt-Feuerheerd 2026. "Collaborative authoritarianism and its unintended consequences." Democratization.

\paragraph{Summary} This paper studies collaborative authoritarianism and its unintended consequences. Why do authoritarian regimes maintain access to open U.S.-owned social media platforms like Twitter and how do pro-government supporters use them?

\paragraph{Context} While the literature on consultative authoritarianism explains how limited public engagement through purpose-built input institutions can be used to gauge public grievances and boost regime legitimacy, we know less about the implications of open social...

\paragraph{Main Argument} Core mechanism maps onto the cluster on interpersonal discussion, networks, and opinion formation, but the exact statement should be verified from the full text if heavily used.

\paragraph{Independent Variable(s)} Not stated in standard IV/DV terms; inferred from title/abstract if needed

\paragraph{Dependent Variable(s)} Not stated in standard IV/DV terms; verify if central to argument

\paragraph{Methods and Design} interviews / focus groups. Observational or descriptive design; verify details from full text if heavily used. Data / sample: Introducing the case of Twitter use in Saudi Arabia, this article draws on around 100 interviews that I conducted with patriotic and nationalist journalists, intellectuals, and Twitter influencers across Saudi Arabia...

\paragraph{Findings} No clean finding sentence detected; verify from full text if central.

\paragraph{Contribution} Provides an external benchmark for mechanisms that may or may not travel to China and WeChat.

\paragraph{Limitations} Design details need closer verification before the paper is used for strong causal claims.

\paragraph{Relevance to WeChat / My Project} Peripheral for a WeChat project; best used as a comparative benchmark rather than core theoretical evidence.

\paragraph{Related Debate} How social ties and discussion networks mediate willingness to speak, share, and update opinions


\subsection{Yu and Shen 2026: Serve the People Public Service Information...}
\textbf{Citation key}: \texttt{yu\_2026\_jcc}.\\
\textbf{Full citation}: Jianjun Yu and Shuyuan Shen 2026. "Serve the People: Public Service Information and Communication Strategies of Local Governments in China." Journal of Contemporary China.

\paragraph{Summary} This paper studies serve the people: public service information and communication strategies of local governments in china. Propaganda remains central to the Chinese government's efforts to maintain legitimacy and public support.1 Yet in the digital age, where citizens navigate an information-rich environment,...

\paragraph{Context} In Russia, for example, coordinated networks of automated and human accounts amplify pro-government messages,3 while Iran employs state and social media to spread CONTACT Shuyuan Shen Shuyuan.Shen@lmu.edu Department of Political Science and International...

\paragraph{Main Argument} Core mechanism maps onto the cluster on censorship, propaganda, and state influence, but the exact statement should be verified from the full text if heavily used.

\paragraph{Independent Variable(s)} Not stated in standard IV/DV terms; inferred from title/abstract if needed

\paragraph{Dependent Variable(s)} Not stated in standard IV/DV terms; verify if central to argument

\paragraph{Methods and Design} survey experiment. survey experiment. Data / sample: In Russia, for example, coordinated networks of automated and human accounts amplify pro-government messages,3 while Iran employs state and social media to spread CONTACT Shuyuan Shen Shuyuan.Shen@lmu.edu Department...

\paragraph{Findings} No clean finding sentence detected; verify from full text if central.

\paragraph{Contribution} Offers comparatively strong leverage on causal mechanisms in digital political communication.

\paragraph{Limitations} Internal validity is strong, but external validity may depend on whether experimental treatments resemble real platform use.

\paragraph{Relevance to WeChat / My Project} Conceptually useful for a WeChat project because it speaks to the Chinese information environment, especially the cluster on censorship, propaganda, and state influence, even if the platform is not WeChat.

\paragraph{Related Debate} Whether digital control works through persuasion, selective amplification, covert influence, or audience management


\subsection{Zakharov 2026: Propaganda to a cynical audience}
\textbf{Citation key}: \texttt{zakharov\_2026\_psrm}.\\
\textbf{Full citation}: Alexei Zakharov 2026. "Propaganda to a cynical audience." Political Science Research and Methods.

\paragraph{Summary} This paper studies propaganda to a cynical audience. Using a model, we explain why propaganda in autocracies can be blatantly false and unconvincing.

\paragraph{Context} Using a model, we explain why propaganda in autocracies can be blatantly false and unconvincing. We model two news outlets that report on a hidden state of the world, motivated by the ex-post beliefs of the audience about the state of the world.

\paragraph{Main Argument} Core mechanism maps onto the cluster on censorship, propaganda, and state influence, but the exact statement should be verified from the full text if heavily used.

\paragraph{Independent Variable(s)} Not stated in standard IV/DV terms; inferred from title/abstract if needed

\paragraph{Dependent Variable(s)} Not stated in standard IV/DV terms; verify if central to argument

\paragraph{Methods and Design} Descriptive or unclear from available parse. Observational or descriptive design; verify details from full text if heavily used. Data / sample: Unclear from available parse; verify from full text if used centrally

\paragraph{Findings} No clean finding sentence detected; verify from full text if central.

\paragraph{Contribution} Provides an external benchmark for mechanisms that may or may not travel to China and WeChat.

\paragraph{Limitations} Mechanisms may not transfer cleanly to WeChat's hybrid private-public architecture without further adaptation.

\paragraph{Relevance to WeChat / My Project} Peripheral for a WeChat project; best used as a comparative benchmark rather than core theoretical evidence.

\paragraph{Related Debate} Whether digital control works through persuasion, selective amplification, covert influence, or audience management



        \section{Important Missing References}

        The reference-list parse across the corpus repeatedly surfaces a smaller set of foundational works that are not themselves in the Information folder but matter for positioning the project. The list below is conservative: it only includes references recoverable from the local writing bibliography and mentioned repeatedly in the corpus.

        \begin{itemize}[leftmargin=1.5em]
        \item \textbf{Zheng2013} (105 corpus mentions). Zheng, Lei (2013), \emph{Social Media in Chinese Government: Drivers, Challenges and Capabilities}. Why it matters: cited across the corpus as a media effects and selective exposure anchor; useful for positioning the WeChat project against the broader China media-effects literature.
\item \textbf{Kuran1995} (99 corpus mentions). Kuran, Timur (1995), \emph{Private Truths, Public Lies: The Social Consequences of Preference Falsification}. Why it matters: cited across the corpus as a opinion expression and self-censorship anchor; useful for positioning the WeChat project against the broader China media-effects literature.
\item \textbf{KingPanRoberts2017} (97 corpus mentions). King, Gary and Pan, Jennifer and Roberts, Margaret E. (2017), \emph{How the Chinese Government Fabricates Social Media Posts for Strategic Distraction, Not Engaged Argument}. Why it matters: cited across the corpus as a media effects and selective exposure anchor; useful for positioning the WeChat project against the broader China media-effects literature.
\item \textbf{ZhengZheng2014} (97 corpus mentions). Zheng, Lei and Zheng, Tuo (2014), \emph{Innovation Through Social Media in the Public Sector: Information and Interactions}. Why it matters: cited across the corpus as a media effects and selective exposure anchor; useful for positioning the WeChat project against the broader China media-effects literature.
\item \textbf{Ng2015} (94 corpus mentions). Ng, Jason Q. (2015), \emph{Politics, Rumors, and Ambiguity: Tracking Censorship on \{WeChat\}'s Public Accounts Platform}. Why it matters: cited across the corpus as a directly about wechat and networked communication anchor; useful for positioning the WeChat project against the broader China media-effects literature.
\item \textbf{KingPanRoberts2013} (91 corpus mentions). King, Gary and Pan, Jennifer and Roberts, Margaret E. (2013), \emph{How Censorship in China Allows Government Criticism but Silences Collective Expression}. Why it matters: cited across the corpus as a censorship, propaganda, and state influence anchor; useful for positioning the WeChat project against the broader China media-effects literature.
\item \textbf{LuPan2021} (84 corpus mentions). Lu, Yingdan and Pan, Jennifer (2021), \emph{Capturing Clicks: How the Chinese Government Uses Clickbait to Compete for Visibility}. Why it matters: cited across the corpus as a foundational reference for china propaganda, censorship, or wechat-adjacent political communication anchor; useful for positioning the WeChat project against the broader China media-effects literature.
\item \textbf{PanXu2018} (79 corpus mentions). Pan, Jennifer and Xu, Yiqing (2018), \emph{China's Ideological Spectrum}. Why it matters: cited across the corpus as a foundational reference for china propaganda, censorship, or wechat-adjacent political communication anchor; useful for positioning the WeChat project against the broader China media-effects literature.
\item \textbf{NoelleNeumann1993} (76 corpus mentions). Noelle-Neumann, Elisabeth (1993), \emph{The Spiral of Silence: Public Opinion---Our Social Skin}. Why it matters: cited across the corpus as a foundational reference for china propaganda, censorship, or wechat-adjacent political communication anchor; useful for positioning the WeChat project against the broader China media-effects literature.
\item \textbf{Pan2020} (76 corpus mentions). Pan, Jennifer (2020), \emph{Welfare for Autocrats: How Social Assistance in China Cares for Its Rulers}. Why it matters: cited across the corpus as a foundational reference for china propaganda, censorship, or wechat-adjacent political communication anchor; useful for positioning the WeChat project against the broader China media-effects literature.
\item \textbf{Mattingly2020} (75 corpus mentions). Mattingly, Daniel C. (2020), \emph{The Art of Political Control in China}. Why it matters: cited across the corpus as a foundational reference for china propaganda, censorship, or wechat-adjacent political communication anchor; useful for positioning the WeChat project against the broader China media-effects literature.
\item \textbf{WangWestwood2024} (66 corpus mentions). Wang, Clara and Westwood, Sean J. (2024), \emph{What Chinese Internet Users ``Like'' to Read: Selective Exposure in a Restricted Information Environment}. Why it matters: cited across the corpus as a media effects and selective exposure anchor; useful for positioning the WeChat project against the broader China media-effects literature.
\item \textbf{Stockmann2013} (61 corpus mentions). Stockmann, Daniela (2013), \emph{Media Commercialization and Authoritarian Rule in China}. Why it matters: cited across the corpus as a media effects and selective exposure anchor; useful for positioning the WeChat project against the broader China media-effects literature.
\item \textbf{TsaiMen2018} (60 corpus mentions). Tsai, Wan-Hsiu Sunny and Men, Linjuan Rita (2018), \emph{Social Messengers as the New Frontier of Organization-Public Engagement: A \{WeChat\} Study}. Why it matters: cited across the corpus as a directly about wechat and networked communication anchor; useful for positioning the WeChat project against the broader China media-effects literature.
\item \textbf{Xu2022WeChat} (48 corpus mentions). Xu, Vivien Weiwei (2022), \emph{\{WeChat\} and the Sharing of News in Networked China}. Why it matters: cited across the corpus as a directly about wechat and networked communication anchor; useful for positioning the WeChat project against the broader China media-effects literature.
        \end{itemize}

        \section{Implications for My Project}

        The literature already establishes several pieces that a WeChat project can build on. First, official digital communication in China is not reducible to crude indoctrination: it often mixes service content, audience management, and softer forms of persuasion. Second, WeChat-specific studies suggest that official accounts function as everyday governance tools as much as propaganda outlets. Third, comparative media-effects research provides stronger leverage on misinformation, selective exposure, and platform design, but often on platforms where the expressive meaning of a click is more public and less institutionally embedded than on WeChat.

        What remains under-studied is the political meaning of interaction itself on WeChat. The literature is much better at studying whether people were exposed to content than at distinguishing low-visibility consumption from socially legible endorsement. It is also better at studying generic social media than at theorizing WeChat's specific mix of semi-private discussion, public-account broadcasting, and state-linked infrastructural power. A WeChat-focused project can therefore contribute by theorizing layered expression, measuring how platform redesign changes the meaning of engagement, and separating service-driven audience accumulation from overt political persuasion.

        \bibliographystyle{plainnat}
        \bibliography{references}
        \end{document}
